% !TEX program = xelatex
% Test document for figset_D.tex (figures/test_fd2), generated by mkdoc.py.
% !TEX program = xelatex
% 编译：xelatex qcd_factorization.tex（运行三次以生成目录与交叉引用）
% 字体：fontset=fandol 随 TeX Live 发行，可在任何平台编译；
%       在 macOS 上可改为 fontset=macnew 以获得更好的中文字体效果。
\documentclass[UTF8,a4paper,11pt,oneside,openany,fontset=fandol]{ctexbook}

\usepackage[margin=2.5cm,headheight=14pt]{geometry}
\usepackage{amsmath,amssymb,amsthm,mathtools,bm}
% \mathbb 1：amssymb 的 msbm 字体没有数字 1，改用 dsfont 的 \mathds{1}；\mathbb{C}、\mathbb{R} 等不变
\usepackage{dsfont}
\let\amsmathbb\mathbb
\renewcommand{\mathbb}[1]{\ifx1#1\mathds{1}\else\amsmathbb{#1}\fi}
\usepackage{slashed}
\usepackage{booktabs,array,multirow,tabularx}
\usepackage{enumitem}
\usepackage{caption,subcaption}
\captionsetup{labelsep=quad}     % 图表标题 “图 5　……”：编号后空一字，不用西文冒号
\usepackage[section]{placeins}   % 浮动体不越过 \section 边界
\usepackage[dvipsnames]{xcolor}
\usepackage{tikz}
\usetikzlibrary{decorations.pathmorphing,decorations.markings,arrows.meta,calc,positioning,shapes.geometric,patterns,backgrounds,fit}
\usepackage{pgfplots}
\pgfplotsset{compat=1.17}
\usepackage[most]{tcolorbox}
\usepackage{etoolbox}            % \patchcmd、\appto（tcolorbox 亦载入，此处显式声明）
\usepackage[colorlinks=true,linkcolor=blue!50!black,citecolor=green!40!black,
            urlcolor=blue!60!black,bookmarksnumbered=true]{hyperref}
\usepackage[nameinlink]{cleveref}
\usepackage[nottoc]{tocbibind}   % 参考文献进入目录与 PDF 书签

% ---------------- 定理类环境 ----------------
% 定理、引理、命题、推论的正文：汉字用楷体，西文保持直立。plain 样式的 \itshape
% 会把对称性名称 “C 宇称/C 偶/C 奇” 中的正体 C 排成斜体，与电荷共轭矩阵 $C$
% 无法区分；上下间距与标题字体同 plain。
% 三种样式的标题都按中文习惯排：附注用全角括号（……），标题后不加西文句点而空一字；
% cdefinition、cremark 的字体与上下间距分别同 amsthm 的 definition、remark 样式。
\newtheoremstyle{cplain}{\topsep}{\topsep}{\kaishu}{0pt}{\bfseries}{}{1em}%
  {\thmname{#1}\thmnumber{ \textup{#2}}\thmnote{\normalfont（#3）}}
\newtheoremstyle{cdefinition}{\topsep}{\topsep}{\normalfont}{0pt}{\bfseries}{}{1em}%
  {\thmname{#1}\thmnumber{ \textup{#2}}\thmnote{\normalfont（#3）}}
\newtheoremstyle{cremark}{\glueexpr\topsep/2\relax}{\glueexpr\topsep/2\relax}%
  {\normalfont}{0pt}{\itshape}{}{1em}%
  {\thmname{#1}\thmnumber{ \textup{#2}}\thmnote{\normalfont（#3）}}
\theoremstyle{cplain}
\newtheorem{theorem}{定理}[chapter]
\newtheorem{lemma}[theorem]{引理}
\newtheorem{proposition}[theorem]{命题}
\newtheorem{corollary}[theorem]{推论}
\newtheorem{conjecture}[theorem]{猜想}
\theoremstyle{cdefinition}
\newtheorem{definition}[theorem]{定义}
\newtheorem{assumption}[theorem]{假设}
\newtheorem{example}[theorem]{例}
\theoremstyle{cremark}
\newtheorem{remark}[theorem]{注}
\renewcommand{\proofname}{证明}
% 证明的标题同样不加句点：amsthm 原定义中的 \@addpunct{.} 改为 0.5em，连同 \labelsep 共空一字
\makeatletter
\renewenvironment{proof}[1][\proofname]{\par\pushQED{\qed}\normalfont\topsep6\p@\@plus6\p@\relax
  \trivlist\item[\hskip\labelsep\itshape #1\hspace{0.5em}]\ignorespaces}{\popQED\endtrivlist\@endpefalse}
\makeatother

\tcolorboxenvironment{theorem}{enhanced,breakable,colback=blue!3!white,
  colframe=blue!45!black,boxrule=0.6pt,arc=2pt,left=5pt,right=5pt,top=3pt,bottom=3pt}

\tcolorboxenvironment{assumption}{enhanced,breakable,colback=orange!4!white,
  colframe=orange!65!black,boxrule=0.6pt,arc=2pt,left=5pt,right=5pt,top=3pt,bottom=3pt}
\tcolorboxenvironment{conjecture}{enhanced,breakable,colback=Plum!4!white,
  colframe=Plum!70!black,boxrule=0.6pt,arc=2pt,left=5pt,right=5pt,top=3pt,bottom=3pt}

\crefname{theorem}{定理}{定理}
\crefname{lemma}{引理}{引理}
\crefname{proposition}{命题}{命题}
\crefname{corollary}{推论}{推论}
\crefname{conjecture}{猜想}{猜想}
\crefname{assumption}{假设}{假设}
\crefname{definition}{定义}{定义}
\crefname{example}{例}{例}
\crefname{remark}{注}{注}
\crefname{equation}{式}{式}
\creflabelformat{equation}{#2(#1)#3}
\crefname{figure}{图}{图}
\crefname{table}{表}{表}
\crefformat{chapter}{#2第~#1~章#3}
\crefformat{section}{#2第~#1~节#3}
\crefformat{subsection}{#2第~#1~小节#3}
\crefformat{appendix}{#2附录~#1#3}
\crefformat{subappendix}{#2附录~#1#3}
\Crefformat{chapter}{#2第~#1~章#3}
\Crefformat{section}{#2第~#1~节#3}
\Crefformat{subsection}{#2第~#1~小节#3}
\Crefformat{appendix}{#2附录~#1#3}
\Crefformat{subappendix}{#2附录~#1#3}
% 多目标引用的中文连接词
\newcommand{\crefpairconjunction}{\CJKecglue 与\CJKecglue}
\newcommand{\crefmiddleconjunction}{、}
\newcommand{\creflastconjunction}{\CJKecglue 与\CJKecglue}
\newcommand{\crefpairgroupconjunction}{\CJKecglue 与\CJKecglue}
\newcommand{\crefmiddlegroupconjunction}{、}
\newcommand{\creflastgroupconjunction}{\CJKecglue 与\CJKecglue}
\newcommand{\crefrangeconjunction}{\CJKecglue 至\CJKecglue}

% \cref 后紧跟汉字时补中西文间距：xeCJK 不能越过超链接的边界自动插入 \CJKecglue
% （源文件中 \cref 与汉字之间不留空格），此处手工补上。
% 最后一组引用若为 section/subsection 类型（“第 n 节/小节”，以汉字结尾）则不补。
\ExplSyntaxOn
\NewCommandCopy \__fact_cref_orig:w \cref
\tl_new:N \l__fact_cref_type_tl
\seq_new:N \l__fact_cref_types_seq
% 标签 #1 的 cleveref 类型，取自 \r@#1@cref = {[类型][..][..]..}{..}；标签未定义时为空
\cs_new:Npn \__fact_cref_type:n #1
  {
    \cs_if_exist:cT { r@ #1 @cref }
      { \exp_last_unbraced:Nv \__fact_cref_type_aux:w { r@ #1 @cref } \q_stop }
  }
\cs_new:Npn \__fact_cref_type_aux:w #1#2 \q_stop
  { \__fact_cref_type_auxii:w #1 [] \q_stop }
\cs_new:Npn \__fact_cref_type_auxii:w #1 [#2] #3 \q_stop
  { \tl_if_empty:nT {#1} {#2} }
% 输出中最后一组引用的类型（cleveref 默认按各类型首次出现的次序分组输出）
\cs_new_protected:Npn \__fact_cref_last_type:n #1
  {
    \seq_clear:N \l__fact_cref_types_seq
    \clist_map_inline:nn {#1}
      {
        \tl_set:Nx \l__fact_cref_type_tl { \__fact_cref_type:n {##1} }
        \seq_if_in:NVF \l__fact_cref_types_seq \l__fact_cref_type_tl
          { \seq_put_right:NV \l__fact_cref_types_seq \l__fact_cref_type_tl }
      }
    \seq_get_right:NNF \l__fact_cref_types_seq \l__fact_cref_type_tl
      { \tl_clear:N \l__fact_cref_type_tl }
  }
% 下一记号若为 CJK 统一汉字（U+3400--U+9FFF），在其前插入 \CJKecglue
\cs_new_protected:Npn \__fact_cref_ecglue:N #1
  {
    \bool_lazy_all:nTF
      {
        { \token_if_letter_p:N #1 }
        { ! \token_if_cs_p:N #1 }
        { \int_compare_p:nNn { `#1 } > { "33FF } }
        { \int_compare_p:nNn { `#1 } < { "A000 } }
      }
      { \CJKecglue #1 }
      { #1 }
  }
\RenewDocumentCommand \cref { s m }
  {
    \IfBooleanTF {#1}
      { \__fact_cref_orig:w * {#2} }
      { \__fact_cref_orig:w {#2} }
    \__fact_cref_last_type:n {#2}
    \str_case:VnF \l__fact_cref_type_tl
      { { chapter } { } { section } { } { subsection } { } }
      { \peek_N_type:TF { \__fact_cref_ecglue:N } { } }
  }
\ExplSyntaxOff

\numberwithin{equation}{chapter}

% 标题中的数学式随标题加粗（含目录中节一级条目）
\ctexset{chapter/number=\arabic{chapter}}   % 章号用阿拉伯数字：“第 5 章”、式 (5.3)
\ctexset{chapter/format+=\boldmath,section/format+=\boldmath,subsection/format+=\boldmath,paragraph/format+=\boldmath}
\makeatletter
\patchcmd{\l@section}{\bfseries}{\bfseries\boldmath}{}%
  {\@latex@warning@no@line{\string\l@section\space not patched: math in TOC section entries stays non-bold}}
\makeatother

% 目录行距略紧并单独成页（正文从新页开始）
\let\factorigtableofcontents\tableofcontents
\renewcommand{\tableofcontents}{%
  \begingroup\linespread{1.1}\selectfont\factorigtableofcontents\endgroup\clearpage}

% 参考文献条目间距略紧。\parsep 须在列表参数中设定（\list 随即把它赋给 \parskip），
% 故写入 thebibliography 在列表参数中执行的 \@openbib@code，而不追加在 \thebibliography 之后
\makeatletter
\appto\@openbib@code{\setlength{\itemsep}{1pt plus 1pt}\setlength{\parsep}{0pt}}
\makeatother

% 页眉：节标题不转为大写（避免 “DIRAC”“FURRY” 之类）；页眉一律用直立体（含目录、
% 参考文献页）：headings 样式原以 \slshape 排页眉，而 FandolSong 没有倾斜体，汉字会
% 退成楷体，西文则倾斜
\makeatletter
\renewcommand{\chaptermark}[1]{\markright{\ifnum\c@secnumdepth>\m@ne\CTEXthechapter\quad\fi #1}}
\renewcommand{\sectionmark}[1]{}
\def\@oddhead{{\upshape\rightmark}\hfil\thepage}
\makeatother

% ---------------- 记号宏 ----------------
\DeclareMathOperator{\Tr}{Tr}
\DeclareMathOperator{\Real}{Re}
\DeclareMathOperator{\Imag}{Im}
\DeclareMathOperator{\Disc}{Disc}
\newcommand{\TrD}{\operatorname{Tr}_{\mathrm{D}}}   % Dirac 迹
\newcommand{\Trc}{\operatorname{Tr}_{\mathrm{c}}}   % 色迹
\newcommand{\sla}[1]{\slashed{#1}}                  % Feynman slash
\newcommand{\dd}{\mathrm{d}}                         % 微分符号
\newcommand{\tp}{\mathsf{T}}                         % 转置
\newcommand{\ii}{\mathrm{i}}                        % 虚数单位（正体）
\newcommand{\ee}{\mathrm{e}}                        % 自然对数底（正体）
\newcommand{\MSbar}{\overline{\mathrm{MS}}}         % MS-bar 方案
\newcommand{\order}[1]{\mathcal{O}\!\left(#1\right)}% 量级
\newcommand{\Lqcd}{\Lambda_{\mathrm{QCD}}}
\newcommand{\as}{\alpha_s}
\newcommand{\pT}{p_{\mathrm{T}}}
\newcommand{\kT}{k_{\mathrm{T}}}
\newcommand{\bT}{b_{\mathrm{T}}}
\newcommand{\qT}{q_{\mathrm{T}}}

% ---------------- 严格性标签 ----------------
% [R] 严格  [A] 全阶论证  [P] 部分  [C] 猜想  [V] 破坏
\newtcbox{\facttagbox}[1][gray]{on line,arc=2pt,boxsep=0pt,left=2.2pt,right=2.2pt,top=1.2pt,bottom=1.2pt,
  boxrule=0.45pt,colback=#1!10!white,colframe=#1!70!black,fontupper=\footnotesize\sffamily\upshape}
\newcommand{\tagR}{\facttagbox[OliveGreen]{R\,严格}}
\newcommand{\tagA}{\facttagbox[RoyalBlue]{A\,全阶论证}}
\newcommand{\tagP}{\facttagbox[BurntOrange]{P\,部分}}
\newcommand{\tagC}{\facttagbox[Plum]{C\,猜想}}
\newcommand{\tagV}{\facttagbox[BrickRed]{V\,破坏}}
% 依赖记录（仿 Lean blueprint 的 \uses）：在定义/假设/引理/定理等环境内部写 \uses{lab1,lab2}，
% 正文中以小号字列出所依赖的条目；附录中的“形式化蓝图”由脚本据此自动生成依赖表。
\newcommand{\uses}[1]{\par\noindent{\footnotesize\color{gray!45!black}依赖：\cref{#1}}\par}
% 只给出证明要点（对 [A]/[P] 结论）时使用
\newenvironment{proofsketch}[1][证明概要]{\begin{proof}[#1]}{\end{proof}}

% ---------------- TikZ 费曼图样式（纯 TikZ，不依赖 tikz-feynman） ----------------
%%% BEGIN STYLES
% =====================================================================
%  Plain-TikZ Feynman-diagram styles (no tikz-feynman).
%  Load in the preamble, after \usepackage{tikz}:
%    \usetikzlibrary{decorations.pathmorphing,decorations.markings,arrows.meta,calc,positioning}
%  Contents of this block (all low-level names start with \fd):
%   (i)  helper registers/macros and three decorations
%          "fd gluon"      coil for open paths,
%          "fd gluon loop" coil for closed paths (full circles),
%          "fd photon"     sine wave;
%        they rescale the nominal segment length so that an integer
%        number of coils / half-waves fits the path exactly (no straight
%        left-overs, evenly spaced coils, also on arcs);
%   (ii) a single \tikzset with the user-level styles:
%          fd             picture style (line width, caps, label font)
%          quark[=pos]    solid line, centred Stealth arrow at pos (default 0.5)
%          gluon          coiled line;   gluon loop[=m]: closed coil, number
%                         of coils a multiple of m (use for full circles)
%          photon         wavy line
%          ghost[=pos]    dotted line with centred arrow
%          mom={..}, mom'={..}   momentum arrow parallel to a path, on its
%                         left / right:  \path[mom={$p$}] (A) -- (B);
%                         keys: mom pos, mom length, mom sep
%          vertex         vertex dot (node style):  \node[vertex] at (A) {};
%          blob           grey filled region:  \draw[blob] (c) circle[radius=5mm];
%          bound state    double line (e.g. a bound state)
% =====================================================================
\newdimen\fdAA \newdimen\fdTW \newdimen\fdHW
\def\fdloopmult{1}
\def\fdpt#1#2#3{\pgfqpoint{\dimexpr#1\fdAA+#3\fdTW\relax}{#2\pgfdecorationsegmentamplitude}}
\def\fdcoilfull{%
  \pgfpathcurveto{\fdpt{0}{0.555}{1}}{\fdpt{0.445}{1}{2}}{\fdpt{1}{1}{3}}%
  \pgfpathcurveto{\fdpt{1.555}{1}{4}}{\fdpt{2}{0.555}{5}}{\fdpt{2}{0}{6}}%
  \pgfpathcurveto{\fdpt{2}{-0.555}{7}}{\fdpt{1.555}{-1}{8}}{\fdpt{1}{-1}{9}}%
  \pgfpathcurveto{\fdpt{0.445}{-1}{10}}{\fdpt{0}{-0.555}{11}}{\fdpt{0}{0}{12}}}
\def\fdcoilhalf{%
  \pgfpathcurveto{\fdpt{0}{0.555}{1}}{\fdpt{0.445}{1}{2}}{\fdpt{1}{1}{3}}%
  \pgfpathcurveto{\fdpt{1.555}{1}{4}}{\fdpt{2}{0.555}{5}}{\fdpt{2}{0}{6}}}
\def\fdcoilopeninit{%
  \fdAA=\pgfdecorationsegmentaspect\pgfdecorationsegmentamplitude\relax
  \pgfmathsetmacro\fdN{max(1,round((\pgfdecoratedpathlength-2*\fdAA)/\pgfdecorationsegmentlength-0.5))}%
  \pgfmathsetlength\pgfdecorationsegmentlength{(\pgfdecoratedpathlength-2*\fdAA-0.3pt)/(\fdN+0.5)}%
  \fdTW=0.0833333\pgfdecorationsegmentlength\relax}%
\def\fdcoilclosedinit{%
  \fdAA=\pgfdecorationsegmentaspect\pgfdecorationsegmentamplitude\relax
  \pgfmathsetmacro\fdN{\fdloopmult*max(1,round(\pgfdecoratedpathlength/(\fdloopmult*\pgfdecorationsegmentlength)))}%
  \pgfmathsetlength\pgfdecorationsegmentlength{(\pgfdecoratedpathlength-0.1pt)/\fdN}%
  \fdTW=0.0833333\pgfdecorationsegmentlength\relax}%
\def\fdphotoninit{%
  \pgfmathsetmacro\fdN{max(1,round(2*\pgfdecoratedpathlength/\pgfdecorationsegmentlength))}%
  \pgfmathsetlength\fdHW{(\pgfdecoratedpathlength-0.1pt)/\fdN}}%
\pgfdeclaredecoration{fd gluon}{init}{%
  \state{init}[width=+0pt,next state=coil,persistent precomputation=\fdcoilopeninit]{}%
  \state{coil}[switch if less than=1.5\pgfdecorationsegmentlength+2\fdAA to last,
               width=+\pgfdecorationsegmentlength]{\fdcoilfull}%
  \state{last}[width=+\dimexpr.5\pgfdecorationsegmentlength+2\fdAA\relax,next state=final]{\fdcoilhalf}%
  \state{final}{\pgfpathlineto{\pgfpointdecoratedpathlast}}}
\pgfdeclaredecoration{fd gluon loop}{init}{%
  \state{init}[width=+0pt,next state=coil,persistent precomputation=\fdcoilclosedinit]{}%
  \state{coil}[switch if less than=+\pgfdecorationsegmentlength to final,
               width=+\pgfdecorationsegmentlength]{\fdcoilfull}%
  \state{final}{\pgfpathlineto{\pgfpointdecoratedpathlast}}}
\pgfdeclaredecoration{fd photon}{init}{%
  \state{init}[width=+0pt,next state=up,persistent precomputation=\fdphotoninit]{}%
  \state{up}[switch if less than=+\fdHW to final,width=+\fdHW,next state=down]{%
    \pgfpathsine{\pgfqpoint{.5\fdHW}{\pgfdecorationsegmentamplitude}}%
    \pgfpathcosine{\pgfqpoint{.5\fdHW}{-\pgfdecorationsegmentamplitude}}}%
  \state{down}[switch if less than=+\fdHW to final,width=+\fdHW,next state=up]{%
    \pgfpathsine{\pgfqpoint{.5\fdHW}{-\pgfdecorationsegmentamplitude}}%
    \pgfpathcosine{\pgfqpoint{.5\fdHW}{\pgfdecorationsegmentamplitude}}}%
  \state{final}{\pgfpathlineto{\pgfpointdecoratedpathlast}}}
\tikzset{
  /pgf/decoration/coil multiple/.store in=\fdloopmult,
  fd/.style={line width=0.8pt,line cap=round,line join=round,font=\small},
  fd arrow/.style={postaction={decorate,solid,line width=0.8pt,decoration={markings,
      mark=at position #1 with {\arrow[xshift=2.9pt]{Stealth[length=5.5pt,width=4.5pt]}}}}},
  quark/.style={draw,solid,line width=0.8pt,fd arrow=#1},
  quark/.default=0.5,
  plain/.style={draw,solid,line width=0.8pt},
  gluon/.style={draw,solid,line width=0.75pt,decorate,
      decoration={fd gluon,amplitude=3.2pt,segment length=4.6pt,aspect=0.7}},
  gluon loop/.style={draw,solid,line width=0.75pt,decorate,
      decoration={fd gluon loop,amplitude=3.2pt,segment length=4.6pt,aspect=0.7,coil multiple=#1}},
  gluon loop/.default=1,
  photon/.style={draw,solid,line width=0.8pt,decorate,
      decoration={fd photon,amplitude=2pt,segment length=6.6pt}},
  ghost/.style={draw,line width=1.3pt,line cap=round,dash pattern=on 0pt off 2.6pt,fd arrow=#1},
  ghost/.default=0.5,
  mom pos/.store in=\fdmompos,mom pos=0.5,
  mom length/.store in=\fdmomlen,mom length=0.55cm,
  mom sep/.store in=\fdmomsep,mom sep=0.25cm,
  mom arrow/.style={draw,solid,line width=0.5pt,-{Stealth[length=3.8pt,width=3.2pt]}},
  mom/.style={decorate,decoration={markings,mark=at position \fdmompos with {%
      \draw[mom arrow] (-0.5*\fdmomlen,\fdmomsep) -- node[auto,inner sep=1.5pt]{#1} (0.5*\fdmomlen,\fdmomsep);}}},
  mom'/.style={decorate,decoration={markings,mark=at position \fdmompos with {%
      \draw[mom arrow] (-0.5*\fdmomlen,-\fdmomsep) -- node[auto,swap,inner sep=1.5pt]{#1} (0.5*\fdmomlen,-\fdmomsep);}}},
  vertex/.style={circle,fill,draw=none,inner sep=0pt,minimum size=3.8pt},
  blob/.style={draw,solid,fill=gray!30,line width=0.8pt},
  bound state/.style={draw,solid,double,double distance=1.6pt,line width=0.6pt},
}
%%% END STYLES

%%% BEGIN STYLES-EXT
% =====================================================================
%  Extension of the plain-TikZ Feynman-diagram library (fd_styles.tex)
%  for the factorization monograph.  Requires fd_styles.tex loaded first.
%  Extra libraries: shapes.geometric, patterns, backgrounds, fit.
%  New user-level styles:
%    wilson[=pos]   Wilson (eikonal) line: double line with centred arrow
%    wilson plain   Wilson line without arrow
%    cut            final-state cut: dashed red vertical line
%    hard blob      hard subgraph (grey);   jet blob (A side, blue);
%    jetB blob      collinear subgraph of hadron B (orange);
%    soft blob      soft subgraph (green);  hadron blob (target, light grey)
%    glauber        Glauber gluon (dark-red coil)
%    opvertex       operator insertion vertex (circled cross), node style
%    scalar         dashed line (e.g. scalar/longitudinal polarization)
%    longit         gluon carrying k^mu (scalar) polarization: blue coil;
%                   optionally mark it with \longitmark{(A)}{(B)} (dot at midpoint)
%    lcaxis         thin arrow for light-cone / spacetime axes
%    region         dashed rounded rectangle marking a momentum region
% Colours are kept light so that the figures print well in greyscale.
% =====================================================================
\tikzset{
  wilson/.style={draw,solid,double,double distance=1.6pt,line width=0.55pt,
      postaction={decorate,solid,decoration={markings,mark=at position #1 with
      {\arrow[xshift=3.2pt,line width=0.6pt]{Stealth[length=6.4pt,width=6pt]}}}}},
  wilson/.default=0.5,
  wilson plain/.style={draw,solid,double,double distance=1.6pt,line width=0.55pt},
  cut/.style={draw,dashed,red!70!black,line width=0.9pt,dash pattern=on 4pt off 2.5pt},
  hard blob/.style={draw,solid,fill=gray!35,line width=0.8pt},
  jet blob/.style={draw,solid,fill=blue!12,line width=0.8pt},
  jetB blob/.style={draw,solid,fill=orange!18,line width=0.8pt},
  soft blob/.style={draw,solid,fill=green!15,line width=0.8pt},
  hadron blob/.style={draw,solid,fill=gray!15,line width=0.8pt},
  glauber/.style={draw,solid,line width=0.75pt,color=red!60!black,decorate,
      decoration={fd gluon,amplitude=3.2pt,segment length=4.6pt,aspect=0.7}},
  opvertex/.style={circle,draw,solid,fill=white,line width=0.8pt,inner sep=0pt,minimum size=7pt,
      path picture={\draw[line width=0.6pt] (path picture bounding box.north west)--(path picture bounding box.south east)
                    (path picture bounding box.south west)--(path picture bounding box.north east);}},
  scalar/.style={draw,solid,line width=0.8pt,dash pattern=on 3pt off 2pt},
  longit/.style={draw,solid,line width=0.75pt,color=blue!60!black,decorate,
      decoration={fd gluon,amplitude=3.2pt,segment length=4.6pt,aspect=0.7}},
  lcaxis/.style={draw,solid,line width=0.5pt,-{Stealth[length=4pt,width=3.4pt]}},
  region/.style={draw,dashed,rounded corners=3pt,line width=0.6pt,gray!70!black},
}
\newcommand{\longitmark}[2]{\fill[blue!60!black] ($#1!0.5!#2$) circle[radius=1.7pt];}
%%% END STYLES-EXT


\begin{document}
\chapter{TikZ 图集 D 测试}
本测试文档逐页排出 \texttt{figset\_D.tex} 中的插图（每页一幅）。交叉引用测试：\cref{fig:heavy-bdpi}、\cref{fig:heavy-qcdf}、\cref{fig:power-borel-plane}、\cref{fig:lamet-ioffe}、\cref{fig:lamet-lattice-geometries}、\cref{fig:lamet-oneloop}、\cref{fig:lattice-matching}、\cref{fig:lamet-dependency}、\cref{fig:lamet-hybrid}、\cref{fig:lamet-fourier}、\cref{fig:lamet-validity}、\cref{fig:yzliu-soft-hqet}、\cref{fig:yzliu-single-line}、\cref{fig:yzliu-theta-strip}、\cref{fig:yzliu-complex-lorentz}、\cref{fig:yzliu-leading-region}、\cref{fig:yzliu-threshold}、\cref{fig:yzliu-borel}、\cref{fig:summary-status-map}、\cref{fig:summary-hypotheses}。
\clearpage
\begin{figure}[htbp]
\centering
\begin{tikzpicture}[fd,
  fd2 heavy/.style={draw,solid,line width=1.3pt,fd arrow=#1},
  fd2 heavy/.default=0.5,
  fd2 meson/.style={bound state,line cap=butt},
  fd2 weak/.style={rectangle,fill=black,draw=none,inner sep=0pt,minimum size=5.6pt},
  fd2 soft/.style={gluon,draw=green!45!black},
  fd2 cross/.pic={\draw[red!75!black,line width=1.1pt,line cap=round] (-0.12,-0.12)--(0.12,0.12) (-0.12,0.12)--(0.12,-0.12);},
  fd2 bracket/.style={draw=gray!55!black,line width=0.5pt},
  fd2 lab/.style={font=\small},
  fd2 note/.style={font=\footnotesize}]
  \coordinate (W) at (0,0.45);
  \coordinate (dq) at (1.449,2.820);   % end of the d line (inside the pion blob)
  \coordinate (uq) at (1.851,2.558);   % end of the u-bar line (inside the pion blob)
  % ---- mesons (double lines) ----
  \draw[fd2 meson] (-5.15,0) -- (-3.9,0);
  \draw[fd2 meson] (3.9,0) -- (5.15,0);
  \draw[fd2 meson] (2.36,3.78) -- (2.84,4.51);
  % ---- soft gluons from the B-D system to the pion pair (decouple) ----
  \coordinate (sd) at ($(W)!0.42!(dq)$);
  \coordinate (su) at ($(W)!0.36!(uq)$);
  \draw[fd2 soft] (-3.42,0.72) -- (sd);
  \draw[fd2 soft] (3.42,0.72) -- (su);
  \pic at ($(-3.42,0.72)!0.5!(sd)$) {fd2 cross};
  \pic at ($(3.42,0.72)!0.5!(su)$) {fd2 cross};
  % ---- quark lines (arrows = fermion-number flow) ----
  \draw[fd2 heavy=0.55] (-3.4,0.45) -- (W);
  \draw[fd2 heavy=0.55] (W) -- (3.4,0.45);
  \draw[quark] (3.4,-0.45) -- (-3.4,-0.45);
  \draw[quark=0.72] (W) -- (dq);
  \draw[quark=0.6] (uq) -- (W);
  % ---- blobs on top ----
  \node[hadron blob,ellipse,minimum width=0.8cm,minimum height=1.8cm,inner sep=0pt] at (-3.6,0) {};
  \node[hadron blob,ellipse,minimum width=0.8cm,minimum height=1.8cm,inner sep=0pt] at (3.6,0) {};
  \node[jet blob,ellipse,minimum width=1.6cm,minimum height=0.72cm,inner sep=0pt,rotate=57] at (1.95,3.15) {};
  \node[fd2 weak] at (W) {};
  % ---- labels ----
  \node[fd2 lab,left] at (-5.15,0) {$\bar B^0$};
  \node[fd2 lab,right] at (5.15,0) {$D^+$};
  \node[fd2 lab,anchor=south west] at (2.78,4.45) {$\pi^-(p)$};
  \node[fd2 lab] at (-1.9,0.2) {$b$};
  \node[fd2 lab] at (1.9,0.2) {$c$};
  \node[fd2 lab] at (-1.9,-0.72) {$\bar d$};
  \path[mom={$xp$},mom sep=0.25cm,mom length=0.5cm,mom pos=0.8] (W) -- (dq);
  \path[mom'={$(1-x)p$},mom sep=0.25cm,mom length=0.5cm,mom pos=0.84] (W) -- (uq);
  \node[fd2 lab,anchor=east] at (1.28,3.42) {$\phi_\pi(x)$};
  \node[fd2 note,anchor=north west] at (0.14,0.36) {$T(x,\mu)$};
  \node[fd2 note,anchor=south west,align=left] at (3.7,0.98) {软胶子在领头幂次\\退耦（色透明性）};
  % ---- the soft B -> D transition ----
  \draw[fd2 bracket] (-4.02,-1.05) -- (-4.02,-1.15) -- (4.02,-1.15) -- (4.02,-1.05);
  \node[fd2 note,below=1pt] at (0,-1.15) {软的 $B\to D$ 跃迁（含旁观者）：形状因子 $F^{B\to D}(0)$};
  \node[fd2 lab] at (0,-2.2) {$\langle D^+\pi^-|O_{\mathbf 0}|\bar B^0\rangle=N\,F^{B\to D}(0)\displaystyle\int_0^1\!\dd x\,T(x,\mu)\,\phi_\pi(x,\mu)\,\big[1+\order{\Lambda/m_b}\big]$};
\end{tikzpicture}
\caption{$\bar B^0\to D^+\pi^-$ 的因子化（\cref{thm:heavy-bdpi}，全阶论证 [A]）。黑方块为弱作用四夸克算符 $O_{\mathbf 0}=[\bar c\gamma^\mu P_Lb][\bar d\gamma_\mu P_Lu]$ 的顶点，粗线为重夸克 $b$、$c$，双线为介子；细线上的箭头为费米子数流向（反夸克 $\bar d$、$\bar u$ 的箭头逆着动量）。
旁观者 $\bar d$ 自 $\bar B^0$ 进入 $D^+$，这一软的跃迁由形状因子 $F^{B\to D}(0)$ 描述；发射的 $\pi^-$ 由弱顶点直接产生的共线 $d\bar u$ 对构成（蓝色团，分布振幅 $\phi_\pi(x)$，$x$ 为 $d$ 夸克的动量分数），它是快速运动的色单态小偶极子：
软胶子与 $d$、$\bar u$ 的程函耦合在领头幂次相消（打叉的绿色胶子，色透明性），硬胶子交换归入短程系数 $T(x,\mu)$。
$N=\frac{\ii}{2}m_BE_\pi f_\pi$；修正被 $\Lambda/m_b$、$\Lambda/m_c$、$\Lambda/E_\pi$ 压低（VF-BQ-6）；$B^-\to D^0\pi^-$ 同理。证明还用到分布振幅端点压低的假设（见正文）。}
\label{fig:heavy-bdpi}
\end{figure}
\clearpage
\begin{figure}[htbp]
\centering
\begin{tikzpicture}[fd,
  fd2 heavy/.style={draw,solid,line width=1.3pt,fd arrow=#1},
  fd2 heavy/.default=0.5,
  fd2 meson/.style={bound state,line cap=butt},
  fd2 weak/.style={rectangle,fill=black,draw=none,inner sep=0pt,minimum size=5.6pt},
  fd2 halo/.style={preaction={draw=white,line width=3.4pt}},
  fd2 endpoint/.style={draw=red!80!black,line width=3.2pt,opacity=0.45,line cap=butt},
  fd2 bracket/.style={draw=gray!55!black,line width=0.5pt},
  fd2 lab/.style={font=\small},
  fd2 note/.style={font=\footnotesize}]
  % ------------- common skeleton: B (left), M1 (right), M2 (emitted, top) -------------
  \foreach \panel/\xs in {a/0, b/8.0} {
  \begin{scope}[xshift=\xs cm]
    \coordinate (W) at (-0.5,0.42);
    \coordinate (q2) at (0.352,2.160);   % quark of M2 (inside the blob)
    \coordinate (a2) at (0.698,1.960);   % antiquark of M2 (inside the blob)
    \draw[fd2 meson] (-3.3,0) -- (-2.75,0);
    \draw[fd2 meson] (2.75,0) -- (3.3,0);
    \draw[fd2 meson] (1.08,3.03) -- (1.43,3.64);
    \if a\panel
      % hard gluon of a vertex correction: M2 antiquark -- light quark q
      \coordinate (g1) at (1.2,0.42);  \coordinate (g2) at ($(W)!0.58!(a2)$);
      \draw[gluon] (g1) -- (g2);
      \draw[fd2 heavy=0.5] (-2.3,0.42) -- (W);
      \draw[quark=0.78] (W) -- (2.3,0.42);
      \draw[quark] (2.3,-0.42) -- (-2.3,-0.42);
      \node[vertex] at (g1) {};  \node[vertex] at (g2) {};
    \else
      % hard-collinear gluon from the spectator to the emitted pair
      \coordinate (g1) at (0.85,-0.42);  \coordinate (g2) at ($(W)!0.52!(q2)$);
      \draw[fd2 endpoint] (0.9,-0.42) -- (2.2,-0.42);
      \draw[gluon] (g1) -- (g2);
      \draw[fd2 heavy=0.5] (-2.3,0.42) -- (W);
      \draw[quark=0.8,fd2 halo] (W) -- (2.3,0.42);
      \draw[quark=0.55] (2.3,-0.42) -- (g1);
      \draw[quark=0.45] (g1) -- (-2.3,-0.42);
      \node[vertex] at (g1) {};  \node[vertex] at (g2) {};
    \fi
    \draw[quark=0.33] (W) -- (q2);
    \draw[quark=0.6] (a2) -- (W);
    \node[hadron blob,ellipse,minimum width=0.66cm,minimum height=1.56cm,inner sep=0pt] at (-2.45,0) {};
    \node[jetB blob,ellipse,minimum width=0.66cm,minimum height=1.56cm,inner sep=0pt] at (2.45,0) {};
    \node[jet blob,ellipse,minimum width=1.35cm,minimum height=0.62cm,inner sep=0pt,rotate=60] at (0.75,2.45) {};
    \node[fd2 weak] at (W) {};
    \node[fd2 lab,left] at (-3.3,0) {$\bar B$};
    \node[fd2 lab,right] at (3.3,0) {$M_1$};
    \node[fd2 lab,anchor=south west] at (1.36,3.55) {$M_2$};
    \node[fd2 lab] at (-1.45,0.66) {$b$};
    \node[fd2 lab] at (1.5,0.66) {$q$};
    \node[fd2 lab,anchor=east] at (0.3,2.72) {$\phi_{M_2}(x)$};
  \end{scope}}
  % ------------- panel (a): labels -------------
  \node[fd2 lab] at (-1.45,-0.66) {$\bar q_s$};
  \node[fd2 note,anchor=west] at (1.05,1.25) {$T^{\rm I}$：硬胶子};
  \draw[fd2 bracket] (-2.82,-1.0) -- (-2.82,-1.1) -- (2.82,-1.1) -- (2.82,-1.0);
  \node[fd2 note,below=1pt] at (0,-1.1) {软的 $B\to M_1$ 跃迁：$F^{B\to M_1}(0)$};
  \node[fd2 lab,align=center] at (0,-2.3) {(a)\ 顶点项\\[2pt]$F^{B\to M_1}(0)\displaystyle\int\!\dd x\,T^{\rm I}(x)\,\phi_{M_2}(x)$};
  % ------------- panel (b): labels -------------
  \begin{scope}[xshift=8.0cm]
    \path[mom={$\xi p_B$},mom sep=0.2cm,mom length=0.5cm,mom pos=0.4] (-2.2,-0.42) -- (0.85,-0.42);
    \path[mom={$\bar y\,p_1$},mom sep=0.2cm,mom length=0.45cm,mom pos=0.52] (0.85,-0.42) -- (2.2,-0.42);
    \node[fd2 note,anchor=south east] at (-0.12,1.06) {$T^{\rm II}$};
    \node[fd2 note,anchor=west] at (0.56,1.02) {$k^2\sim m_b\Lambda$};
    \node[fd2 note] at (-2.45,-1.1) {$\phi_B(\xi)$};
    \node[fd2 note] at (2.6,-1.12) {$\phi_{M_1}(y)$};
    \node[fd2 note,text=red!70!black] at (1.05,-0.74) {端点区 $\bar y\to0$};
    \node[fd2 lab,align=center] at (0,-2.3) {(b)\ 硬旁观者散射\\[2pt]$\displaystyle\int\!\dd\xi\,\dd x\,\dd y\,T^{\rm II}(\xi,x,y)\,\phi_B(\xi)\,\phi_{M_1}(y)\,\phi_{M_2}(x)$};
  \end{scope}
\end{tikzpicture}
\caption{无粲两体衰变 $\bar B\to M_1M_2$ 的 QCD 因子化（BBNS；\cref{prop:heavy-qcdf-lp}，部分 [P]），各只画一个代表图。黑方块为弱作用四夸克算符的顶点，粗线为 $b$ 夸克，箭头为费米子数流向；
$M_2$（蓝）为由弱顶点发射的介子，$M_1$（橙）为接收旁观反夸克 $\bar q_s$ 的介子。(a)~形状因子项：硬胶子连接发射对与 $b\to q$ 线（顶点修正），软的跃迁由 $F^{B\to M_1}(0)$ 描述。
(b)~硬旁观者散射：旁观者（动量 $\xi p_B$，$\xi\sim\Lambda/m_b$）放出一个硬共线胶子（虚度 $\sim m_b\Lambda$），以动量分数 $\bar y=1-y$ 进入 $M_1$。
红色标出端点区 $\bar y\to0$（旁观反夸克以趋于零的动量分数进入 $M_1$）：扭度 3 的硬旁观者项在此对数发散，$X_H=\int_0^1\frac{\dd y}{1-y}\Phi_p(y)$；弱湮灭项（未画出）有同类端点发散 $X_A=\int_0^1\frac{\dd y}{y}$（VF-BQ-3）。
这些幂次修正的标准共线因子化不成立（\cref{prop:heavy-qcdf-endpoint}，[V]），$X_{A,H}$ 的参数化只是模型。}
\label{fig:heavy-qcdf}
\end{figure}
\clearpage
\begin{figure}[htbp]
\centering
\begin{tikzpicture}[fd,x=1.3cm,y=1cm,
  fd2 pole/.pic={\draw[line width=1.05pt,line cap=round] (-0.09,-0.09)--(0.09,0.09) (-0.09,0.09)--(0.09,-0.09);},
  fd2 ax/.style={lcaxis,gray!60!black},
  fd2 contour/.style={draw=#1,line width=1.0pt,postaction={decorate,decoration={markings,
      mark=at position 0.5 with {\arrow{Stealth[length=5.5pt,width=4.6pt]}}}}},
  fd2 pointer/.style={draw=#1,line width=0.45pt,-{Stealth[length=3.5pt,width=3pt]}},
  fd2 note/.style={font=\footnotesize},
  fd2 lab/.style={font=\small}]
  % axes
  \draw[fd2 ax] (-3.75,0) -- (6.05,0) node[right] {$\Real u$};
  \draw[fd2 ax] (0,-0.75) -- (0,1.75) node[left] {$\Imag u$};
  \foreach \u in {-3,-2,-1,1,2,3,4,5} \draw[gray!60!black,line width=0.4pt] (\u,-0.05) -- (\u,0.05);
  \foreach \u in {-3,-2,-1,1,2,3,4,5} \node[fd2 note] at (\u,-0.5) {$\u$};
  \node[fd2 note,anchor=north east] at (-0.03,-0.04) {$0$};
  % Borel integration paths above and below the positive axis
  \draw[fd2 contour=blue!60!black] (0,0) .. controls (0.3,0.22) and (0.55,0.24) .. (0.9,0.24) -- (5.5,0.24);
  \draw[fd2 contour=blue!60!black] (0,0) .. controls (0.3,-0.22) and (0.55,-0.24) .. (0.9,-0.24) -- (5.5,-0.24);
  \node[fd2 lab,text=blue!50!black,anchor=west] at (5.5,0.28) {$C_+$};
  \node[fd2 lab,text=blue!50!black,anchor=west] at (5.5,-0.28) {$C_-$};
  % IR renormalons (u = 2,3,...), UV renormalons (u = -1,-2,...), no pole at u = 1
  \foreach \u in {2,3,4,5} \pic[red!70!black] at (\u,0) {fd2 pole};
  \foreach \u in {-1,-2,-3} \pic[gray!55!black] at (\u,0) {fd2 pole};
  \draw[line width=0.9pt,fill=white] (1,0) circle[radius=0.085];
  % annotations
  \node[fd2 note,text=red!60!black] at (3.5,0.72) {红外重整化子：$u=2,3,\dots$};
  \node[fd2 note,align=center,text=gray!40!black] at (-1.95,0.82) {紫外重整化子：$u=-1,-2,\dots$\\（符号交错，不在积分路径上）};
  \draw[fd2 pointer=gray!60!black] (0.78,-1.02) -- (0.97,-0.12);
  \node[fd2 note,anchor=north east] at (0.95,-0.98) {无 $u=1$ 极点：不存在量纲为 2 的规范不变算符};
  \draw[fd2 pointer=red!60!black] (2.12,-1.0) -- (2.02,-0.12);
  \node[fd2 note,align=left,anchor=north west,text=red!50!black] at (1.75,-0.96) {$u=2$：沿 $C_+$ 与 $C_-$ 积分之差 $\propto(\Lambda^2/Q^2)^2$，\\对应 OPE 中的凝聚项 $\langle\frac{\alpha_s}{\pi}G^2\rangle/Q^4$};
  \node[fd2 lab] at (1.15,-2.7) {$D(Q^2)=\displaystyle\int_0^\infty\!\dd t\,\ee^{-t/\alpha_s(Q)}B[D](t)$，\quad $u=\dfrac{\beta_0t}{4\pi}$，\quad $\ee^{-t/\alpha_s(Q)}=\Big(\dfrac{\Lambda^2}{Q^2}\Big)^{u}$};
\end{tikzpicture}
\caption{大 $\beta_0$（泡链）近似下 Adler 函数 $D(Q^2)$ 的 Borel 平面（\cref{prop:power-adler-large-beta0}，部分 [P]；Borel 变换按\cref{def:power-borel}：$R=\sum_nr_n\alpha_s^{n+1}$，$B[R](t)=\sum_nr_nt^n/n!$，$u=\beta_0t/(4\pi)$，单圈跑动下 $\ee^{-t/\alpha_s(Q)}=(\Lambda^2/Q^2)^u$）。
红外重整化子极点位于正实轴 $u=2,3,\dots$（红叉），使沿正实轴的 Borel 积分不确定：绕到极点上方（$C_+$）或下方（$C_-$）的两种积分之差在 $u=u_0$ 处正比于 $(\Lambda^2/Q^2)^{u_0}$。
$u=1$ 处没有极点（空心圆），这与不存在量纲为 2 的规范不变算符相一致；$u=2$ 的极点与 OPE 中的胶子凝聚项 $\langle\frac{\alpha_s}{\pi}G^2\rangle/Q^4$ 相对应。紫外重整化子位于 $u=-1,-2,\dots$（灰叉）。
只有大 $\beta_0$ 极限中的这一结构是确立的；红外重整化子与高扭度之间的一般对应是猜想（\cref{conj:power-renormalon-ht}，[C]）。}
\label{fig:power-borel-plane}
\end{figure}
\clearpage
\begin{figure}[htbp]
\centering
\begin{tikzpicture}
\begin{axis}[width=0.8\textwidth,height=6.4cm,xmin=0,xmax=16.6,ymin=0,ymax=1.25,
  axis lines=left,axis line style={gray!60!black},clip=false,
  xlabel={Ioffe 时间 $\lambda=P\cdot z$},ylabel={$|z|/\mathrm{fm}$},
  xtick={0,2,...,16},ytick={0,0.2,...,1.2},
  yticklabel style={/pgf/number format/fixed,/pgf/number format/fixed zerofill,/pgf/number format/precision=1},
  tick label style={font=\scriptsize},label style={font=\small},
  every axis x label/.style={at={(axis description cs:0.5,-0.13)},anchor=north},
  every axis y label/.style={at={(axis description cs:-0.075,0.5)},rotate=90,anchor=south}]
  % lattice data region (schematic) and short-distance band
  \fill[gray!9] (axis cs:0,0) rectangle (axis cs:15,1.0);
  \draw[gray!50,line width=0.4pt,dash pattern=on 3pt off 2pt] (axis cs:0,1.0) -- (axis cs:15,1.0) -- (axis cs:15,0);
  \fill[blue!11] (axis cs:0,0) rectangle (axis cs:16.6,0.3);
  % pseudo-distributions: fixed z^2
  \draw[black!75,line width=0.7pt,dotted] (axis cs:0,0.1) -- (axis cs:16.6,0.1);
  \draw[black!75,line width=0.7pt,dotted] (axis cs:0,0.2) -- (axis cs:16.6,0.2);
  \draw[black!75,line width=0.7pt,dotted] (axis cs:0,0.3) -- (axis cs:16.6,0.3);
  % quasi-distributions: fixed P^z, |z| = 0.1973 lambda/(P^z/GeV) fm
  \addplot[red!65!black,line width=0.9pt,dashed,domain=0:6.335,samples=2] {0.1973*x/1};
  \addplot[red!65!black,line width=0.9pt,dashed,domain=0:12.67,samples=2] {0.1973*x/2};
  \addplot[red!65!black,line width=0.9pt,dashed,domain=0:16.6,samples=2] {0.1973*x/3};
  % light-cone ITD on the lambda axis
  \draw[blue!55!black,line width=2.2pt,line cap=butt] (axis cs:0,0) -- (axis cs:16.6,0);
  % labels
  \node[font=\scriptsize,text=red!55!black,anchor=south] at (axis cs:6.335,1.25) {$P^z=1$ GeV};
  \node[font=\scriptsize,text=red!55!black,anchor=south] at (axis cs:12.67,1.25) {$P^z=2$ GeV};
  \node[font=\scriptsize,text=red!55!black,anchor=west] at (axis cs:16.6,1.092) {$P^z=3$ GeV};
  \node[font=\footnotesize,text=red!50!black,align=left,anchor=west,inner sep=1.2pt] at (axis cs:0.6,1.135) {准分布：固定 $P^z$\\[-1pt]（射线 $|z|=|\lambda|/P^z$）};
  \draw[red!55!black,line width=0.55pt,-{Stealth[length=4.5pt,width=3.8pt]}] (axis cs:6.9,0.97) to[bend left=14] node[pos=0.55,above right=-1pt,font=\footnotesize] {$P^z\to\infty$} (axis cs:13.4,0.6);
  \node[font=\footnotesize,anchor=west] at (axis cs:16.6,0.3) {$|z|=0.3$ fm};
  \node[font=\footnotesize,anchor=west] at (axis cs:16.6,0.2) {$|z|=0.2$ fm};
  \node[font=\footnotesize,anchor=west] at (axis cs:16.6,0.1) {$|z|=0.1$ fm};
  \node[font=\footnotesize,anchor=west,align=left] at (axis cs:16.6,0.52) {赝分布：\\[-1pt]固定 $z^2$};
  \node[font=\footnotesize,text=blue!45!black,anchor=north west,align=left] at (axis cs:16.6,-0.02) {光锥 ITD $I(\lambda)$\\[-1pt]（$z^2=0$）};
  \node[font=\footnotesize,fill=blue!11,inner sep=1.5pt] at (axis cs:9.6,0.15) {短距离因子化（SDF）};
  \node[font=\footnotesize,text=gray!45!black,anchor=south east] at (axis cs:14.85,0.33) {格点数据区（示意）};
\end{axis}
\end{tikzpicture}
\caption{Ioffe 时间 $\lambda=P\cdot z$（$=-P^zz$，等时类空间隔 $z^\mu=z\hat e_z^\mu$）与间隔 $|z|$ 平面上的三种表示（图中只画 $\lambda\ge0$；$\hbar c=0.1973$ GeV\,fm）。光锥 ITD $I(\lambda)=\int_{-1}^1\dd x\,\ee^{\ii x\lambda}f(x)$ 位于 $z^2=0$（粗线）；
固定 $z^2$ 的水平线（点线）给出赝分布：对 $\lambda$ 作 Fourier 变换，支集为 $[-1,1]$；固定 $P^z$ 的射线 $|z|=|\lambda|/P^z$（虚线，$P^z=1,2,3$ GeV）给出准分布：$z$ 与 $\lambda$ 同时变化，$P^z$ 增大时射线趋向光锥轴。
短距离因子化只用阴影带（$|z|\lesssim0.3$ fm）内的数据，准分布的 Fourier 变换则需要整条射线（长距离部分见正文）；灰色矩形示意格点数据所及的范围。}
\label{fig:lamet-ioffe}
\end{figure}
\clearpage
\begin{figure}[htbp]
\centering
\begin{tikzpicture}[fd,
  fd2 grid/.style={draw=gray!30,line width=0.3pt},
  fd2 warrow/.style={postaction={decorate,decoration={markings,
      mark=at position #1 with {\arrow[xshift=3.2pt]{Stealth[length=6.4pt,width=6pt]}}}}},
  fd2 site/.style={circle,fill=black,draw=none,inner sep=0pt,minimum size=5pt},
  fd2 osite/.style={circle,draw=black,fill=white,line width=0.8pt,inner sep=0pt,minimum size=5pt},
  fd2 dim/.style={draw=gray!55!black,line width=0.5pt,{Stealth[length=3.2pt,width=2.8pt]}-{Stealth[length=3.2pt,width=2.8pt]}},
  fd2 ax/.style={lcaxis,gray!60!black},
  fd2 lab/.style={font=\small,fill=white,inner sep=1pt},
  fd2 note/.style={font=\footnotesize,fill=white,inner sep=1pt}]
  % ======================= (a) quasi-PDF operator =======================
  \begin{scope}
    \draw[fd2 grid,step=0.3] (-0.6,-0.6) grid (2.4,1.8);
    \draw[wilson plain] (0,0) -- (1.8,0);
    \path[fd2 warrow=0.56] (0,0) -- (1.8,0);
    \node[fd2 site] at (0,0) {};  \node[fd2 osite] at (1.8,0) {};
    \node[fd2 lab,anchor=north] at (0,-0.12) {$\psi(0)$};
    \node[fd2 lab,anchor=north] at (1.8,-0.12) {$\bar\psi(z\hat e_z)$};
    \draw[fd2 dim] (0,0.36) -- node[fd2 note,above=1pt] {$z$} (1.8,0.36);
    \draw[fd2 dim] (2.1,1.32) -- node[fd2 note,above=1pt] {$a$} (2.4,1.32);
    \draw[fd2 ax] (-0.45,1.05) -- (0.2,1.05) node[fd2 note,right=0pt] {$\hat z$};
    \draw[fd2 ax] (-0.45,1.05) -- (-0.45,1.62) node[fd2 note,right=1pt] {$\vec b_\perp$};
    \node[font=\small] at (0.9,-1.15) {$\bar\psi(z\hat e_z)\,\Gamma\,W[z\hat e_z,0]\,\psi(0)$};
    \node[font=\small] at (0.9,-1.75) {(a)\ 准 PDF：直 Wilson 线};
  \end{scope}
  % ======================= (b) quasi-TMD operator: space-like staple =======================
  \begin{scope}[xshift=4.4cm]
    \draw[fd2 grid,step=0.3] (-0.6,-0.6) grid (2.7,1.8);
    \draw[wilson plain] (0,0) -- (2.1,0) -- (2.1,1.2) -- (0.9,1.2);
    \path[fd2 warrow=0.52] (0,0) -- (2.1,0);
    \path[fd2 warrow=0.55] (2.1,0) -- (2.1,1.2);
    \path[fd2 warrow=0.6] (2.1,1.2) -- (0.9,1.2);
    \node[fd2 site] at (0,0) {};  \node[fd2 osite] at (0.9,1.2) {};
    \node[fd2 lab,anchor=north] at (0,-0.12) {$\psi(0)$};
    \node[fd2 lab,anchor=south] at (0.75,1.33) {$\bar\psi(z\hat e_z+\vec b_\perp)$};
    \draw[fd2 dim] (0,0.36) -- node[fd2 note,above=1pt] {$L$} (2.1,0.36);
    \draw[fd2 dim] (-0.33,0) -- node[fd2 note,left=1pt] {$b_\perp$} (-0.33,1.2);
    \draw[red!60!black,line width=0.6pt,-{Stealth[length=4.5pt,width=3.8pt]}] (2.25,0.6) -- (2.85,0.6) node[font=\footnotesize,text=red!55!black,right=1pt] {$L\to\infty$};
    \node[font=\small] at (1.05,-1.15) {$\bar\psi(z\hat e_z+\vec b_\perp)\,\Gamma\,\mathcal W_z\,\psi(0)$};
    \node[font=\small] at (1.05,-1.75) {(b)\ 准 TMD：类空订书钉};
  \end{scope}
  % ======================= (c) rectangular Wilson loop =======================
  \begin{scope}[xshift=9.45cm]
    \draw[fd2 grid,step=0.3] (-0.6,-0.6) grid (4.8,1.8);
    \draw[wilson plain] (0,0) rectangle (4.2,1.2);
    \path[fd2 warrow=0.5] (0,0) -- (4.2,0);
    \path[fd2 warrow=0.55] (4.2,0) -- (4.2,1.2);
    \path[fd2 warrow=0.5] (4.2,1.2) -- (0,1.2);
    \path[fd2 warrow=0.55] (0,1.2) -- (0,0);
    \draw[fd2 dim] (0,-0.36) -- (4.2,-0.36);
    \node[fd2 note] at (2.1,-0.36) {$2L$};
    \draw[fd2 dim] (-0.33,0) -- node[fd2 note,left=1pt] {$b_\perp$} (-0.33,1.2);
    \node[font=\small,fill=white,inner sep=1.2pt] at (2.1,0.6) {$Z_E(2L,b_\perp)$};
    \node[font=\small] at (2.1,-1.15) {$\frac1{N_c}\Tr\langle0|W_\perp\,\mathcal W_z(\vec b_\perp;2L)|0\rangle$};
    \node[font=\small] at (2.1,-1.75) {(c)\ 矩形 Wilson 圈};
  \end{scope}
\end{tikzpicture}
\caption{格点上的算符几何（等时面 $t=0$，晶格间距 $a$；实心点为 $\psi$、空心点为 $\bar\psi$ 所在格点，双线为规范链，箭头为路径方向，即 $W[b,a]$ 由 $a$ 指向 $b$）：
(a)~准 PDF 算符的直 Wilson 线 $W[z\hat e_z,0]$；(b)~准 TMD 算符的订书钉形 Wilson 线 $\mathcal W_z$：自 $\psi$ 沿 $+\hat z$ 走长度 $L$，经横向连接（间隔 $b_\perp$），再沿 $-\hat z$ 回到 $\bar\psi$（图示为 DY 型；SIDIS 型沿 $-\hat z$）；
(c)~矩形 Wilson 圈 $Z_E(2L,b_\perp)$，$\sqrt{Z_E}$ 除去订书钉的线性发散、尖点发散与捏点奇点，使 $L\to\infty$ 的极限存在。手征对称性不严格的格点费米子还引起同维混合 $\Gamma\to\Gamma'$（见正文）。
注意：这里的订书钉是类空的，与\cref{fig:tmd-staple}中 Collins TMD 的类光订书钉不同。}
\label{fig:lamet-lattice-geometries}
\end{figure}
\clearpage
\begin{figure}[htbp]
\centering
\begin{tikzpicture}[fd,
  fd2 light/.style={draw=gray!60},
  fd2 bracket/.style={draw=gray!55!black,line width=0.5pt},
  fd2 lab/.style={font=\small},
  fd2 note/.style={font=\footnotesize}]
  \foreach \panel/\xs/\wpos/\ttl in {a/0/0.5/{(a)\ 顶点}, b/3.1/0.32/{(b)\ 帆}, c/6.2/0.68/{(c)\ 帆（镜像）}, d/9.3/0.5/{(d)\ Wilson 线自能}, e/12.4/0.5/{(e)\ 外线自能}} {
  \begin{scope}[xshift=\xs cm]
    \coordinate (O1) at (-0.85,0);   \coordinate (O2) at (0.85,0);
    \coordinate (I) at (-1.35,-1.35); \coordinate (F) at (1.35,-1.35);
    \coordinate (i1) at ($(I)!0.5!(O1)$); \coordinate (f1) at ($(O2)!0.5!(F)$);
    \if a\panel \draw[gluon] (i1) -- (f1); \fi
    \if b\panel \draw[gluon] (i1) to[bend right=28] (0.32,0); \fi
    \if c\panel \draw[gluon] (-0.32,0) to[bend right=28] (f1); \fi
    \if d\panel \draw[gluon] (-0.48,0) arc[start angle=180,end angle=360,radius=0.48]; \fi
    \if e\panel
      \draw[gluon,fd2 light] ($(I)!0.22!(O1)$) to[bend left=80] ($(I)!0.78!(O1)$);
    \fi
    \draw[wilson=\wpos] (O1) -- (O2);
    \if e\panel
      \draw[quark=0.88,fd2 light] (I) -- (O1);
    \else
      \draw[quark=0.3] (I) -- (O1);
    \fi
    \draw[quark=0.7] (O2) -- (F);
    \if a\panel \node[vertex] at (i1) {}; \node[vertex] at (f1) {}; \fi
    \if b\panel \node[vertex] at (i1) {}; \node[vertex] at (0.32,0) {}; \fi
    \if c\panel \node[vertex] at (-0.32,0) {}; \node[vertex] at (f1) {}; \fi
    \if d\panel \node[vertex] at (-0.48,0) {}; \node[vertex] at (0.48,0) {}; \fi
    \if e\panel
      \node[vertex,fill=gray!60] at ($(I)!0.22!(O1)$) {}; \node[vertex,fill=gray!60] at ($(I)!0.78!(O1)$) {};
    \fi
    \node[opvertex] at (O1) {};  \node[opvertex] at (O2) {};
    \node[fd2 note,above=2pt] at (O1) {$0$};
    \node[fd2 note,above=1pt] at (O2) {$z\hat e_z$};
    \node[fd2 lab] at (0,-1.78) {\ttl};
  \end{scope}}
  % external momenta (same in all panels)
  \path[mom={$p$},mom sep=0.22cm,mom length=0.4cm,mom pos=0.24] (-1.35,-1.35) -- (-0.85,0);
  \path[mom={$p$},mom sep=0.22cm,mom length=0.4cm,mom pos=0.76] (0.85,0) -- (1.35,-1.35);
  % unphysical-region contributions (xi > 1), VF-LAM-7
  \node[fd2 lab] at (0,-2.62) {$a\displaystyle\int_0^1\!\dd y\,\frac{1-y}{|\xi-y|}$};
  \draw[fd2 bracket] (1.75,-2.08) -- (1.75,-2.16) -- (7.75,-2.16) -- (7.75,-2.08);
  \node[fd2 lab] at (4.65,-2.62) {(b)+(c)：$\dfrac{a}{1-\xi}\displaystyle\int_0^1\!\dd y\,\frac{\xi+y}{|\xi-y|}$};
  \node[fd2 lab] at (9.3,-2.62) {$-\dfrac{a}{|1-\xi|}$};
  \node[fd2 note,text=gray!45!black] at (12.4,-2.62) {（只出现在 $\delta(1-\xi)$ 中）};
\end{tikzpicture}
\caption{准 PDF 的单圈图（Feynman 规范，在壳无质量夸克 $p=(p^z,0,0,p^z)$，$\Gamma=\gamma^0$；入射与出射夸克动量同为 $p$，前向矩阵元，无割线）。双线为等时、沿 $\hat z$ 的类空 Wilson 线 $W[z\hat e_z,0]$（箭头为路径方向），$\otimes$ 为算符端点 $\psi(0)$、$\bar\psi(z\hat e_z)$：
(a)~顶点图，(b)(c)~“帆”图，(d)~Wilson 线自能，(e)~外线自能。图下为非物理区 $\xi>1$ 中各图的贡献（$y$ 为 Feynman 参数，$a=\frac{\alpha_sC_F}{2\pi}$），在 $d=4$ 中都有限，其和为
$\tilde q^{(1)}=a\big[\frac{1+\xi^2}{1-\xi}\ln\frac{\xi}{\xi-1}+1\big]$（VF-LAM-7，与\cite{IzubuchiEtAl2018}附录 B 一致）；维数正规化中 Wilson 线自能没有线性发散。}
\label{fig:lamet-oneloop}
\end{figure}
\clearpage
\begin{figure}[htbp]
\centering
\begin{tikzpicture}[fd,
  fd2 qty/.style={draw=black!70,fill=white,rounded corners=2pt,line width=0.6pt,
      minimum width=4.3cm,minimum height=0.95cm,align=center,font=\small},
  fd2 frame/.style={draw=#1!60!black,fill=#1!4,rounded corners=4pt,line width=0.6pt},
  fd2 match/.style={draw=black!80,line width=1.0pt,{Stealth[length=7pt,width=6pt]}-{Stealth[length=7pt,width=6pt]}},
  fd2 ft/.style={draw=gray!55!black,line width=0.6pt,{Stealth[length=4.5pt,width=3.8pt]}-{Stealth[length=4.5pt,width=3.8pt]}},
  fd2 note/.style={font=\footnotesize},
  fd2 lab/.style={font=\small}]
  % column frames
  \draw[fd2 frame=gray] (-7.45,-2.85) rectangle (-1.85,2.45);
  \draw[fd2 frame=blue] (1.85,-2.85) rectangle (7.45,2.45);
  \node[fd2 lab,font=\small\bfseries] at (-4.6,2.12) {格点：等时（欧氏）矩阵元};
  \node[fd2 lab,font=\small\bfseries] at (4.6,2.12) {光锥分布（$\MSbar$）};
  % quantities
  \node[fd2 qty] (q)  at (-4.6,1.05)  {准 PDF $\tilde q_\Gamma(x,P^z)$\\[-1pt]{\footnotesize 固定 $P^z$}};
  \node[fd2 qty] (h)  at (-4.6,-1.05) {准 ITD $\tilde h(\lambda,z^2)$\\[-1pt]{\footnotesize $\propto\langle P|\bar\psi(z\hat e_z)\Gamma W[z\hat e_z,0]\psi(0)|P\rangle_R$}};
  \node[fd2 qty] (f)  at (4.6,1.05)   {光锥 PDF $f(y,\mu)$\\[-1pt]{\footnotesize 支集 $[-1,1]$}};
  \node[fd2 qty] (I)  at (4.6,-1.05)  {光锥 ITD $I(\lambda,\mu)$\\[-1pt]{\footnotesize $\int_{-1}^1\dd x\,\ee^{\ii x\lambda}f(x,\mu)$}};
  \draw[fd2 ft] (q) -- node[fd2 note,right=2pt,align=left] {Fourier：沿射线\\$z=-\lambda/P^z$} (h);
  \draw[fd2 ft] (f) -- node[fd2 note,left=2pt] {Fourier} (I);
  \node[fd2 note,align=center,text=gray!30!black] at (-4.65,-2.2) {非微扰重整化（除去线性发散 $\ee^{-\delta m|z|}$）\\[-1pt]与方案选择，见\cref{fig:lamet-hybrid}};
  % matching arrows
  \draw[fd2 match] (q.east) -- node[fd2 lab,above=2pt] {LaMET 匹配} node[fd2 lab,below=1pt] {$C\big(\frac xy,\frac{\mu}{|y|P^z}\big)$} (f.west);
  \draw[fd2 match] (h.east) -- node[fd2 lab,above=2pt] {SDF 匹配} node[fd2 lab,below=1pt] {$\mathcal C(\alpha,-z^2\mu^2)$} (I.west);
  % formulas
  \node[fd2 lab] at (0,-3.55) {LaMET：$\tilde q(x,P^z)=\displaystyle\int_{-1}^{1}\frac{\dd y}{|y|}\,C\Big(\frac xy,\frac{\mu}{|y|P^z}\Big)f(y,\mu)+\order{\frac{\Lambda^2}{x^2P_z^2},\frac{\Lambda^2}{(1-x)^{2-\delta}P_z^2}}$};
  \node[fd2 lab] at (0,-4.65) {SDF：$\tilde h(\lambda,z^2)=\displaystyle\int_{-1}^{1}\!\dd\alpha\,\mathcal C(\alpha,-z^2\mu^2)\,I(\alpha\lambda,\mu)+\order{z^2\Lambda^2}$\quad（$|z|\lesssim0.2$--$0.3$ fm）};
  % status footer
  \node[draw=gray!60,fill=gray!5,rounded corners=2pt,inner sep=4pt,font=\small] at (0,-5.6)
    {微扰匹配：全阶论证\hspace{0.25em}{\color{RoyalBlue!80!black}\textbf{[A]}}；强子层次：依赖编号假设\hspace{0.25em}{\color{Plum!80!black}\textbf{[C]}}};
\end{tikzpicture}
\caption{LaMET 与短距离因子化（SDF）的结构。格点可计算的等时关联函数（左）经微扰匹配与光锥 PDF（右）相联系：准 PDF 与准 ITD 是同一矩阵元的两种 Fourier 表示，LaMET 在固定 $P^z$ 下作动量空间匹配，匹配核的宗量是部分子动量 $|y|P^z$（而非强子动量 $P^z$）；
SDF 在固定 $z^2$ 的短距离处作坐标空间匹配（VF-LAM-6）。误差为幂次修正：动量空间中在两端点分别按 $\Lambda^2/(x^2P_z^2)$ 与 $\Lambda^2/((1-x)^{2-\delta}P_z^2)$ 增大（$\delta$ 由扭度 4 分布的端点行为决定）。
部分子态的匹配定理是全阶论证（\cref{thm:lattice-lamet}），强子层次的陈述依赖于编号假设 \cref{asm:lamet-uniform,asm:lamet-twist4,asm:lamet-large-z}（\cref{thm:lamet-hadronic}）。}
\label{fig:lattice-matching}
\end{figure}
\clearpage
\begin{figure}[htbp]
\centering
\begin{tikzpicture}[fd,
  fd2 node/.style={rectangle,rounded corners=1.5pt,draw=black!65,line width=0.5pt,font=\tiny\ttfamily,
      inner sep=1.5pt,minimum height=0.62cm,minimum width=2.05cm,align=center},
  fd2 R/.style={fd2 node,fill=OliveGreen!15,draw=OliveGreen!65!black},
  fd2 A/.style={fd2 node,fill=RoyalBlue!15,draw=RoyalBlue!65!black},
  fd2 math/.style={fd2 node,fill=white},
  fd2 cond/.style={fd2 node,fill=OliveGreen!15,draw=OliveGreen!65!black,dash pattern=on 2.2pt off 1.4pt},
  fd2 asm/.style={fd2 node,draw=orange!65!black,fill=orange!8,dash pattern=on 2.2pt off 1.4pt},
  fd2 asmC/.style={fd2 node,draw=red!60!black,fill=red!9,dash pattern=on 2.2pt off 1.4pt,line width=0.7pt},
  fd2 grade/.style={font=\tiny\sffamily\bfseries,text=#1!75!black,inner sep=0.6pt,anchor=north east},
  fd2 use/.style={draw=gray!60,line width=0.5pt,-{Stealth[length=3.4pt,width=2.8pt]}},
  fd2 prem/.style={draw=orange!65!black,line width=0.6pt,dash pattern=on 2.5pt off 1.5pt,-{Stealth[length=3.6pt,width=3pt]}},
  fd2 neck/.style={draw=red!60!black,line width=1.3pt,-{Stealth[length=4.5pt,width=4pt]}}]
  % ---------------- nodes (layers left to right) ----------------
  \node[fd2 asm]  (aux) at (0,2.0)     {lamet-aux-renorm};
  \node[fd2 R]    (sup) at (0,-1.0)    {lamet-support};
  \node[fd2 A]    (ren) at (2.3,2.0)   {lattice-\\renormalizability};
  \node[fd2 R]    (ent) at (2.3,-1.0)  {lamet-entire};
  \node[fd2 A]    (cfp) at (4.6,3.2)   {lamet-cfp};
  \node[fd2 asm]  (ope) at (4.6,2.0)   {lamet-ope};
  \node[fd2 asmC] (tw4) at (4.6,-0.6)  {lamet-twist4};
  \node[fd2 math] (fou) at (4.6,-1.7)  {lamet-fourier};
  \node[fd2 A]    (lcs) at (4.6,-2.8)  {lamet-lcs};
  \node[fd2 cond] (coo) at (6.9,1.0)   {lamet-coordinate};
  \node[fd2 cond] (pmo) at (9.2,2.0)   {lamet-parton-\\momentum};
  \node[fd2 cond] (tmc) at (9.2,0.95)  {lamet-tmc};
  \node[fd2 asmC] (lgz) at (9.2,-0.6)  {lamet-large-z};
  \node[fd2 A]    (lat) at (11.5,2.6)  {lattice-lamet};
  \node[fd2 asmC] (uni) at (11.5,0.75) {lamet-uniform};
  \node[fd2 cond] (had) at (13.8,1.0)  {lamet-hadronic};
  % grade letters: tags of the assumptions, content grades of the R_cond theorems
  \foreach \n/\g/\c in {aux/A/RoyalBlue,ope/A/RoyalBlue,tw4/C/Plum,lgz/C/Plum,uni/C/Plum,
                        coo/A/RoyalBlue,pmo/A/RoyalBlue,tmc/A/RoyalBlue,had/C/Plum}
    \node[fd2 grade=\c] at ($(\n.north east)+(-0.03,-0.02)$) {\g};
  % ---------------- edges (transitive reduction; premise -> dependent) ----------------
  \begin{scope}[on background layer]
    \draw[fd2 neck] (aux) -- (ren);
    \draw[fd2 use] (ren) -- (ope);
    \draw[fd2 use] (ren.north east) -- (cfp.west);
    \draw[fd2 use] (sup) -- (ent);
    \draw[fd2 use] (ent.north east) .. controls (3.4,0.15) and (5.1,0.45) .. (coo.south west);
    \draw[fd2 use] (ent.east) -- (tw4.west);
    \draw[fd2 use] (ent.east) -- (fou.west);
    \draw[fd2 use] (ent.south east) -- (lcs.west);
    \draw[fd2 prem] (ope.south east) -- (coo.north west);
    \draw[fd2 prem] (ope) -- (pmo);
    \draw[fd2 prem] (ope.east) .. controls (6.6,1.85) and (7.9,1.62) .. (tmc.north west);
    \draw[fd2 use] (coo.north east) -- (pmo.south west);
    \draw[fd2 use] (coo) -- (tmc);
    \draw[fd2 use] (coo.south east) -- (lgz.north west);
    \draw[fd2 use] (tw4) -- (lgz);
    \draw[fd2 use] (pmo.east) -- (lat.south west);
    \draw[fd2 use] (cfp.east) -- (lat.north west);
    \draw[fd2 use] (tmc.east) -- (uni.west);
    \draw[fd2 use] (lat.south east) -- (had.north);
    \draw[fd2 prem] (uni) -- (had.west);
    \draw[fd2 prem] (lgz.east) -- (had.south west);
    \draw[fd2 prem] (tw4.south east) .. controls (7.6,-1.45) and (12.6,-1.35) .. (had.south);
    \draw[fd2 use] (fou.east) .. controls (8.6,-2.05) and (13.9,-1.95) .. (had.south east);
  \end{scope}
  % ---------------- legend ----------------
  \begin{scope}[shift={(0.2,-3.85)}]
    \node[fd2 R,minimum width=0.5cm,minimum height=0.3cm] at (-0.9,0) {};    \node[font=\scriptsize,anchor=west] at (-0.6,0) {[R]};
    \node[fd2 A,minimum width=0.5cm,minimum height=0.3cm] at (0.3,0) {};     \node[font=\scriptsize,anchor=west] at (0.6,0) {[A]};
    \node[fd2 math,minimum width=0.5cm,minimum height=0.3cm] at (1.5,0) {};  \node[font=\scriptsize,anchor=west] at (1.8,0) {纯数学};
    \node[fd2 cond,minimum width=0.5cm,minimum height=0.3cm] at (3.3,0) {};  \node[font=\scriptsize,anchor=west] at (3.6,0) {$R_{\mathrm{cond}}$（角标：内容等级）};
    \node[fd2 asm,minimum width=0.5cm,minimum height=0.3cm] at (7.6,0) {};   \node[font=\scriptsize,anchor=west] at (7.9,0) {[A] 级假设};
    \node[fd2 asmC,minimum width=0.5cm,minimum height=0.3cm] at (9.85,0) {}; \node[font=\scriptsize,anchor=west] at (10.15,0) {[C] 级假设（强子层次的前提）};
    \draw[fd2 use] (-1.2,-0.55) -- (-0.6,-0.55);   \node[font=\scriptsize,anchor=west] at (-0.5,-0.55) {\texttt{\string\uses} 依赖};
    \draw[fd2 prem] (2.1,-0.55) -- (2.7,-0.55);    \node[font=\scriptsize,anchor=west] at (2.8,-0.55) {假设作为 $R_{\mathrm{cond}}$ 定理的前提};
    \draw[fd2 neck] (7.6,-0.55) -- (8.2,-0.55);    \node[font=\scriptsize,anchor=west] at (8.3,-0.55) {[A] 陈述的瓶颈假设};
  \end{scope}
\end{tikzpicture}
\caption{LaMET 匹配定理的依赖结构（\cref{tab:lamet-dependency} 的图示）。节点为去掉类型前缀的规范标签，箭头由前提指向依赖者，只画传递约化后的边（可由其他路径推出的依赖从略，假设到条件定理的前提边一律保留）；
填充色为严格性标签（绿 [R]、蓝 [A]，白色为纯数学引理），绿色虚线框为 $R_{\mathrm{cond}}$ 条件定理，右上角字母为其内容等级；橙色与红色虚线框分别为 [A] 级与 [C] 级编号假设。
第~7、8~章的前提（\cref{thm:dis-ope,thm:pdf-moments-local,thm:pdf-dglap}）未画出。
部分子态的链 \cref{thm:lamet-coordinate}$\,\to\,$\cref{thm:lamet-parton-momentum}$\,\to\,$\cref{thm:lattice-lamet} 只依赖 [A] 级输入；强子层次的陈述 \cref{thm:lamet-hadronic} 另需三个 [C] 级假设（红色节点），其内容等级因而为 [C]。}
\label{fig:lamet-dependency}
\end{figure}
\clearpage
\begin{figure}[htbp]
\centering
\begin{tikzpicture}[fd,
  fd2 ax/.style={lcaxis,gray!60!black},
  fd2 ratio/.style={draw=red!65!black,line width=1.1pt},
  fd2 hyb/.style={draw=blue!60!black,line width=1.6pt},
  fd2 note/.style={font=\footnotesize},
  fd2 lab/.style={font=\small}]
  % x(z) = 10 z  (|z| in fm);  y = 2.4 (f - 0.4),  f(z) = 1 + 0.25 ln(z/0.2) + 0.5 max(0,z-0.45)^2  (schematic)
  % ======================= (a) the two |z| intervals =======================
  \begin{scope}[yshift=5.1cm]
    \fill[blue!12] (0,0) rectangle (2.5,0.62);
    \fill[orange!14] (2.5,0) rectangle (12,0.62);
    \draw[gray!60!black,line width=0.5pt] (0,0) rectangle (12,0.62);
    \node[fd2 lab] at (1.25,0.31) {比值方案};
    \node[fd2 lab] at (7.25,0.31) {Wilson 线质量减除：$\tilde h\,\ee^{-\delta m|z|}Z_{\rm hyb}$};
    \node[fd2 note,anchor=south] at (1.25,0.7) {$\tilde h(\lambda,z^2)/\tilde h(0,z^2)$};
    \node[fd2 note,anchor=south] at (7.25,0.7) {$Z_{\rm hyb}=\ee^{\delta m z_S}/\tilde h(0,z_S^2)$ 由 $z_S$ 处的连续性确定};
    \draw[fd2 ax] (0,-0.02) -- (12.7,-0.02) node[right,font=\small] {$|z|$};
    \foreach \t/\l in {0/0,12/1.2\,\mathrm{fm}} {\draw[gray!60!black,line width=0.45pt] (\t,-0.1) -- (\t,0.06); \node[fd2 note,below=2pt] at (\t,-0.06) {$\l$};}
    \draw[gray!60!black,line width=0.45pt] (2.5,-0.1) -- (2.5,0.06);
    \draw[gray!55!black,line width=0.45pt] (2.0,-0.3) -- (3.0,-0.3);
    \draw[gray!55!black,line width=0.45pt] (2.0,-0.36) -- (2.0,-0.24) (3.0,-0.36) -- (3.0,-0.24);
    \node[fd2 note,anchor=west] at (3.08,-0.31) {$z_S\approx0.2$--$0.3$ fm};
    \node[fd2 lab] at (-1.25,0.31) {(a)};
  \end{scope}
  % ======================= (b) conversion factor to MS-bar (schematic) =======================
  \fill[gray!14] (4.5,0) rectangle (12,4.0);
  \node[fd2 note,align=center,text=gray!35!black] at (9.6,0.38) {非微扰区：$\alpha_s(1/|z|)$ 不再小};
  \draw[gray!65!black,line width=0.5pt,dash pattern=on 3pt off 2pt] (2.5,0) -- (2.5,5.1);
  \draw[fd2 ratio] plot[domain=0.03:0.45,samples=60,smooth] ({10*\x},{2.4*(0.6+0.25*ln(\x/0.2))});
  \draw[fd2 ratio,dash pattern=on 4pt off 2pt] plot[domain=0.45:1.2,samples=40,smooth] ({10*\x},{2.4*(0.6+0.25*ln(\x/0.2)+0.5*(\x-0.45)^2)});
  \draw[fd2 hyb] plot[domain=0.03:0.25,samples=40,smooth] ({10*\x},{2.4*(0.6+0.25*ln(\x/0.2))}) -- (12,{2.4*(0.6+0.25*ln(1.25))});
  \fill[black] (2.5,{2.4*(0.6+0.25*ln(1.25))}) circle[radius=1.8pt];
  \draw[fd2 ax] (0,0) -- (12.7,0) node[right,font=\small] {$|z|/\mathrm{fm}$};
  \draw[fd2 ax] (0,0) -- (0,4.25);
  \node[fd2 note,rotate=90,anchor=south] at (-0.12,2.1) {到 $\MSbar$ 的换算因子（示意）};
  \foreach \t/\l in {0/0,2/0.2,4/0.4,6/0.6,8/0.8,10/1.0,12/1.2} {\draw[gray!60!black,line width=0.45pt] (\t,-0.06) -- (\t,0.06); \node[font=\scriptsize,below=2pt] at (\t,-0.04) {$\l$};}
  \node[fd2 note,text=red!55!black,anchor=south east] at (11.6,3.25) {比值方案、RI/MOM：依赖 $\ln(z^2\mu^2)$};
  \node[fd2 note,text=blue!50!black,anchor=north west] at (4.65,1.42) {混合方案：$|z|>z_S$ 时保持在 $z_S$ 处的值（处处微扰）};
  \node[fd2 lab] at (-1.25,2.1) {(b)};
\end{tikzpicture}
\caption{混合重整化方案（示意）。(a)~$|z|\le z_S$ 用比值方案，紫外发散（包括 Wilson 线的线性发散）在比值中相消；$|z|>z_S$ 只减去 Wilson 线的线性发散与对数因子（$\delta m$ 的符号约定同 \cref{prop:lattice-renormalizability}：$O^R=Z_j^{-2}\ee^{-\delta m|z|}O^B$）。
(b)~到 $\MSbar$ 的换算因子：混合方案中它等于 $z_S$ 处的值，在所有 $|z|$ 上都是微扰的，因此动量空间匹配成立；比值方案与 RI/MOM 方案的换算因子依赖于 $\ln(z^2\mu^2)$，在大 $|z|$ 处含 $\alpha_s(1/|z|)$（灰色区），动量空间匹配没有依据（\cref{rem:lamet-ratio-large-z}）。
单圈混合核满足 $\int\dd\xi\,C_{\rm hyb}=1$（\cref{prop:lamet-hybrid-kernel}，VF-LAM-10）。}
\label{fig:lamet-hybrid}
\end{figure}
\clearpage
\begin{figure}[htbp]
\centering
\begin{tikzpicture}
\begin{semilogyaxis}[name=main,width=0.67\textwidth,height=6.6cm,xmin=0,xmax=20,ymin=4e-3,ymax=1.6,
  log origin=infty,axis lines=left,axis line style={gray!60!black},
  xlabel={Ioffe 时间 $\lambda$},ylabel={$|\tilde h(\lambda,-\lambda^2/P_z^2)|/A$},
  xtick={0,5,10,15,20},
  tick label style={font=\scriptsize},label style={font=\small}]
  % data window lambda <= lambda_L = z_L P^z = 10.14
  \fill[blue!9] (axis cs:0,4e-3) rectangle (axis cs:10.14,1.6);
  \draw[blue!55!black,line width=0.6pt,dash pattern=on 3pt off 2pt] (axis cs:10.14,4e-3) -- (axis cs:10.14,1.6);
  % truncated tail
  \addplot[draw=none,pattern=north east lines,pattern color=red!60!black,domain=10.14:20,samples=60] {exp(-0.25*x)} \closedcycle;
  % the bound |h| <= A exp(-m lambda/P^z), m = 0.5 GeV, P^z = 2 GeV
  \addplot[black,line width=1.1pt,domain=0:20,samples=120] {exp(-0.25*x)};
  \node[font=\footnotesize,anchor=south west] at (axis cs:2.0,0.64) {$A\,\ee^{-m\lambda/P^z}$};
  \node[font=\footnotesize,text=blue!45!black,align=center] at (axis cs:5.0,0.02) {数据窗 $|\lambda|\le\lambda_L$\\$\lambda_L=z_LP^z\approx10.1$\\（$z_L=1$ fm）};
  \node[font=\footnotesize,text=red!55!black,align=center] at (axis cs:15.2,0.42) {截断的尾部（斜线）：\\误差 $\le\dfrac{AP^z}{\pi m}\,\ee^{-mz_L}$\\$\ee^{-mz_L}=0.079$};
  \draw[red!55!black,line width=0.5pt,-{Stealth[length=4pt,width=3.4pt]}] (axis cs:14.6,0.115) -- (axis cs:13.4,0.016);
\end{semilogyaxis}
% ---------------- inset: analyticity strip in the complex x plane ----------------
\begin{scope}[shift={($(main.outer east)+(2.05,0.35)$)},fd]
  \fill[white] (-1.85,-0.95) rectangle (2.1,1.05);
  \draw[gray!55,line width=0.4pt,rounded corners=2pt] (-1.85,-0.95) rectangle (2.1,1.05);
  \fill[orange!20] (-1.6,-0.32) rectangle (1.6,0.32);
  \draw[orange!60!black,line width=0.5pt,dash pattern=on 2.5pt off 1.5pt] (-1.6,0.32) -- (1.6,0.32) (-1.6,-0.32) -- (1.6,-0.32);
  \draw[lcaxis,gray!60!black] (-1.7,0) -- (1.75,0) node[right,font=\scriptsize] {$\Real x$};
  \draw[lcaxis,gray!60!black] (0,-0.5) -- (0,0.88) node[right,font=\scriptsize] {$\Imag x$};
  \foreach \t/\l in {-1/-1,1/1} {\draw[gray!60!black,line width=0.4pt] (\t,-0.05) -- (\t,0.05); \node[font=\scriptsize,fill=orange!20,inner sep=0.5pt] at (\t,-0.17) {$\l$};}
  \node[font=\scriptsize,anchor=west] at (0.05,0.5) {$\Imag x=m/P^z$};
  \node[font=\scriptsize,align=center] at (0.12,-0.7) {$\tilde q(x,P^z)$ 在带内全纯};
\end{scope}
\end{tikzpicture}
\caption{Fourier 步骤的截断误差与解析带（示意，数值取自 VF-LAM-17）。若准 ITD 沿射线满足 $|\tilde h(\lambda,-\lambda^2/P_z^2)|\le A\,\ee^{-m|\lambda|/P^z}$（指数衰减假设），则准 PDF 在带形 $|\Imag x|<m/P^z$ 中全纯（内插图，\cref{lem:lamet-strip}），在有限 $P^z$ 下没有端点奇异性；
只用数据窗 $|\lambda|\le\lambda_L=z_LP^z$（蓝色；图中 $m=0.5$ GeV，$P^z=2$ GeV）作 Fourier 变换时，截去的尾部（斜线）引起的误差不超过 $\frac{AP^z}{\pi m}\ee^{-mz_L}$（\cref{prop:lamet-truncation}）。对 $m=0.5$ GeV、$z_L=1$ fm，$\ee^{-mz_L}=0.079$：对现有的 $z_L$ 这一压低并不是参数上小的。QCD 中这种衰减本身是猜想（\cref{conj:lamet-decay-qcd}）。}
\label{fig:lamet-fourier}
\end{figure}
\clearpage
\begin{figure}[htbp]
\centering
\begin{tikzpicture}
\begin{semilogyaxis}[width=0.8\textwidth,height=6.4cm,xmin=0,xmax=1,ymin=5e-3,ymax=2,
  axis lines=left,axis line style={gray!60!black},clip=false,
  xlabel={$x$},ylabel={相对幂次修正（示意）},
  xtick={0,0.2,0.5,0.8,1},
  xticklabels={$0$,$x_0$,$0.5$,$1-x_0$,$1$},
  extra x ticks={0.1,0.9},extra x tick labels={$\frac{\Lambda}{P^z}$,$1-\frac{\Lambda}{P^z}$},
  extra x tick style={tick label style={yshift=-9pt}},
  ytick={0.01,0.1,1},yticklabels={$10^{-2}$,$10^{-1}$,$1$},
  tick label style={font=\scriptsize},label style={font=\small},
  legend style={font=\scriptsize,draw=gray!40,fill=white,at={(0.5,0.97)},anchor=north},
  legend cell align=left]
  % excluded end-point strips (Lambda/P^z = 0.1 in this sketch)
  \fill[gray!22] (axis cs:0,5e-3) rectangle (axis cs:0.1,2);
  \fill[gray!22] (axis cs:0.9,5e-3) rectangle (axis cs:1,2);
  \draw[gray!60!black,line width=0.4pt,dash pattern=on 2pt off 2pt] (axis cs:0,1) -- (axis cs:1,1);
  \addplot[blue!60!black,line width=1.1pt,domain=0.0708:0.9,samples=120] {0.01/x^2};
  \addlegendentry{$\Lambda^2/(x^2P_z^2)$}
  \addplot[red!65!black,line width=1.1pt,dashed,domain=0.1:0.9293,samples=120] {0.01/(1-x)^2};
  \addlegendentry{$\Lambda^2/((1-x)^2P_z^2)$（$\delta=0$）}
  \addplot[red!65!black,line width=1.3pt,dotted,domain=0.1:0.995,samples=160] {0.01/(1-x)};
  \addlegendentry{$\Lambda^2/((1-x)P_z^2)$（$\delta=1$）}
  \node[font=\scriptsize,align=center,rotate=90] at (axis cs:0.05,0.06) {$xP^z\lesssim\Lambda$：不适用};
  \node[font=\scriptsize,align=center,rotate=90] at (axis cs:0.95,0.022) {$(1-x)P^z\lesssim\Lambda$};
  % fixed-order domain and its extension by threshold resummation
  \draw[gray!50!black,line width=0.5pt] (axis cs:0.2,2.6) -- (axis cs:0.2,3.0) -- (axis cs:0.8,3.0) -- (axis cs:0.8,2.6);
  \node[font=\scriptsize,anchor=south] at (axis cs:0.5,3.0) {固定阶匹配一致成立：$[x_0,1-x_0]$，$x_0P^z\gg\Lambda$};
  \draw[red!55!black,line width=0.7pt,-{Stealth[length=4.5pt,width=3.8pt]}] (axis cs:0.8,2.75) -- (axis cs:0.9,2.75);
  \node[font=\scriptsize,text=red!50!black,anchor=south west,align=left] at (axis cs:0.82,3.05) {阈值重求和：大 $x$ 边界推到\\$(1-x)P^z\gg\Lambda$};
\end{semilogyaxis}
\end{tikzpicture}
\caption{LaMET 预言的适用范围（示意，图中取 $\Lambda/P^z=0.1$）。端点区要求硬标度 $2xP^z$ 与旁观者动量 $(1-x)P^z$ 都远大于 $\Lambda$（灰色带不适用）。相对幂次修正在小 $x$ 处按 $\Lambda^2/(x^2P_z^2)$ 增大；在大 $x$ 处，两种计数规则 $\Lambda^2/((1-x)^{2-\delta}P_z^2)$（$\delta=0$ 虚线，$\delta=1$ 点线）对应扭度 4 分布 $\rho_4/q\sim(1-x)^\delta$ 的不同端点行为（\cref{cor:lamet-endpoint}，VF-LAM-12），$\delta$ 在 QCD 中未知。
固定阶匹配只在 $[x_0,1-x_0]$ 上一致成立；阈值重求和（\cref{thm:yzliu-threshold}）把大 $x$ 边界推到 $(1-x)P^z\gg\Lambda$（\cref{rem:lamet-domain}）。}
\label{fig:lamet-validity}
\end{figure}
\clearpage
\begin{figure}[htbp]
\centering
\begin{tikzpicture}[fd,
  fd2 ax/.style={lcaxis,gray!60!black},
  fd2 lc/.style={draw=gray!55,line width=0.5pt,dash pattern=on 1.2pt off 1.6pt},
  fd2 warrow/.style={postaction={decorate,decoration={markings,
      mark=at position #1 with {\arrow[xshift=3.2pt]{Stealth[length=6.4pt,width=6pt]}}}}},
  fd2 heavy/.style={draw,solid,line width=1.4pt,fd arrow=#1},
  fd2 heavy/.default=0.5,
  fd2 sheet/.style={fill=gray!13,draw=none},
  fd2 lab/.style={font=\small},
  fd2 note/.style={font=\footnotesize}]
  % geometry: beta = 0.55, t = t' = 1.85; perspective b_perp vector B = (1.7,0.55)
  \foreach \panel/\xs in {a/0, b/7.6} {
  \begin{scope}[xshift=\xs cm]
    \coordinate (O) at (0,0);          \coordinate (OB) at (1.7,0.55);
    \coordinate (P) at (-1.0175,-1.85); \coordinate (PB) at (0.6825,-1.3);
    \coordinate (F) at (-1.0175,1.85);  \coordinate (FB) at (0.6825,2.4);
    % axes and light cone (b_perp = 0)
    \draw[fd2 lc] (-1.75,-1.75) -- (1.75,1.75);  \draw[fd2 lc] (-1.75,1.75) -- (1.75,-1.75);
    \draw[fd2 ax] (-1.75,0) -- (2.75,0) node[below left=1pt and -3pt,font=\footnotesize] {$z$};
    \draw[fd2 ax] (0,-2.25) -- (0,2.95) node[left,font=\footnotesize] {$t$};
    \draw[gray!60!black,line width=0.5pt,dash pattern=on 3pt off 2pt,-{Stealth[length=4pt,width=3.4pt]}] (O) -- (2.45,0.79) node[right,font=\footnotesize] {$\vec b_\perp$};
    \if a\panel
      % two staple-shaped Wilson lines: past (v) and future (v'), closed at -t and t'
      \draw[wilson plain] (OB) -- (PB) -- (P) -- (O) -- (F) -- (FB) -- cycle;
      \path[fd2 warrow=0.55] (P) -- (O);    \path[fd2 warrow=0.6] (O) -- (F);
      \path[fd2 warrow=0.55] (FB) -- (OB);  \path[fd2 warrow=0.6] (OB) -- (PB);
      \path[fd2 warrow=0.55] (F) -- (FB);   \path[fd2 warrow=0.55] (PB) -- (P);
      \node[opvertex] at (O) {};  \node[opvertex] at (OB) {};
      \node[fd2 lab,anchor=east] at (-0.6,-1.0) {$v$};
      \node[fd2 lab,anchor=east] at (-0.6,1.0) {$v'$};
      \node[fd2 note,anchor=north east] at ($(P)+(0.02,-0.04)$) {$-t$};
      \node[fd2 note,anchor=south east] at ($(F)+(0.02,0.04)$) {$t'$};
      \node[fd2 note,anchor=north west] at (0.05,-0.02) {$0$};
      \node[fd2 lab,align=center] at (0.45,-3.1) {$S=\displaystyle\lim_{t,t'\to\infty}\frac{W(t,t',b_\perp,Y,Y')}{\sqrt{Z(2t,b_\perp,Y)\,Z(2t',b_\perp,Y')}}$};
      \node[fd2 note,align=left,anchor=west] at (1.95,1.75) {时序乘积 $\mathcal T$};
      \node[fd2 lab] at (0.45,-4.2) {(a)\ 订书钉形 Wilson 回路};
    \else
      % HQET picture: static colour sources Q (at 0) and Qbar (at b_perp)
      \fill[fd2 sheet] (P) -- (O) -- (F) -- (FB) -- (OB) -- (PB) -- cycle;
      \draw[fd2 heavy=0.55] (P) -- (O);    \draw[fd2 heavy=0.55] (O) -- (F);
      \draw[fd2 heavy=0.5] (FB) -- (OB);   \draw[fd2 heavy=0.55] (OB) -- (PB);
      \draw[gray!60!black,line width=0.6pt,dash pattern=on 2.5pt off 1.5pt] (O) -- (OB);
      \node[opvertex] at (O) {};  \node[opvertex] at (OB) {};
      \node[fd2 lab,anchor=east] at (-0.6,-1.0) {$v$};
      \node[fd2 lab,anchor=east] at (-0.6,1.0) {$v'$};
      \node[fd2 note,anchor=east] at (-1.05,-1.9) {$Q$};
      \node[fd2 note,anchor=west] at (0.72,-1.38) {$\bar Q$};
      \node[fd2 note,anchor=south,inner sep=1pt] at (0.62,0.44) {$J$（$t=0$）};
      \node[fd2 lab,anchor=north] at (-0.15,-2.15) {$|Q\bar Q,\vec b_\perp\rangle_v$};
      \node[fd2 lab,anchor=west] at (0.8,2.6) {${}_{v'}\langle Q\bar Q,\vec b_\perp|$};
      \node[draw=gray!55,fill=gray!4,rounded corners=2pt,font=\footnotesize,align=left,anchor=north west,inner sep=3pt] at (1.75,-0.55)
        {欧氏表示：\\$\ii\tilde v\cdot D_E=D_\tau-\ii\beta D_z$\\（$\tau$：欧氏时间）};
      \node[fd2 lab] at (0.45,-4.2) {(b)\ HQET 形状因子};
    \fi
  \end{scope}}
  \node[fd2 lab,align=center] at (8.05,-3.1) {$S={}_{v'}\langle Q\bar Q,\vec b_\perp|\,J\,|Q\bar Q,\vec b_\perp\rangle_v$};
\end{tikzpicture}
\caption{离光锥 TMD 软函数作为快速色源对的形状因子（$v=\gamma(1,\beta,\mathbf 0_\perp)$、$v'=\gamma'(1,-\beta',\mathbf 0_\perp)$，$\beta=\tanh Y$、$\beta'=\tanh Y'$，坐标 $(t,z,\vec b_\perp)$，$\vec b_\perp$ 方向按透视画出）。
(a)~两个订书钉形 Wilson 回路（$t<0$ 沿 $v$，$t>0$ 沿 $v'$，在 $-t$ 与 $t'$ 处由横向连接闭合）的时序真空期望，除以 $\sqrt{ZZ}$（$Z$ 为矩形回路）后取 $t,t'\to\infty$；箭头为路径方向：$0$ 处沿世界线指向未来（色源 $Q$），$\vec b_\perp$ 处指向过去（$\bar Q$）。
(b)~等价的 HQET 图像：间隔 $b_\perp$ 的色单态静态 $Q\bar Q$ 对以速度 $v$ 运动，在 $t=0$ 被等时的流 $J$ 踢到 $v'$，软函数是 ${}_{v'}\langle Q\bar Q,\vec b_\perp|J|Q\bar Q,\vec b_\perp\rangle_v$，因而有欧氏表示（\cref{thm:yzliu-soft-hqet}，全阶论证 [A]；非微扰的形状因子形式另需\cref{asm:yzliu-heavy-pair}）。
大快度下 $\ln S=(Y+Y')\tilde K+\mathcal D+\order{\ee^{-(Y+Y')}}$，本征软函数 $S_I=\ee^{-\mathcal D}$（VF-YZL-2）。}
\label{fig:yzliu-soft-hqet}
\end{figure}
\clearpage
\begin{figure}[htbp]
\centering
\begin{tikzpicture}[fd,
  fd2 warrow/.style={postaction={decorate,decoration={markings,
      mark=at position #1 with {\arrow[xshift=3.2pt]{Stealth[length=6.4pt,width=6pt]}}}}},
  fd2 dipole/.style={draw=black!80,line width=2.4pt,dash pattern=on 6pt off 2.5pt,line cap=butt},
  fd2 web/.style={gluon,draw=orange!70!black},
  fd2 vtx/.style={gluon,draw=blue!60!black},
  fd2 big/.style={font=\large},
  fd2 lab/.style={font=\small},
  fd2 note/.style={font=\footnotesize},
  fd2 staple/.pic={
    \draw[wilson plain] (0,0) -- (2.4,0) -- (2.4,0.8) -- (0.45,0.8);
    \path[fd2 warrow=0.55] (0,0) -- (2.4,0);  \path[fd2 warrow=0.6] (2.4,0.8) -- (0.45,0.8);
    \draw[quark=0.55] (-0.5,-0.42) -- (0,0);  \draw[quark=0.6] (0.45,0.8) -- (-0.05,1.22);
    \node[opvertex] at (0,0) {};  \node[opvertex] at (0.45,0.8) {};}]
  % ======================= (a) single-line notation =======================
  \pic at (0,0) {fd2 staple};
  \node[fd2 note,anchor=north] at (1.2,-0.08) {$L$};
  \node[fd2 note,anchor=west] at (2.45,0.4) {$b_\perp$};
  \node[fd2 big] at (3.65,0.4) {$\equiv$};
  \draw[quark=0.55] (4.2,-0.02) -- (4.7,0.4);  \draw[quark=0.6] (4.7,0.4) -- (4.2,0.82);
  \draw[fd2 dipole] (4.7,0.4) -- (7.4,0.4);
  \draw[black!80,line width=1.0pt] (7.4,0.2) -- (7.4,0.6);
  \node[opvertex,minimum size=8pt] at (4.7,0.4) {};
  \node[fd2 note,anchor=south] at (6.2,0.58) {偶极线：量纲 $\frac32$、自旋 $0$};
  \node[fd2 note,anchor=north] at (6.5,0.2) {（一对订书钉腿合成一条线）};
  \node[fd2 lab] at (3.65,-1.0) {(a)\ 单线记号：订书钉的两条腿在 $\hat z$ 方向的演化合成一个色单态偶极子};
  % ======================= (b) dipole 2PI decomposition =======================
  \begin{scope}[yshift=-3.3cm]
    % left: full correlator with a self-interaction web and an operator--link vertex
    \draw[fd2 web] (1.6,0) -- (1.6,0.8);
    \draw[fd2 vtx] (-0.3,-0.25) to[bend right=40] (0.75,0);
    \pic at (0,0) {fd2 staple};
    \node[vertex] at (1.6,0) {}; \node[vertex] at (1.6,0.8) {}; \node[vertex] at (0.75,0) {};
    \node[fd2 big] at (3.2,0.4) {$=$};
    % exponentiated webs
    \node[fd2 big] at (3.95,0.4) {$\exp\Big\{$};
    \begin{scope}[xshift=4.6cm]
      \draw[fd2 web] (1.2,0) -- (1.2,0.8);
      \draw[wilson plain] (0,0) -- (2.4,0) -- (2.4,0.8) -- (0,0.8);
      \node[vertex] at (1.2,0) {}; \node[vertex] at (1.2,0.8) {};
    \end{scope}
    \node[fd2 lab,anchor=west] at (7.1,0.4) {$+\,\cdots\Big\}$};
    \node[fd2 big] at (8.75,0.4) {$\times$};
    % operator--link 2PI vertex part
    \begin{scope}[xshift=9.75cm]
      \draw[fd2 vtx] (-0.3,-0.25) to[bend right=40] (0.75,0);
      \pic at (0,0) {fd2 staple};
      \node[vertex] at (0.75,0) {};
    \end{scope}
    \node[fd2 note,text=orange!60!black] at (5.8,-0.42) {自相互作用 web：$\ee^{\Phi(\mathcal W_z)}=\langle0|\mathcal W_z|0\rangle$};
    \node[fd2 note,text=blue!50!black] at (10.95,-0.75) {算符—规范线 2PI 顶点};
    \node[fd2 lab,align=center] at (5.6,-1.65) {$\ee^{\Phi(\mathcal W_z)}\sim\ee^{-LE(b_\perp)}$，\quad $\sqrt{Z_E(2L,b_\perp)}\sim\ee^{-LE(b_\perp)}$，\quad $E=2\delta m+V(b_\perp)$};
    \node[fd2 lab] at (5.6,-2.45) {(b)\ 偶极子 2PI 分解：除以 $\sqrt{Z_E}$ 后自相互作用 web 相消，$L\to\infty$ 的极限存在};
  \end{scope}
\end{tikzpicture}
\caption{单线记号与偶极子 2PI 分解（示意；双线为类空 Wilson 线，箭头为路径方向，$\otimes$ 为夸克场端点）。(a)~准 TMD 的订书钉形 Wilson 线的两条腿沿 $\hat z$ 平行演化，可看作一个色单态“偶极子”，在单线记号中画成一条量纲 $\frac32$、自旋 $0$ 的线（用于\cref{prop:yzliu-qtmd-powercount}的幂次计数）。
(b)~部分指数化引理（\cref{lem:yzliu-partial-exp}）：关联函数等于规范线真空期望值 $\ee^{\Phi(\mathcal W_z)}$（其 web 的指数，橙色胶子为一例）乘以算符—规范线两粒子不可约顶点之和（蓝色胶子为一例）；$\ee^{\Phi(\mathcal W_z)}$ 与 $\sqrt{Z_E}$ 都按 $\ee^{-LE(b_\perp)}$ 依赖于 $L$，在比值中相消（\cref{prop:yzliu-pinch-ultrasoft}）。}
\label{fig:yzliu-single-line}
\end{figure}
\clearpage
\begin{figure}[htbp]
\centering
\begin{tikzpicture}[fd,x=1cm,y=0.75cm,
  fd2 ax/.style={lcaxis,gray!60!black},
  fd2 tick/.style={draw=gray!60!black,line width=0.45pt},
  fd2 pt/.style={circle,fill=black,draw=none,inner sep=0pt,minimum size=5.4pt},
  fd2 lab/.style={font=\footnotesize,align=left,anchor=west,inner sep=1.5pt},
  fd2 val/.style={font=\footnotesize,anchor=east,inner sep=1.5pt},
  fd2 note/.style={font=\footnotesize}]
  \def\Yo{3.0}
  % ---- the strip |Im theta| <= pi and its boundary lines ----
  \fill[blue!7] (0,-3.1416) rectangle (11.4,3.1416);
  \draw[blue!50!black,line width=0.75pt] (0,3.1416) -- (11.4,3.1416);
  \draw[blue!50!black,line width=0.75pt] (0,-3.1416) -- (11.4,-3.1416);
  \node[fd2 note,text=blue!45!black,anchor=south east] at (11.4,3.1416) {$\Imag\theta=\pi$};
  \node[fd2 note,text=blue!45!black,anchor=north east] at (11.4,-3.1416) {$\Imag\theta=-\pi$};
  \node[fd2 note,text=blue!45!black,align=center] at (9.35,1.62) {$S(\theta)$ 在 $|\Imag\theta|<\pi$ 中全纯，\\连续到边界};
  % ---- axes ----
  \draw[fd2 ax] (-0.7,0) -- (11.85,0) node[right] {$\Real\theta$};
  \draw[fd2 ax] (0,-4.05) -- (0,4.15) node[left] {$\Imag\theta$};
  \foreach \y/\t in {3.1416/{\pi},1.5708/{\frac{\pi}{2}},-1.5708/{-\frac{\pi}{2}},-3.1416/{-\pi}} {
    \draw[fd2 tick] (-0.07,\y) -- (0.07,\y);
    \node[font=\scriptsize,anchor=east,inner sep=1pt] at (-0.1,\y) {$\t$};}
  \node[font=\scriptsize,anchor=north east,inner sep=1.5pt] at (0,0) {$0$};
  % ---- the line Re theta = Y_0 ----
  \draw[gray!60,line width=0.5pt,dash pattern=on 3pt off 2pt] (\Yo,-3.85) -- (\Yo,3.85);
  \draw[fd2 tick] (\Yo,-0.07) -- (\Yo,0.07);
  \node[font=\scriptsize,anchor=north,inner sep=1.5pt] at (\Yo,-3.85) {$\Real\theta=Y_0$};
  % ---- the six configurations (VF-YZL-4) ----
  \foreach \y in {3.1416,1.5708,0,-1.5708,-3.1416} \node[fd2 pt] at (\Yo,\y) {};
  \node[fd2 val] at (\Yo-0.12,3.1416+0.32) {$Y_0+\ii\pi$};
  \node[fd2 val] at (\Yo-0.12,1.5708+0.32) {$Y_0+\frac{\ii\pi}{2}$};
  \node[fd2 val,fill=blue!7] at (\Yo-0.12,0.36) {$Y_0$};
  \node[fd2 val] at (\Yo-0.12,-1.5708+0.32) {$Y_0-\frac{\ii\pi}{2}$};
  \node[fd2 val] at (\Yo-0.12,-3.1416+0.32) {$Y_0-\ii\pi$};
  \node[fd2 lab] at (\Yo+0.15,3.1416+0.36) {类空/类空，过去/未来};
  \node[fd2 lab] at (\Yo+0.15,1.5708+0.36) {类时/类空，过去/未来};
  \node[fd2 lab] at (\Yo+0.15,0.62) {类空/类空，过去/过去（Collins 的 DY 软因子 $S_{\mathrm{DY}}$）};
  \node[fd2 lab] at (\Yo+0.15,-0.42) {类时/类时，过去/未来（形状因子 $S_{\mathrm{FF}}$）};
  \node[fd2 lab] at (\Yo+0.15,-1.5708+0.36) {类时/类空，过去/过去};
  \node[fd2 lab] at (\Yo+0.15,-3.1416+0.36) {类时/类时，过去/过去};
  % ---- light-cone limit ----
  \draw[red!60!black,line width=0.8pt,-{Stealth[length=5.5pt,width=4.6pt]}] (7.3,-2.05) -- (10.9,-2.05);
  \node[fd2 note,text=red!55!black,align=center,anchor=south] at (9.1,-2.0) {$Y\to\infty$：光锥极限};
  \node[fd2 note,text=red!55!black,align=center,anchor=north] at (9.1,-2.12) {（假设对 $\Imag\theta$ 一致）};
\end{tikzpicture}
\caption{离光锥软函数对双曲角的解析性。单一时序的减除软函数只通过双曲角 $\theta(v_1,v_2)$（$\cosh\theta=v_1\cdot v_2/\sqrt{v_1^2v_2^2}$，$\pm\ii0$ 处方见正文）依赖于两个方向，是带形 $-\pi\le\Imag\theta\le\pi$（阴影，$\Real\theta\ge0$）上同一函数 $S(\theta)$ 的取值；
在固定的 $\Real\theta=Y_0$ 上标出两个方向为类空/类时、过去/未来指向的六种组合所对应的 $\theta$（VF-YZL-4），其中 $\theta=Y_0$ 一点同时代表 Collins 的类空 DY 软因子与类时形状因子软函数，
由此在微扰论各阶 $S_{\mathrm{DY}}=S_{\mathrm{FF}}$（\cref{thm:yzliu-analyticity}，全阶论证 [A]；前提为维数正规化微扰论中的 Wightman 结构，\cref{asm:yzliu-wightman-pt}）。
光锥极限 $Y\to\infty$ 能否越过整个带形延拓，取决于该极限对 $\Imag\theta$ 的一致性；它只在 $\Imag\theta\in\{0,\pm\frac{\pi}{2},\pm\pi\}$ 处得到微扰验证（\cref{asm:yzliu-lightcone-uniform}，部分 [P]）。}
\label{fig:yzliu-theta-strip}
\end{figure}
\clearpage
\begin{figure}[htbp]
\centering
\begin{tikzpicture}[fd,scale=1.08,
  fd2 ax/.style={lcaxis,gray!60!black},
  fd2 lc/.style={draw=gray!55,line width=0.5pt,dash pattern=on 3pt off 2pt},
  fd2 wl/.style={wilson=#1},
  fd2 wlc/.style={draw=violet!70!black,solid,double,double distance=1.6pt,line width=0.55pt,
      postaction={decorate,solid,decoration={markings,mark=at position #1 with
      {\arrow[xshift=3.2pt,line width=0.6pt,violet!70!black]{Stealth[length=6.4pt,width=6pt]}}}}},
  fd2 vtx/.style={circle,fill=black,draw=none,inner sep=0pt,minimum size=4pt},
  fd2 vtxc/.style={circle,fill=violet!70!black,draw=none,inner sep=0pt,minimum size=4pt},
  fd2 shift/.style={draw=#1,line width=0.6pt,-{Stealth[length=4pt,width=3.4pt]}},
  fd2 lab/.style={font=\small},
  fd2 note/.style={font=\footnotesize}]
  % ===================== left: S, DY shape, exponential regulator =====================
  \begin{scope}
    \draw[fd2 ax] (-2.15,1.0) -- (-1.55,1.0) node[right=-1pt,fd2 note] {$z$};
    \draw[fd2 ax] (-2.15,1.0) -- (-2.15,1.6) node[above=-1pt,fd2 note] {$t$};
    \draw[fd2 lc] (-1.85,-1.85) -- (1.25,1.25);
    \draw[fd2 lc] (1.85,-1.85) -- (-1.25,1.25);
    % amplitude T[W^dag_nbar(0) W_n(0)]: past-pointing lines meeting at 0
    \draw[fd2 wl=0.55] (-1.5,-1.5) -- (0,0);
    \draw[fd2 wl=0.72] (0,0) -- (1.5,-1.5);
    \node[fd2 vtx] at (0,0) {};
    % conjugate amplitude Tbar[W^dag_n W_nbar] at b - i nu e_t (drawn displaced)
    \draw[fd2 wlc=0.55] (1.5,-2.65) -- (0,-1.15);
    \draw[fd2 wlc=0.55] (0,-1.15) -- (-1.5,-2.65);
    \node[fd2 vtxc] at (0,-1.15) {};
    \draw[fd2 shift=violet!60!black,dotted,line width=0.8pt] (0,-0.1) -- (0,-1.04);
    \node[font=\scriptsize,text=violet!55!black,anchor=west,inner sep=1pt] at (0.04,-0.9) {$-\ii\nu e_t$};
    % labels
    \node[fd2 note,anchor=west] at (0.1,0.02) {$0$};
    \node[fd2 note,anchor=north east] at (-1.42,-1.45) {$n$};
    \node[fd2 note,anchor=north west] at (1.42,-1.45) {$\bar n$};
    \node[fd2 note,anchor=south] at (0,0.14) {$T[W_{\bar n}^\dagger W_n](0)$};
    \node[fd2 note,anchor=north,align=center,text=violet!55!black] at (0,-2.75) {$\bar T[W_n^\dagger W_{\bar n}](\vec b_\perp-\ii\nu e_t)$\\[1pt]（虚时间位移）};
    \node[fd2 lab,align=center] at (0,-4.2) {$S$：DY 形软因子（指数正规子）};
  \end{scope}
  % ===================== right: S_t, single time ordering, space-like form factor =====================
  \begin{scope}[xshift=9.6cm]
    \draw[fd2 ax] (0.55,1.0) -- (1.15,1.0) node[right=-1pt,fd2 note] {$z$};
    \draw[fd2 ax] (0.55,1.0) -- (0.55,1.6) node[above=-1pt,fd2 note] {$t$};
    \draw[fd2 lc] (-1.85,-1.85) -- (1.0,1.0);
    \draw[fd2 lc] (1.0,-1.0) -- (-1.85,1.85);
    % pair at 0: incoming along n (past), outgoing along nbar (future)
    \draw[fd2 wl=0.55] (-1.5,-1.5) -- (0,0);
    \draw[fd2 wl=0.6] (0,0) -- (-1.5,1.5);
    \node[fd2 vtx] at (0,0) {};
    % pair at b - nu e_z (real space-like shift)
    \draw[fd2 wlc=0.5] (-2.7,1.5) -- (-1.2,0);
    \draw[fd2 wlc=0.55] (-1.2,0) -- (-2.7,-1.5);
    \node[fd2 vtxc] at (-1.2,0) {};
    \draw[fd2 shift=violet!60!black] (-0.1,0) -- (-1.1,0);
    \node[font=\scriptsize,text=violet!55!black,anchor=south,inner sep=1pt] at (-0.68,0.05) {$-\nu e_z$};
    % labels
    \node[fd2 note,anchor=west] at (0.1,0) {$0$};
    \node[fd2 note,anchor=north west] at (-1.42,-1.45) {$n$};
    \node[fd2 note,anchor=south west] at (-1.42,1.45) {$\bar n$};
    \node[fd2 note,anchor=north,align=center,text=violet!55!black] at (-2.0,-1.62) {$\vec b_\perp-\nu e_z$\\[1pt]（实的类空位移）};
    \node[fd2 note,anchor=south,align=center] at (-1.2,1.95) {$T[W_n^\dagger W_{\bar n+}^\dagger\,W_{\bar n+}W_n]$（单一时序）};
    \node[fd2 lab,align=center] at (-1.1,-4.2) {$S_t$：单一时序的类空形状因子};
  \end{scope}
  % ===================== the complex Lorentz transformation =====================
  \draw[line width=1.5pt,{Stealth[length=8pt,width=8pt]}-{Stealth[length=8pt,width=8pt]}] (2.25,0.05) -- (6.3,0.05);
  \node[fd2 lab,anchor=south] at (4.27,0.22) {$\Lambda(t,z)=(\ii z,\ii t)$};
  \node[fd2 note,anchor=north,align=left] at (4.27,-0.2) {$\Lambda(-\ii\lambda n)=\lambda n$\\[1pt]$\Lambda(-\ii\lambda\bar n)=-\lambda\bar n$\\[1pt]$\Lambda(-\nu e_z)=-\ii\nu e_t$};
  % ===================== invariant lengths =====================
  \node[fd2 lab] at (4.3,-5.0) {$x^2_{B,ij}=-\nu^2-b_\perp^2-\sqrt2\,\ii\nu\big(\lambda^L_i-\lambda^R_j\big)$：实部 $\le-b_\perp^2-\nu^2<0$，阈下、无割线};
\end{tikzpicture}
\caption{时序无关性的解析延拓图像（$(z,t)$ 平面示意；双线为类光 Wilson 线，箭头为路径方向；$n$、$\bar n$ 分别沿 $+z$、$-z$，即本书 $n_A$、$n_B$ 的方向；横向间隔 $\vec b_\perp$ 垂直于纸面）。
左：指数正规子下的 DY 形软因子 $S$——振幅中的 $T[W_{\bar n}^\dagger W_n](0)$ 与共轭振幅中位于 $\vec b_\perp-\ii\nu e_t$ 的 $\bar T[W_n^\dagger W_{\bar n}]$（紫色，虚时间位移画成竖直错开）；
右：单一时序的类空形状因子 $S_t$——位于 $0$ 与 $\vec b_\perp-\nu e_z$ 的一对类光色源，沿 $n$ 入射、沿 $\bar n$ 出射。
二者由复 Lorentz 变换 $\Lambda(t,z)=(\ii z,\ii t)$（快度为 $\ii\pi/2$ 的推动）相联系，所有不变间隔的实部都为负（阈下表示，例如图下的 $x^2_{B,ij}$）。
在维数正规化中胶子 Wightman 函数满足谱条件与局域对易性的假设下（\cref{asm:yzliu-wightman-pt}；复 Lorentz 协变性来自 Bargmann--Hall--Wightman 定理），微扰论各阶 $S=S_t$（\cref{thm:yzliu-time-ordering}，全阶论证 [A]；VF-YZL-12）。}
\label{fig:yzliu-complex-lorentz}
\end{figure}
\clearpage
\begin{figure}[htbp]
\centering
\begin{tikzpicture}[fd,scale=0.95,
  fd2 soft/.style={gluon,draw=green!45!black},
  fd2 hadron/.style={bound state,line cap=butt},
  fd2 dipole/.style={draw=green!40!black,line width=2.2pt,dash pattern=on 5pt off 3pt,line cap=butt},
  fd2 lab/.style={font=\small},
  fd2 note/.style={font=\footnotesize}]
  % ======================= (a) double-line notation =======================
  \begin{scope}
    % staple (space-like): from psi(0) along +z to L, transverse link, back along -z to psibar
    \draw[wilson=0.5] (1.94,1.75) -- (4.3,1.75);
    \draw[wilson=0.55] (4.3,1.75) -- (4.3,2.8);
    \draw[wilson=0.42] (4.3,2.8) -- (0.69,2.8);
    % soft subgraph S: gluons to the staple and to C
    \draw[fd2 soft] (2.98,1.07) -- (2.55,1.75);
    \draw[fd2 soft] (3.55,1.05) -- (4.3,2.25);
    \draw[fd2 soft] (2.85,0.55) -- (2.4,0.3);
    \draw[fd2 soft] (3.1,0.3) -- (2.6,0.0);
    \node[soft blob,circle,minimum size=0.9cm,inner sep=0pt] at (3.35,0.68) {$S$};
    % collinear quark lines (arrows: fermion-number flow) and longitudinal gluons
    \draw[quark=0.55] (1.68,0.48) -- (1.7,1.47);
    \draw[quark=0.5] (0.45,2.52) -- (0.55,0.5);
    \draw[longit] (2.12,0.42) -- (1.88,1.55);   \longitmark{(2.12,0.42)}{(1.88,1.55)}
    \draw[longit] (-0.02,0.38) -- (0.26,2.58);  \longitmark{(-0.02,0.38)}{(0.26,2.58)}
    \node[fd2 note,text=blue!55!black] at (2.33,1.05) {$\cdots$};
    \node[fd2 note,text=blue!55!black] at (-0.27,1.5) {$\cdots$};
    % collinear subgraph C (hadron P along +z; forward matrix element)
    \draw[fd2 hadron] (-0.95,-0.75) -- (0.25,-0.2);
    \draw[fd2 hadron] (2.55,-0.2) -- (3.75,-0.75);
    \node[jet blob,ellipse,minimum width=2.7cm,minimum height=1.05cm,inner sep=0pt] at (1.4,0.0) {$C$};
    % hard vertices at the two quark--link junctions
    \node[hard blob,circle,minimum size=0.54cm,inner sep=0pt] at (1.7,1.75) {\footnotesize$H$};
    \node[hard blob,circle,minimum size=0.54cm,inner sep=0pt] at (0.45,2.8) {\footnotesize$H$};
    % labels
    \node[fd2 note,anchor=north west] at (-0.95,-0.75) {$P$};
    \node[fd2 note,anchor=north east] at (3.75,-0.75) {$P$};
    \node[fd2 note,anchor=east] at (1.4,1.78) {$0$};
    \node[fd2 note,anchor=south] at (0.45,3.08) {$z\hat e_z+\vec b_\perp$};
    \node[fd2 note,anchor=west] at (4.38,2.27) {$L\to\infty$};
    \node[fd2 lab] at (1.9,-1.5) {(a)\ 双线记号};
  \end{scope}
  % ======================= (b) single-line notation =======================
  \begin{scope}[xshift=7.15cm]
    % dipole line H -- S and soft gluons S -- C
    \draw[fd2 dipole] (1.33,2.3) -- (3.12,2.3);
    \draw[fd2 soft] (3.35,1.86) -- (2.42,0.3);
    \draw[fd2 soft] (3.8,1.9) -- (2.66,0.02);
    \node[soft blob,circle,minimum size=0.9cm,inner sep=0pt] at (3.55,2.3) {$S$};
    % collinear quark lines and longitudinal gluons into H
    \draw[quark=0.55] (0.75,0.48) -- (0.75,1.97);
    \draw[quark=0.55] (1.15,1.97) -- (1.2,0.5);
    \draw[longit] (1.75,0.42) -- (1.34,2.07);   \longitmark{(1.75,0.42)}{(1.34,2.07)}
    \draw[longit] (0.12,0.4) -- (0.56,2.07);    \longitmark{(0.12,0.4)}{(0.56,2.07)}
    % C and H
    \draw[fd2 hadron] (-0.95,-0.75) -- (0.25,-0.2);
    \draw[fd2 hadron] (2.55,-0.2) -- (3.75,-0.75);
    \node[jet blob,ellipse,minimum width=2.7cm,minimum height=1.05cm,inner sep=0pt] at (1.4,0.0) {$C$};
    \node[hard blob,circle,minimum size=0.75cm,inner sep=0pt] at (0.95,2.3) {$H$};
    \node[fd2 note,anchor=north west] at (-0.95,-0.75) {$P$};
    \node[fd2 note,anchor=north east] at (3.75,-0.75) {$P$};
    \node[fd2 note,anchor=south,text=green!30!black] at (2.25,2.82) {偶极线：量纲 $\frac32$、自旋 $0$};
    \node[fd2 lab] at (1.4,-1.5) {(b)\ 单线记号};
    % power counting
    \node[draw=gray!60,rounded corners=2pt,fill=gray!5,inner sep=4pt,align=left,font=\footnotesize,anchor=west] at (4.35,0.95)
      {$F(H)=-\frac72+2+\frac32=0$\\[3pt]
       $2$：两条共线夸克线\\[1pt]
       $\frac32$：一条偶极线\\[3pt]
       $F(SC)\ge0$};
  \end{scope}
\end{tikzpicture}
\caption{准 TMD 的领头区域（示意；幂次计数参数为 $1/(b_\perp P^z)$）。
(a)~双线记号：共线子图 $C$（强子动量 $P$ 沿 $+\hat z$，前向矩阵元）、位于两个夸克—规范线连接点 $0$ 与 $z\hat e_z+\vec b_\perp$ 的硬顶点 $H$、类空订书钉（双线，箭头为路径方向：自 $\psi(0)$ 沿 $+\hat z$ 到 $L\to\infty$，经横向连接回到 $\bar\psi$；与\cref{fig:tmd-staple}中的类光订书钉不同）与软子图 $S$；细线箭头为费米子数流向；$C$ 与 $H$ 之间可有任意多条纵向共线胶子（带点的蓝色胶子），$S$ 与 $C$ 之间可有任意多条软胶子（绿色）。
(b)~单线记号：一对订书钉线合成一条量纲 $\frac32$、自旋 $0$ 的偶极线（粗虚线）。
幂次计数 $F(H)=0$、$F(SC)\ge0$ 表明没有超领头区域，且硬动量不能在 $0$ 与 $\vec b_\perp$ 之间传递，因而匹配中没有 $x$ 卷积（VF-YZL-8；\cref{prop:yzliu-qtmd-powercount}，全阶论证 [A]）。}
\label{fig:yzliu-leading-region}
\end{figure}
\clearpage
\begin{figure}[htbp]
\centering
\begin{tikzpicture}[fd,
  fd2 ax/.style={lcaxis,gray!60!black},
  fd2 soft/.style={gluon,draw=green!45!black},
  fd2 hadron/.style={bound state,line cap=butt},
  fd2 hv/.style={hard blob,circle,minimum size=0.42cm,inner sep=0pt},
  fd2 lab/.style={font=\small},
  fd2 note/.style={font=\footnotesize}]
  % ======================= (a) DIS at x -> 1 =======================
  \begin{scope}
    % remnant and soft radiation between jet and remnant
    \draw[draw=gray!60!black,line width=0.8pt] (1.45,-0.28) -- (4.85,-0.98);
    \draw[fd2 soft] (3.98,1.12) -- (3.82,-0.76);
    \draw[fd2 soft] (4.55,1.46) -- (4.5,-0.9);
    % struck quark, photon, outgoing jet
    \draw[quark=0.5] (1.5,0.1) -- (2.95,0.38);
    \draw[photon] (2.45,1.75) -- (2.95,0.38);
    \draw[quark=0.6] (2.95,0.38) -- (3.6,1.06);
    \node[jetB blob,ellipse,minimum width=1.75cm,minimum height=0.5cm,inner sep=0pt,rotate=40] at (4.12,1.55) {};
    \node[fd2 hv] at (2.95,0.38) {};
    % target
    \draw[fd2 hadron] (-0.75,0) -- (0.3,0);
    \node[hadron blob,ellipse,minimum width=1.3cm,minimum height=0.85cm,inner sep=0pt] at (0.9,0) {$f$};
    % labels
    \node[fd2 note,anchor=south] at (-0.25,0.05) {$P$};
    \node[fd2 note,anchor=north] at (2.15,0.12) {$xP$};
    \node[fd2 note,anchor=east] at (2.45,1.6) {$q$};
    \node[fd2 note,anchor=north west] at (3.02,0.3) {$H$};
    \node[fd2 note,anchor=south west] at (4.6,2.02) {窄喷注 $J$};
    \node[fd2 note,anchor=north,text=green!30!black] at (4.15,-1.08) {软辐射 $S$};
    \node[fd2 note,anchor=north east,text=gray!35!black] at (2.6,-0.55) {残余 $(1-x)P$};
    \node[fd2 note,align=center] at (2.05,-1.95) {$H$：$Q^2$；\ $J$：$p_J^2\sim(1-x)Q^2$；\ $S$：能量 $\sim(1-x)Q$};
    \node[fd2 lab] at (2.05,-2.6) {(a)\ $x\to1$ 的 DIS};
  \end{scope}
  % ======================= (b) quasi-PDF kernel at xi -> 1 =======================
  \begin{scope}[xshift=9.25cm]
    \draw[fd2 ax] (-2.3,1.0) -- (-1.7,1.0) node[right=-1pt,fd2 note] {$z$};
    \draw[fd2 ax] (-2.3,1.0) -- (-2.3,1.6) node[above=-1pt,fd2 note] {$t$};
    % space-like jet function: link + semi-hard gluons (dashed box)
    \draw[region] (-0.95,-1.05) rectangle (3.95,1.05);
    \draw[fd2 soft] (-0.62,-0.62) .. controls (-0.05,-0.82) and (0.85,-0.62) .. (1.15,0);
    \draw[fd2 soft] (1.85,0) .. controls (2.15,0.62) and (3.0,0.82) .. (3.62,0.62);
    \draw[wilson=0.42] (0,0) -- (3.0,0);
    % collinear quark: in along n at 0, out along n at z
    \draw[quark=0.55] (-1.35,-1.35) -- (0,0);
    \draw[quark=0.5] (3.0,0) -- (4.35,1.35);
    \node[fd2 hv] at (0,0) {};
    \node[fd2 hv] at (3.0,0) {};
    % labels
    \node[fd2 note,anchor=south east] at (-1.02,-1.0) {$p$};
    \node[fd2 note,anchor=east] at (-0.25,0.02) {$H$};
    \node[fd2 note,anchor=north west] at (3.1,-0.1) {$H$};
    \node[fd2 note,anchor=south] at (0.0,0.22) {$0$};
    \node[fd2 note,anchor=south] at (3.0,0.2) {$z$};
    \node[fd2 note,anchor=north] at (2.35,-0.14) {$W[z\hat e_z,0]$};
    \node[fd2 note,anchor=south west,text=gray!35!black] at (-0.95,1.07) {虚线框：类空喷注函数 $\tilde J$};
    \node[fd2 note,align=center] at (1.5,-1.95) {$H$：硬，$\zeta_z=4p_z^2$；\ 绿色：半硬胶子，$|1-\xi|\,p^z$};
    \node[fd2 lab] at (1.5,-2.6) {(b)\ $\xi\to1$ 的准 PDF 匹配核};
  \end{scope}
  % ======================= threshold factorization =======================
  \node[fd2 lab] at (6.85,-3.55) {$C\Big(\xi,\dfrac{p^z}{\mu}\Big)\Big|_{\xi\to1}=H\Big(\dfrac{4p_z^2}{\mu^2}\Big)\,p^z\,J_f\Big(\dfrac{(1-\xi)p^z}{\mu},\dfrac{4p_z^2}{\mu^2}\Big)\big[1+\order{1-\xi}\big]$，\quad $J_f=\tilde J\,\ee^{\ii\,\mathrm{sgn}(z)A}$};
\end{tikzpicture}
\caption{阈值因子化的时空图像（示意）。
(a)~$x\to1$ 的 DIS：被击出的夸克在硬顶点 $H$ 吸收虚光子后形成窄喷注 $J$，喷注与靶残余之间只有软辐射 $S$。
(b)~$\xi\to1$ 的准 PDF 匹配核（$\Gamma=\gamma^0$，$\MSbar$，非单态）：共线夸克（动量 $p$，沿类光方向；细线箭头为费米子数流向）在类空 Wilson 线 $W[z\hat e_z,0]$（双线，箭头为路径方向）的两个连接点处发生硬相互作用（灰色顶点 $H$，标度 $\zeta_z=4p_z^2$），并与静态规范线交换半硬胶子（绿色，标度 $|1-\xi|p^z$）；虚线框内的部分组成类空喷注函数 $\tilde J$，乘上相位因子后为 $J_f$。
$H$ 与准 TMD 因子化中的硬核相同（\cref{thm:yzliu-threshold}，全阶论证 [A]；VF-LAM-14，VF-YZL-9）。
这一因子化形式在单圈精确成立；精确的两圈核在阈值极限下与它的全部 41 个奇异项相容（其中两个常数由比较定出）。}
\label{fig:yzliu-threshold}
\end{figure}
\clearpage
\begin{figure}[htbp]
\centering
\begin{tikzpicture}[fd,x=2.05cm,y=1cm,
  fd2 ax/.style={lcaxis,gray!60!black},
  fd2 tick/.style={draw=gray!60!black,line width=0.45pt},
  fd2 pole/.pic={\draw[line width=1.1pt,line cap=round] (-0.045,-0.09)--(0.045,0.09) (-0.045,0.09)--(0.045,-0.09);},
  fd2 contour/.style={draw=blue!60!black,line width=0.9pt,postaction={decorate,decoration={markings,
      mark=at position 0.62 with {\arrow{Stealth[length=5pt,width=4.2pt]}}}}},
  fd2 pointer/.style={draw=#1,line width=0.45pt,-{Stealth[length=3.5pt,width=3pt]}},
  fd2 note/.style={font=\footnotesize},
  fd2 lab/.style={font=\small}]
  \foreach \panel/\xs in {a/0,b/3.75} {
  \begin{scope}[xshift=\xs*2.05cm]
    \draw[fd2 ax] (-0.45,0) -- (2.55,0) node[right] {$\Real u$};
    \draw[fd2 ax] (0,-0.8) -- (0,1.25) node[left] {$\Imag u$};
    \foreach \u/\t in {0.5/{\frac12},1/{1},1.5/{\frac32},2/{2}} {
      \draw[fd2 tick] (\u,-0.06) -- (\u,0.06);
      \node[font=\scriptsize,anchor=north,inner sep=1.5pt] at (\u,-0.5) {$\t$};}
    \node[font=\scriptsize,anchor=north east,inner sep=1.5pt] at (0,0) {$0$};
    \draw[fd2 contour] (0,0) .. controls (0.1,0.2) and (0.2,0.26) .. (0.32,0.26) -- (2.35,0.26);
    \if b\panel
      \draw[fd2 contour] (0,0) .. controls (0.1,-0.2) and (0.2,-0.26) .. (0.32,-0.26) -- (2.35,-0.26);
      \node[font=\footnotesize,text=blue!50!black,anchor=west,inner sep=1.5pt] at (2.35,0.3) {$C_+$};
      \node[font=\footnotesize,text=blue!50!black,anchor=west,inner sep=1.5pt] at (2.35,-0.3) {$C_-$};
    \fi
  \end{scope}}
  % ---------------- (a) real part (modulus): regular at u = 1/2 ----------------
  \draw[line width=0.9pt,fill=white,draw=gray!50!black] (0.5,0) circle[radius=0.075cm*1.6];
  \draw[fd2 pointer=gray!55!black] (0.95,0.83) -- (0.56,0.12);
  \node[fd2 note,anchor=west,align=left,text=gray!35!black] at (0.97,0.83) {$u=\frac12$ 处正则（无线性重整化子）};
  \node[fd2 lab,align=center] at (1.05,-1.95) {(a)\ 核的实部};
  % ---------------- (b) imaginary part (phase A): simple pole at u = 1/2 ----------------
  \begin{scope}[xshift=3.75*2.05cm]
    \pic[red!70!black] at (0.5,0) {fd2 pole};
    \draw[fd2 pointer=red!60!black] (0.95,0.83) -- (0.55,0.12);
    \node[fd2 note,anchor=west,align=left,text=red!50!black] at (0.97,0.83) {单极点 $\propto\dfrac{\mathrm{sgn}(z)}{u-\frac12}$};
    \node[fd2 note,anchor=north,align=center,text=blue!45!black] at (1.2,-0.95) {$C_+$ 与 $C_-$ 之差 $\propto\Lambda/p^z$（线性幂次）};
    \node[fd2 lab,align=center] at (1.05,-1.95) {(b)\ 核的虚部（相位 $A$）};
  \end{scope}
  % ---------------- formulas ----------------
  \node[fd2 lab,align=center] at (2.92,-3.25) {$\Imag B(u)\to\dfrac{\mu}{p^z}\,\dfrac{C_F\ee^{5/6}}{2\pi}\,\dfrac{\mathrm{sgn}(z)}{u-\frac12}$\\[6pt]
    $A_{\mathrm{asym}}=-2\dfrac{\mu}{p^z}\displaystyle\sum_n r_n\alpha_s^{n+1}$，\quad $r_n=N_m\Big(\dfrac{\beta_0}{2\pi}\Big)^{\!n}n!$，\quad $N_m=\dfrac{C_F\ee^{5/6}}{\pi}\simeq0.977$};
\end{tikzpicture}
\caption{大 $\beta_0$（泡链）近似下重—轻 Sudakov 核的 Borel 平面（$u=\beta_0t/(4\pi)$，约定同\cref{def:power-borel}；只画出 $u=\frac12$ 附近，其余奇点未画出）。
线性红外重整化子位于 $u=\frac12$，只出现在核的虚部即相位 $A$ 中 (b)，使沿正实轴的 Borel 积分产生正比于 $\Lambda/p^z$ 的不确定性（绕极点上方 $C_+$ 或下方 $C_-$）；核的实部在该处正则 (a)；轻—轻 Sudakov 核没有这一奇点。
它与次领头幂次阈值软算符（及含维数 1 的 Wilson 线尖点算符的软因子）的紫外重整化子相消（VF-YZL-11；\cref{prop:yzliu-linear-renormalon}，部分 [P]）。
归一化 $N_m$ 与极点质量、Wilson 线线性重整化子相同；它在大 $\beta_0$ 之外的普适性是猜想（\cref{conj:yzliu-renormalon-norm}，[C]）。}
\label{fig:yzliu-borel}
\end{figure}
\clearpage
\begin{figure}[htbp]
\centering
\begin{tikzpicture}[fd,x=1cm,y=1cm,
  fd2 dot/.style={circle,fill=#1,draw=#1!70!black,line width=0.4pt,inner sep=0pt,minimum size=0.176cm},
  fd2 hol/.style={circle,fill=white,draw=#1,line width=0.85pt,inner sep=0pt,minimum size=0.166cm},
  fd2 row/.style={font=\footnotesize,anchor=west,inner sep=0pt},
  fd2 ids/.style={font=\tiny,anchor=west,inner sep=0pt,text=gray!45!black},
  fd2 fn/.style={font=\tiny,anchor=west,inner sep=0.5pt}]
  \fill[gray!7] (0,-1.440) rectangle (15.950,-0.720);
  \fill[gray!7] (0,-2.880) rectangle (15.950,-2.160);
  \fill[gray!7] (0,-4.320) rectangle (15.950,-3.600);
  \fill[gray!7] (0,-5.760) rectangle (15.950,-5.040);
  \fill[gray!7] (0,-7.200) rectangle (15.950,-6.480);
  \fill[gray!7] (0,-8.640) rectangle (15.950,-7.920);
  \fill[gray!7] (0,-10.080) rectangle (15.950,-9.360);
  \fill[gray!7] (0,-11.520) rectangle (15.950,-10.800);
  \fill[gray!7] (0,-12.960) rectangle (15.950,-12.240);
  \fill[gray!7] (0,-14.400) rectangle (15.950,-13.680);
  \node[inner sep=0pt] at (5.600,0.420) {\tagR};
  \node[inner sep=0pt] at (7.900,0.420) {\tagA};
  \node[inner sep=0pt] at (10.200,0.420) {\tagP};
  \node[inner sep=0pt] at (12.500,0.420) {\tagC};
  \node[inner sep=0pt] at (14.800,0.420) {\tagV};
  \node[font=\footnotesize\bfseries,anchor=west,inner sep=0pt] at (0,0.420) {过程类};
  \draw[gray!60,line width=0.6pt] (0,0.000) -- (15.950,0.000);
  \draw[gray!55,line width=0.6pt] (0,-5.040) -- (15.950,-5.040);
  \draw[gray!55,line width=0.6pt] (0,-9.360) -- (15.950,-9.360);
  \draw[gray!55,line width=0.6pt] (0,-12.240) -- (15.950,-12.240);
  \draw[gray!60,line width=0.6pt] (0,-14.400) -- (15.950,-14.400);
  \draw[gray!35,line width=0.4pt] (4.450,0.000) -- (4.450,-14.400);
  \draw[gray!35,line width=0.4pt] (6.750,0.000) -- (6.750,-14.400);
  \draw[gray!35,line width=0.4pt] (9.050,0.000) -- (9.050,-14.400);
  \draw[gray!35,line width=0.4pt] (11.350,0.000) -- (11.350,-14.400);
  \draw[gray!35,line width=0.4pt] (13.650,0.000) -- (13.650,-14.400);
  \draw[gray!35,line width=0.4pt] (15.950,0.000) -- (15.950,-14.400);
  \node[fd2 row] at (0.06,-0.250) {$e^+e^-$ 总截面};
  \node[fd2 ids] at (0.06,-0.560) {II.1--7};
  \node[fd2 dot=OliveGreen] at (4.660,-0.360) {};  % II.1
  \node[fd2 dot=OliveGreen] at (4.888,-0.360) {};  % II.2
  \node[fd2 dot=OliveGreen] at (5.116,-0.360) {};  % II.4
  \node[fd2 dot=RoyalBlue] at (6.960,-0.360) {};  % II.3
  \node[fd2 dot=RoyalBlue] at (7.188,-0.360) {};  % II.5
  \node[fd2 dot=Plum] at (11.560,-0.360) {};  % II.6
  \node[fd2 dot=Plum] at (11.788,-0.360) {};  % II.7
  \node[fd2 row] at (0.06,-0.970) {DIS 矩、OPE 与极化求和规则};
  \node[fd2 ids] at (0.06,-1.280) {II.8--15,\,17--19,\,28,\,29；III.21--24};
  \node[fd2 dot=OliveGreen] at (4.660,-1.080) {};  % II.17
  \node[fd2 dot=OliveGreen] at (4.888,-1.080) {};  % II.18
  \node[fd2 dot=OliveGreen] at (5.116,-1.080) {};  % II.28
  \node[fd2 dot=OliveGreen] at (5.344,-1.080) {};  % III.21
  \node[fd2 hol=OliveGreen] at (5.632,-1.080) {};  % II.12
  \node[fd2 hol=OliveGreen] at (5.860,-1.080) {};  % II.13
  \node[fd2 hol=OliveGreen] at (6.088,-1.080) {};  % II.14
  \node[fd2 hol=OliveGreen] at (6.316,-1.080) {};  % II.15
  \node[fd2 fn] at (6.434,-0.990) {(a)};
  \node[fd2 dot=RoyalBlue] at (6.960,-1.080) {};  % II.8
  \node[fd2 dot=RoyalBlue] at (7.188,-1.080) {};  % II.9
  \node[fd2 dot=RoyalBlue] at (7.416,-1.080) {};  % II.10
  \node[fd2 dot=RoyalBlue] at (7.644,-1.080) {};  % II.11
  \node[fd2 dot=RoyalBlue] at (7.872,-1.080) {};  % II.19
  \node[fd2 hol=RoyalBlue] at (8.160,-1.080) {};  % II.29
  \node[fd2 dot=Plum] at (11.560,-1.080) {};  % III.23
  \node[fd2 dot=BrickRed] at (13.860,-1.080) {};  % III.22
  \node[fd2 dot=BrickRed] at (14.088,-1.080) {};  % III.24
  \node[fd2 row] at (0.06,-1.690) {DIS $x$ 空间};
  \node[fd2 ids] at (0.06,-2.000) {II.16,\,25；III.10--15；VII.30};
  \node[fd2 dot=OliveGreen] at (4.660,-1.800) {};  % II.25
  \node[fd2 dot=OliveGreen] at (4.888,-1.800) {};  % III.14
  \node[fd2 hol=OliveGreen] at (5.176,-1.800) {};  % II.16
  \node[fd2 fn] at (5.294,-1.710) {(a)};
  \node[fd2 dot=RoyalBlue] at (6.960,-1.800) {};  % III.13
  \node[fd2 dot=RoyalBlue] at (7.188,-1.800) {};  % III.15
  \node[fd2 dot=RoyalBlue] at (7.416,-1.800) {};  % VII.30
  \node[fd2 hol=RoyalBlue] at (7.704,-1.800) {};  % III.10
  \node[fd2 hol=RoyalBlue] at (7.932,-1.800) {};  % III.11
  \node[fd2 hol=Plum] at (11.560,-1.800) {};  % III.12
  \node[fd2 fn] at (11.678,-1.710) {(c)};
  \node[fd2 row] at (0.06,-2.410) {SIA、SIDIS 与光生};
  \node[fd2 ids] at (0.06,-2.720) {III.16--19};
  \node[fd2 dot=RoyalBlue] at (6.960,-2.520) {};  % III.16
  \node[fd2 dot=RoyalBlue] at (7.188,-2.520) {};  % III.17
  \node[fd2 dot=RoyalBlue] at (7.416,-2.520) {};  % III.18
  \node[fd2 dot=Plum] at (11.560,-2.520) {};  % III.19
  \node[fd2 row] at (0.06,-3.130) {DY 与色单态末态};
  \node[fd2 ids] at (0.06,-3.440) {III.25--30,\,33,\,34};
  \node[fd2 dot=OliveGreen] at (4.660,-3.240) {};  % III.27
  \node[fd2 dot=RoyalBlue] at (6.960,-3.240) {};  % III.26
  \node[fd2 dot=RoyalBlue] at (7.188,-3.240) {};  % III.29
  \node[fd2 dot=RoyalBlue] at (7.416,-3.240) {};  % III.30
  \node[fd2 dot=RoyalBlue] at (7.644,-3.240) {};  % III.33
  \node[fd2 hol=RoyalBlue] at (7.932,-3.240) {};  % III.25
  \node[fd2 dot=BurntOrange] at (9.260,-3.240) {};  % III.28
  \node[fd2 dot=BrickRed] at (13.860,-3.240) {};  % III.34
  \node[fd2 row] at (0.06,-3.850) {$pp\to$ 喷注/强子（有色末态）};
  \node[fd2 ids] at (0.06,-4.160) {III.31,\,32；IV.24--26};
  \node[fd2 dot=Plum] at (11.560,-3.960) {};  % III.31
  \node[fd2 dot=Plum] at (11.788,-3.960) {};  % III.32
  \node[fd2 dot=Plum] at (12.016,-3.960) {};  % IV.25
  \node[fd2 dot=BrickRed] at (13.860,-3.960) {};  % IV.24
  \node[fd2 dot=BrickRed] at (14.088,-3.960) {};  % IV.26
  \node[fd2 row] at (0.06,-4.570) {阈值重新因子化};
  \node[fd2 ids] at (0.06,-4.880) {III.35--37};
  \node[fd2 dot=RoyalBlue] at (6.960,-4.680) {};  % III.35
  \node[fd2 dot=RoyalBlue] at (7.188,-4.680) {};  % III.36
  \node[fd2 dot=BurntOrange] at (9.260,-4.680) {};  % III.37
  \node[fd2 row] at (0.06,-5.290) {TMD：DY、SIDIS、$e^+e^-$};
  \node[fd2 ids] at (0.06,-5.600) {IV.1--9；VII.48};
  \node[fd2 dot=OliveGreen] at (4.660,-5.400) {};  % IV.2
  \node[fd2 dot=OliveGreen] at (4.888,-5.400) {};  % IV.6
  \node[fd2 dot=RoyalBlue] at (6.960,-5.400) {};  % IV.3
  \node[fd2 dot=RoyalBlue] at (7.188,-5.400) {};  % IV.4
  \node[fd2 dot=RoyalBlue] at (7.416,-5.400) {};  % IV.5
  \node[fd2 dot=RoyalBlue] at (7.644,-5.400) {};  % IV.7
  \node[fd2 dot=RoyalBlue] at (7.872,-5.400) {};  % IV.8
  \node[fd2 hol=RoyalBlue] at (8.160,-5.400) {};  % IV.1
  \node[fd2 dot=BurntOrange] at (9.260,-5.400) {};  % IV.9
  \node[fd2 dot=Plum] at (11.560,-5.400) {};  % VII.48
  \node[fd2 row] at (0.06,-6.010) {TMD：$h_1h_2\to H_1H_2X$};
  \node[fd2 ids] at (0.06,-6.320) {IV.22,\,23};
  \node[fd2 hol=RoyalBlue] at (6.960,-6.120) {};  % IV.22
  \node[fd2 dot=BrickRed] at (13.860,-6.120) {};  % IV.23
  \node[fd2 row] at (0.06,-6.730) {SCET、事件形状与喷注};
  \node[fd2 ids] at (0.06,-7.040) {IV.10--17；VII.23--25};
  \node[fd2 dot=OliveGreen] at (4.660,-6.840) {};  % IV.10
  \node[fd2 dot=RoyalBlue] at (6.960,-6.840) {};  % IV.11
  \node[fd2 dot=RoyalBlue] at (7.188,-6.840) {};  % IV.12
  \node[fd2 dot=RoyalBlue] at (7.416,-6.840) {};  % IV.13
  \node[fd2 dot=RoyalBlue] at (7.644,-6.840) {};  % IV.15
  \node[fd2 dot=RoyalBlue] at (7.872,-6.840) {};  % VII.23
  \node[fd2 dot=BurntOrange] at (9.260,-6.840) {};  % IV.14
  \node[fd2 dot=BurntOrange] at (9.488,-6.840) {};  % IV.16
  \node[fd2 dot=BurntOrange] at (9.716,-6.840) {};  % VII.24
  \node[fd2 dot=Plum] at (11.560,-6.840) {};  % IV.17
  \node[fd2 dot=BrickRed] at (13.860,-6.840) {};  % VII.25
  \node[fd2 row] at (0.06,-7.450) {小 $x$};
  \node[fd2 ids] at (0.06,-7.760) {IV.18--21；VII.31};
  \node[fd2 dot=RoyalBlue] at (6.960,-7.560) {};  % IV.18
  \node[fd2 dot=RoyalBlue] at (7.188,-7.560) {};  % IV.20
  \node[fd2 dot=BurntOrange] at (9.260,-7.560) {};  % IV.19
  \node[fd2 dot=Plum] at (11.560,-7.560) {};  % IV.21
  \node[fd2 dot=BrickRed] at (13.860,-7.560) {};  % VII.31
  \node[fd2 row] at (0.06,-8.170) {衍射};
  \node[fd2 ids] at (0.06,-8.480) {III.20；IV.28,\,29};
  \node[fd2 dot=RoyalBlue] at (6.960,-8.280) {};  % IV.28
  \node[fd2 dot=BrickRed] at (13.860,-8.280) {};  % III.20
  \node[fd2 dot=BrickRed] at (14.088,-8.280) {};  % IV.29
  \node[fd2 row] at (0.06,-8.890) {双部分子散射与多重相互作用};
  \node[fd2 ids] at (0.06,-9.200) {IV.27,\,30--32；VII.22,\,46};
  \node[fd2 dot=RoyalBlue] at (6.960,-9.000) {};  % IV.30
  \node[fd2 dot=RoyalBlue] at (7.188,-9.000) {};  % VII.22
  \node[fd2 dot=Plum] at (11.560,-9.000) {};  % IV.31
  \node[fd2 dot=Plum] at (11.788,-9.000) {};  % IV.32
  \node[fd2 dot=Plum] at (12.016,-9.000) {};  % VII.46
  \node[fd2 dot=BrickRed] at (13.860,-9.000) {};  % IV.27
  \node[fd2 row] at (0.06,-9.610) {硬遍举过程};
  \node[fd2 ids] at (0.06,-9.920) {V.1--8};
  \node[fd2 dot=RoyalBlue] at (6.960,-9.720) {};  % V.2
  \node[fd2 dot=RoyalBlue] at (7.188,-9.720) {};  % V.3
  \node[fd2 dot=RoyalBlue] at (7.416,-9.720) {};  % V.5
  \node[fd2 dot=RoyalBlue] at (7.644,-9.720) {};  % V.6
  \node[fd2 hol=RoyalBlue] at (7.932,-9.720) {};  % V.1
  \node[fd2 dot=BurntOrange] at (9.260,-9.720) {};  % V.4
  \node[fd2 dot=BurntOrange] at (9.488,-9.720) {};  % V.8
  \node[fd2 dot=BrickRed] at (13.860,-9.720) {};  % V.7
  \node[fd2 row] at (0.06,-10.330) {$B$ 介子衰变};
  \node[fd2 ids] at (0.06,-10.640) {V.9--13};
  \node[fd2 dot=RoyalBlue] at (6.960,-10.440) {};  % V.9
  \node[fd2 dot=RoyalBlue] at (7.188,-10.440) {};  % V.10
  \node[fd2 dot=RoyalBlue] at (7.416,-10.440) {};  % V.13
  \node[fd2 dot=BurntOrange] at (9.260,-10.440) {};  % V.11
  \node[fd2 dot=BrickRed] at (13.860,-10.440) {};  % V.12
  \node[fd2 row] at (0.06,-11.050) {重夸克偶素（NRQCD）};
  \node[fd2 ids] at (0.06,-11.360) {V.14--16};
  \node[fd2 dot=RoyalBlue] at (6.960,-11.160) {};  % V.14
  \node[fd2 dot=RoyalBlue] at (7.188,-11.160) {};  % V.16
  \node[fd2 dot=Plum] at (11.560,-11.160) {};  % V.15
  \node[fd2 row] at (0.06,-11.770) {幂次修正、NLP 与重整化子};
  \node[fd2 ids] at (0.06,-12.080) {V.17--24,\,26；VII.26,\,47};
  \node[fd2 dot=OliveGreen] at (4.660,-11.880) {};  % V.21
  \node[fd2 dot=RoyalBlue] at (6.960,-11.880) {};  % V.17
  \node[fd2 dot=RoyalBlue] at (7.188,-11.880) {};  % VII.26
  \node[fd2 dot=BurntOrange] at (9.260,-11.880) {};  % V.18
  \node[fd2 dot=BurntOrange] at (9.488,-11.880) {};  % V.20
  \node[fd2 dot=BurntOrange] at (9.716,-11.880) {};  % V.22
  \node[fd2 dot=BurntOrange] at (9.944,-11.880) {};  % V.23
  \node[fd2 dot=BurntOrange] at (10.172,-11.880) {};  % V.26
  \node[fd2 dot=Plum] at (11.560,-11.880) {};  % V.19
  \node[fd2 dot=Plum] at (11.788,-11.880) {};  % V.24
  \node[fd2 dot=Plum] at (12.016,-11.880) {};  % VII.47
  \node[fd2 row] at (0.06,-12.490) {格点与 LaMET（共线）};
  \node[fd2 ids] at (0.06,-12.800) {VII.1--9,\,16,\,17,\,33--38,\,40--42};
  \node[fd2 dot=OliveGreen] at (4.660,-12.600) {};  % VII.1
  \node[fd2 dot=OliveGreen] at (4.888,-12.600) {};  % VII.6
  \node[fd2 dot=OliveGreen] at (5.116,-12.600) {};  % VII.7
  \node[fd2 dot=RoyalBlue] at (6.960,-12.600) {};  % VII.2
  \node[fd2 dot=RoyalBlue] at (7.188,-12.600) {};  % VII.3
  \node[fd2 dot=RoyalBlue] at (7.416,-12.600) {};  % VII.16
  \node[fd2 dot=RoyalBlue] at (7.644,-12.600) {};  % VII.34
  \node[fd2 dot=RoyalBlue] at (7.872,-12.600) {};  % VII.35
  \node[fd2 dot=RoyalBlue] at (8.100,-12.600) {};  % VII.38
  \node[fd2 hol=RoyalBlue] at (8.388,-12.600) {};  % VII.5
  \node[fd2 hol=RoyalBlue] at (8.616,-12.600) {};  % VII.33
  \node[fd2 dot=BurntOrange] at (9.260,-12.600) {};  % VII.17
  \node[fd2 dot=Plum] at (11.560,-12.600) {};  % VII.4
  \node[fd2 dot=Plum] at (11.788,-12.600) {};  % VII.36
  \node[fd2 dot=Plum] at (12.016,-12.600) {};  % VII.40
  \node[fd2 dot=Plum] at (12.244,-12.600) {};  % VII.41
  \node[fd2 dot=Plum] at (12.472,-12.600) {};  % VII.42
  \node[fd2 hol=Plum] at (12.760,-12.600) {};  % VII.8
  \node[fd2 hol=Plum] at (12.988,-12.600) {};  % VII.9
  \node[fd2 fn] at (13.106,-12.510) {(b)};
  \node[fd2 dot=BrickRed] at (13.860,-12.600) {};  % VII.37
  \node[fd2 row] at (0.06,-13.210) {准 TMD 与软函数};
  \node[fd2 ids] at (0.06,-13.520) {VII.10--15,\,39};
  \node[fd2 dot=RoyalBlue] at (6.960,-13.320) {};  % VII.10
  \node[fd2 dot=RoyalBlue] at (7.188,-13.320) {};  % VII.11
  \node[fd2 dot=RoyalBlue] at (7.416,-13.320) {};  % VII.12
  \node[fd2 dot=RoyalBlue] at (7.644,-13.320) {};  % VII.13
  \node[fd2 dot=RoyalBlue] at (7.872,-13.320) {};  % VII.14
  \node[fd2 dot=RoyalBlue] at (8.100,-13.320) {};  % VII.15
  \node[fd2 dot=Plum] at (11.560,-13.320) {};  % VII.39
  \node[fd2 row] at (0.06,-13.930) {Landau 理论、区域与 OPE};
  \node[fd2 ids] at (0.06,-14.240) {I.10--12,\,22,\,23；V.25；VII.18--21,\,43--45};
  \node[fd2 dot=OliveGreen] at (4.660,-14.040) {};  % I.10
  \node[fd2 dot=OliveGreen] at (4.888,-14.040) {};  % I.11
  \node[fd2 dot=OliveGreen] at (5.116,-14.040) {};  % I.12
  \node[fd2 dot=OliveGreen] at (5.344,-14.040) {};  % I.22
  \node[fd2 dot=OliveGreen] at (5.572,-14.040) {};  % V.25
  \node[fd2 dot=RoyalBlue] at (6.960,-14.040) {};  % VII.18
  \node[fd2 hol=RoyalBlue] at (7.248,-14.040) {};  % VII.19
  \node[fd2 dot=BurntOrange] at (9.260,-14.040) {};  % I.23
  \node[fd2 dot=Plum] at (11.560,-14.040) {};  % VII.20
  \node[fd2 dot=Plum] at (11.788,-14.040) {};  % VII.45
  \node[fd2 dot=BrickRed] at (13.860,-14.040) {};  % VII.21
  \node[fd2 dot=BrickRed] at (14.088,-14.040) {};  % VII.43
  \node[fd2 dot=BrickRed] at (14.316,-14.040) {};  % VII.44
  \node[fd2 dot=gray,anchor=center] at (0.15,-14.820) {};
  \node[font=\footnotesize,anchor=west,inner sep=0pt] at (0.32,-14.820) {表~\ref{tab:summary-master} 中的一条陈述，画在其规范标签所在列（$R_{\mathrm{PT}}$、$R_{\mathrm{def}}$、$R_{\mathrm{OPE}}$ 均在 [R] 列）；颜色同标签};
  \node[fd2 hol=darkgray,anchor=center] at (0.15,-15.240) {};
  \node[font=\footnotesize,anchor=west,inner sep=0pt] at (0.32,-15.240) {条件定理 $[\mathrm R,R_{\mathrm{cond}}]$，画在其内容等级所在列（全局假设 \texttt{asm:intro-*} 不降低内容等级）};
  \node[font=\footnotesize,anchor=north west,inner sep=0pt,align=left,text width=15.85cm] at (0.06,-15.590) {[A] 列中的空心圆：部分子层次的条件定理，前提为 [A] 级技术假设（\texttt{asm:appM-*}、\texttt{asm:appN-*}、\texttt{asm:appO-*}、\cref{asm:lamet-ope}）或 [A] 级定理。(a)~前提只含全局非微扰假设，II.14、II.15 的瓶颈为 \cref{asm:intro-np-ope,asm:intro-growth}；\ (b)~强子层次，瓶颈为 [C] 级假设：VII.8 为 \cref{asm:lamet-uniform,asm:lamet-twist4,asm:lamet-large-z}，VII.9 为 \cref{asm:lamet-gap}；\ (c)~部分子层次的正规子路线，内容等级经由 \cref{asm:appN-ir-regulator} 为 [C]。};
\end{tikzpicture}
\caption{各类过程中因子化陈述的地位一览（与表~\ref{tab:summary-master} 一致）。每个圆点代表该表的一行：实心圆画在其规范标签所在的列，空心圆为条件定理 $[\mathrm R,R_{\mathrm{cond}}]$，画在其内容等级所在的列，[C] 列的空心圆是内容等级为 [C] 的条件定理（VII.8、VII.9 为强子层次的陈述），图下注明其瓶颈假设；左栏小字为归入该类的表行编号。基础、方法与部分子分布性质的陈述（区块 I 的大部分、II.20--II.24、II.26、II.27、III.1--III.9、III.38）以及 VII.27--VII.29、VII.32 不属于某一过程类，未画出。同一过程类往往同时含有不同等级的陈述：标签属于陈述，而不属于过程。}
\label{fig:summary-status-map}
\end{figure}
\clearpage
\begin{figure}[htbp]
\centering
\begin{tikzpicture}[fd,x=1cm,y=1cm,
  fd2 asm/.style={draw=orange!65!black,dashed,fill=orange!8,rounded corners=1.5pt,line width=0.5pt,minimum height=0.34cm,inner sep=1.5pt,font=\tiny\ttfamily,anchor=center},
  fd2 thm/.style={draw=#1!65!black,fill=#1!15,rounded corners=1.5pt,line width=0.5pt,minimum height=0.34cm,inner sep=1.5pt,font=\tiny\ttfamily,anchor=center},
  fd2 rcond/.style={fd2 thm=OliveGreen,dashed},
  fd2 tag/.style={font=\tiny\sffamily\bfseries,text=#1!75!black,inner sep=0.5pt},
  fd2 head/.style={font=\scriptsize\bfseries,inner sep=1pt},
  fd2 e/.style={draw=gray!60,line width=0.35pt,-{Stealth[length=2.6pt,width=2.2pt]}},
  fd2 b/.style={draw=red!60!black,line width=1.0pt,opacity=0.75,-{Stealth[length=3.4pt,width=3pt]}}]
  \draw[gray!45,line width=0.4pt,rounded corners=3pt,fill=gray!3] (-0.120,-6.304) rectangle (3.120,-0.174);
  \node[fd2 head] at (1.500,-0.394) {减除法（附录~N）};
  \draw[gray!45,line width=0.4pt,rounded corners=3pt,fill=gray!3] (-0.120,-12.809) rectangle (3.120,-9.689);
  \node[fd2 head] at (1.500,-9.909) {全局非微扰假设};
  \draw[gray!45,line width=0.4pt,rounded corners=3pt,fill=gray!3] (-0.120,-14.719) rectangle (3.120,-12.889);
  \node[fd2 head] at (1.500,-13.109) {OPE（附录~M）};
  \draw[gray!45,line width=0.4pt,rounded corners=3pt,fill=gray!3] (12.630,-5.870) rectangle (15.870,-3.180);
  \node[fd2 head] at (14.250,-3.400) {Glauber（附录~O）};
  \draw[gray!45,line width=0.4pt,rounded corners=3pt,fill=gray!3] (12.630,-10.360) rectangle (15.870,-5.950);
  \node[fd2 head] at (14.250,-6.170) {LaMET（第~22~章）};
  \draw[gray!45,line width=0.4pt,rounded corners=3pt,fill=gray!3] (12.630,-12.270) rectangle (15.870,-10.440);
  \node[fd2 head] at (14.250,-10.660) {Liu（第~23~章）};
  \draw[gray!45,line width=0.4pt,rounded corners=3pt,fill=gray!3] (12.630,-13.750) rectangle (15.870,-12.350);
  \node[fd2 head] at (14.250,-12.570) {其他前沿（第~24~章）};
  \draw[gray!45,line width=0.4pt,rounded corners=3pt,fill=gray!3] (5.955,-4.210) rectangle (9.945,0.200);
  \node[fd2 head] at (7.950,-0.020) {共线因子化};
  \draw[gray!45,line width=0.4pt,rounded corners=3pt,fill=gray!3] (5.955,-6.730) rectangle (9.945,-4.470);
  \node[fd2 head] at (7.950,-4.690) {Glauber 区与 DY 型过程};
  \draw[gray!45,line width=0.4pt,rounded corners=3pt,fill=gray!3] (5.955,-12.260) rectangle (9.945,-6.990);
  \node[fd2 head] at (7.950,-7.210) {前沿};
  \draw[gray!45,line width=0.4pt,rounded corners=3pt,fill=gray!3] (5.955,-13.920) rectangle (9.945,-12.520);
  \node[fd2 head] at (7.950,-12.740) {OPE 与非微扰矩};
  \draw[fd2 e] (3.030,-10.389) .. controls (4.930,-10.389) and (4.155,-0.930) .. (6.055,-0.930);
  \draw[fd2 e] (3.030,-10.389) .. controls (4.930,-10.389) and (4.155,-1.360) .. (6.055,-1.360);
  \draw[fd2 e] (3.030,-10.389) .. controls (4.930,-10.389) and (4.155,-1.790) .. (6.055,-1.790);
  \draw[fd2 e] (3.030,-4.744) .. controls (4.930,-4.744) and (4.155,-1.790) .. (6.055,-1.790);
  \draw[fd2 e] (3.030,-10.389) .. controls (4.930,-10.389) and (4.155,-2.220) .. (6.055,-2.220);
  \draw[fd2 e] (3.030,-10.389) .. controls (4.930,-10.389) and (4.155,-2.650) .. (6.055,-2.650);
  \draw[fd2 e] (3.030,-10.389) .. controls (4.930,-10.389) and (4.155,-3.080) .. (6.055,-3.080);
  \draw[fd2 e] (3.030,-4.314) .. controls (4.930,-4.314) and (4.155,-3.510) .. (6.055,-3.510);
  \draw[fd2 e] (3.030,-1.734) .. controls (4.930,-1.734) and (4.155,-3.510) .. (6.055,-3.510);
  \draw[fd2 e] (3.030,-2.594) .. controls (4.930,-2.594) and (4.155,-3.510) .. (6.055,-3.510);
  \draw[fd2 e] (3.030,-4.744) .. controls (4.930,-4.744) and (4.155,-3.510) .. (6.055,-3.510);
  \draw[fd2 e] (3.030,-1.304) .. controls (4.930,-1.304) and (4.155,-3.510) .. (6.055,-3.510);
  \draw[fd2 e] (3.030,-2.164) .. controls (4.930,-2.164) and (4.155,-3.510) .. (6.055,-3.510);
  \draw[fd2 e] (3.030,-3.024) .. controls (4.930,-3.024) and (4.155,-3.510) .. (6.055,-3.510);
  \draw[fd2 e] (3.030,-5.604) .. controls (4.930,-5.604) and (4.155,-3.510) .. (6.055,-3.510);
  \draw[fd2 e] (3.030,-3.454) .. controls (4.930,-3.454) and (4.155,-3.510) .. (6.055,-3.510);
  \draw[fd2 e] (3.030,-6.034) .. controls (4.930,-6.034) and (4.155,-3.510) .. (6.055,-3.510);
  \draw[fd2 e] (3.030,-0.874) .. controls (4.930,-0.874) and (4.155,-3.510) .. (6.055,-3.510);
  \draw[fd2 e] (3.030,-2.164) .. controls (4.930,-2.164) and (4.155,-3.940) .. (6.055,-3.940);
  \draw[fd2 e] (3.030,-4.314) .. controls (4.930,-4.314) and (4.155,-3.940) .. (6.055,-3.940);
  \draw[fd2 e] (3.030,-10.389) .. controls (4.930,-10.389) and (4.155,-5.170) .. (6.055,-5.170);
  \draw[fd2 e] (12.720,-4.740) .. controls (10.820,-4.740) and (11.745,-5.600) .. (9.845,-5.600);
  \draw[fd2 e] (12.720,-5.170) .. controls (10.820,-5.170) and (11.745,-5.600) .. (9.845,-5.600);
  \draw[fd2 e] (3.030,-1.734) .. controls (4.930,-1.734) and (4.155,-5.600) .. (6.055,-5.600);
  \draw[fd2 e] (3.030,-2.594) .. controls (4.930,-2.594) and (4.155,-5.600) .. (6.055,-5.600);
  \draw[fd2 e] (3.030,-3.024) .. controls (4.930,-3.024) and (4.155,-5.600) .. (6.055,-5.600);
  \draw[fd2 e] (3.030,-10.389) .. controls (4.930,-10.389) and (4.155,-6.030) .. (6.055,-6.030);
  \draw[fd2 e] (3.030,-10.389) .. controls (4.930,-10.389) and (4.155,-6.460) .. (6.055,-6.460);
  \draw[fd2 e] (3.030,-10.819) .. controls (4.930,-10.819) and (4.155,-8.980) .. (6.055,-8.980);
  \draw[fd2 e] (12.720,-9.660) .. controls (10.820,-9.660) and (11.745,-9.410) .. (9.845,-9.410);
  \draw[fd2 e] (3.030,-10.819) .. controls (4.930,-10.819) and (4.155,-13.650) .. (6.055,-13.650);
  \draw[fd2 e] (3.030,-11.249) .. controls (4.930,-11.249) and (4.155,-13.650) .. (6.055,-13.650);
  \draw[fd2 e] (3.030,-12.109) .. controls (4.930,-12.109) and (4.155,-13.650) .. (6.055,-13.650);
  \draw[fd2 b] (3.030,-10.389) .. controls (4.930,-10.389) and (4.155,-0.500) .. (6.055,-0.500);
  \draw[fd2 b] (3.030,-1.734) .. controls (4.930,-1.734) and (4.155,-0.500) .. (6.055,-0.500);
  \draw[fd2 b] (3.030,-0.874) .. controls (4.930,-0.874) and (4.155,-1.790) .. (6.055,-1.790);
  \draw[fd2 b] (3.030,-4.744) .. controls (4.930,-4.744) and (4.155,-2.220) .. (6.055,-2.220);
  \draw[fd2 b] (3.030,-1.734) .. controls (4.930,-1.734) and (4.155,-2.220) .. (6.055,-2.220);
  \draw[fd2 b] (3.030,-2.594) .. controls (4.930,-2.594) and (4.155,-2.220) .. (6.055,-2.220);
  \draw[fd2 b] (3.030,-5.604) .. controls (4.930,-5.604) and (4.155,-2.220) .. (6.055,-2.220);
  \draw[fd2 b] (3.030,-3.024) .. controls (4.930,-3.024) and (4.155,-2.220) .. (6.055,-2.220);
  \draw[fd2 b] (3.030,-1.304) .. controls (4.930,-1.304) and (4.155,-2.220) .. (6.055,-2.220);
  \draw[fd2 b] (3.030,-2.164) .. controls (4.930,-2.164) and (4.155,-2.220) .. (6.055,-2.220);
  \draw[fd2 b] (3.030,-4.744) .. controls (4.930,-4.744) and (4.155,-2.650) .. (6.055,-2.650);
  \draw[fd2 b] (3.030,-1.734) .. controls (4.930,-1.734) and (4.155,-2.650) .. (6.055,-2.650);
  \draw[fd2 b] (3.030,-2.594) .. controls (4.930,-2.594) and (4.155,-2.650) .. (6.055,-2.650);
  \draw[fd2 b] (3.030,-5.604) .. controls (4.930,-5.604) and (4.155,-2.650) .. (6.055,-2.650);
  \draw[fd2 b] (3.030,-3.024) .. controls (4.930,-3.024) and (4.155,-2.650) .. (6.055,-2.650);
  \draw[fd2 b] (3.030,-1.304) .. controls (4.930,-1.304) and (4.155,-2.650) .. (6.055,-2.650);
  \draw[fd2 b] (3.030,-2.164) .. controls (4.930,-2.164) and (4.155,-2.650) .. (6.055,-2.650);
  \draw[fd2 b] (3.030,-4.744) .. controls (4.930,-4.744) and (4.155,-3.080) .. (6.055,-3.080);
  \draw[fd2 b] (3.030,-1.734) .. controls (4.930,-1.734) and (4.155,-3.080) .. (6.055,-3.080);
  \draw[fd2 b] (3.030,-2.594) .. controls (4.930,-2.594) and (4.155,-3.080) .. (6.055,-3.080);
  \draw[fd2 b] (3.030,-5.604) .. controls (4.930,-5.604) and (4.155,-3.080) .. (6.055,-3.080);
  \draw[fd2 b] (3.030,-3.024) .. controls (4.930,-3.024) and (4.155,-3.080) .. (6.055,-3.080);
  \draw[fd2 b] (3.030,-1.304) .. controls (4.930,-1.304) and (4.155,-3.080) .. (6.055,-3.080);
  \draw[fd2 b] (3.030,-2.164) .. controls (4.930,-2.164) and (4.155,-3.080) .. (6.055,-3.080);
  \draw[fd2 b] (3.030,-3.884) .. controls (4.930,-3.884) and (4.155,-3.510) .. (6.055,-3.510);
  \draw[fd2 b] (3.030,-5.174) .. controls (4.930,-5.174) and (4.155,-3.940) .. (6.055,-3.940);
  \draw[fd2 b] (3.030,-6.034) .. controls (4.930,-6.034) and (4.155,-5.170) .. (6.055,-5.170);
  \draw[fd2 b] (3.030,-4.744) .. controls (4.930,-4.744) and (4.155,-5.170) .. (6.055,-5.170);
  \draw[fd2 b] (12.720,-4.310) .. controls (10.820,-4.310) and (11.745,-5.170) .. (9.845,-5.170);
  \draw[fd2 b] (12.720,-3.880) .. controls (10.820,-3.880) and (11.745,-5.170) .. (9.845,-5.170);
  \draw[fd2 b] (12.720,-3.880) .. controls (10.820,-3.880) and (11.745,-5.600) .. (9.845,-5.600);
  \draw[fd2 b] (12.720,-5.600) .. controls (10.820,-5.600) and (11.745,-5.600) .. (9.845,-5.600);
  \draw[fd2 b] (12.720,-4.310) .. controls (10.820,-4.310) and (11.745,-5.600) .. (9.845,-5.600);
  \draw[fd2 b] (12.720,-3.880) .. controls (10.820,-3.880) and (11.745,-6.030) .. (9.845,-6.030);
  \draw[fd2 b] (12.720,-5.600) .. controls (10.820,-5.600) and (11.745,-6.030) .. (9.845,-6.030);
  \draw[fd2 b] (12.720,-4.310) .. controls (10.820,-4.310) and (11.745,-6.030) .. (9.845,-6.030);
  \draw[fd2 b] (12.720,-4.740) .. controls (10.820,-4.740) and (11.745,-6.030) .. (9.845,-6.030);
  \draw[fd2 b] (12.720,-5.170) .. controls (10.820,-5.170) and (11.745,-6.030) .. (9.845,-6.030);
  \draw[fd2 b] (3.030,-4.744) .. controls (4.930,-4.744) and (4.155,-6.030) .. (6.055,-6.030);
  \draw[fd2 b] (3.030,-5.604) .. controls (4.930,-5.604) and (4.155,-6.030) .. (6.055,-6.030);
  \draw[fd2 b] (12.720,-3.880) .. controls (10.820,-3.880) and (11.745,-6.460) .. (9.845,-6.460);
  \draw[fd2 b] (12.720,-5.600) .. controls (10.820,-5.600) and (11.745,-6.460) .. (9.845,-6.460);
  \draw[fd2 b] (12.720,-4.310) .. controls (10.820,-4.310) and (11.745,-6.460) .. (9.845,-6.460);
  \draw[fd2 b] (12.720,-4.740) .. controls (10.820,-4.740) and (11.745,-6.460) .. (9.845,-6.460);
  \draw[fd2 b] (12.720,-5.170) .. controls (10.820,-5.170) and (11.745,-6.460) .. (9.845,-6.460);
  \draw[fd2 b] (3.030,-4.744) .. controls (4.930,-4.744) and (4.155,-6.460) .. (6.055,-6.460);
  \draw[fd2 b] (3.030,-5.604) .. controls (4.930,-5.604) and (4.155,-6.460) .. (6.055,-6.460);
  \draw[fd2 b] (12.720,-7.080) .. controls (10.820,-7.080) and (11.745,-7.690) .. (9.845,-7.690);
  \draw[fd2 b] (12.720,-6.650) .. controls (10.820,-6.650) and (11.745,-7.690) .. (9.845,-7.690);
  \draw[fd2 b] (12.720,-7.080) .. controls (10.820,-7.080) and (11.745,-8.120) .. (9.845,-8.120);
  \draw[fd2 b] (12.720,-6.650) .. controls (10.820,-6.650) and (11.745,-8.120) .. (9.845,-8.120);
  \draw[fd2 b] (12.720,-7.510) .. controls (10.820,-7.510) and (11.745,-8.550) .. (9.845,-8.550);
  \draw[fd2 b] (12.720,-7.940) .. controls (10.820,-7.940) and (11.745,-8.550) .. (9.845,-8.550);
  \draw[fd2 b] (12.720,-8.370) .. controls (10.820,-8.370) and (11.745,-8.550) .. (9.845,-8.550);
  \draw[fd2 b] (12.720,-8.800) .. controls (10.820,-8.800) and (11.745,-8.980) .. (9.845,-8.980);
  \draw[fd2 b] (12.720,-9.230) .. controls (10.820,-9.230) and (11.745,-9.410) .. (9.845,-9.410);
  \draw[fd2 b] (12.720,-10.090) .. controls (10.820,-10.090) and (11.745,-9.410) .. (9.845,-9.410);
  \draw[fd2 b] (12.720,-11.140) .. controls (10.820,-11.140) and (11.745,-9.840) .. (9.845,-9.840);
  \draw[fd2 b] (12.720,-11.570) .. controls (10.820,-11.570) and (11.745,-10.270) .. (9.845,-10.270);
  \draw[fd2 b] (12.720,-11.570) .. controls (10.820,-11.570) and (11.745,-10.700) .. (9.845,-10.700);
  \draw[fd2 b] (12.720,-12.000) .. controls (10.820,-12.000) and (11.745,-11.130) .. (9.845,-11.130);
  \draw[fd2 b] (12.720,-13.050) .. controls (10.820,-13.050) and (11.745,-11.560) .. (9.845,-11.560);
  \draw[fd2 b] (12.720,-13.480) .. controls (10.820,-13.480) and (11.745,-11.990) .. (9.845,-11.990);
  \draw[fd2 b] (3.030,-13.589) .. controls (4.930,-13.589) and (4.155,-13.220) .. (6.055,-13.220);
  \draw[fd2 b] (3.030,-14.019) .. controls (4.930,-14.019) and (4.155,-13.220) .. (6.055,-13.220);
  \draw[fd2 b] (3.030,-14.449) .. controls (4.930,-14.449) and (4.155,-13.220) .. (6.055,-13.220);
  \draw[fd2 b] (3.030,-11.679) .. controls (4.930,-11.679) and (4.155,-13.650) .. (6.055,-13.650);
  \draw[fd2 b] (3.030,-12.539) .. controls (4.930,-12.539) and (4.155,-13.650) .. (6.055,-13.650);
  \node[fd2 asm,minimum width=3.00cm] at (1.500,-0.874) {};
  \node[font=\tiny\ttfamily,anchor=west,inner sep=0pt] at (0.070,-0.874) {appN-endpoint};
  \node[fd2 tag=RoyalBlue,anchor=east] at (2.950,-0.874) {A};
  \node[fd2 asm,minimum width=3.00cm] at (1.500,-1.304) {};
  \node[font=\tiny\ttfamily,anchor=west,inner sep=0pt] at (0.070,-1.304) {appN-uv};
  \node[fd2 tag=RoyalBlue,anchor=east] at (2.950,-1.304) {A};
  \node[fd2 asm,minimum width=3.00cm] at (1.500,-1.734) {};
  \node[font=\tiny\ttfamily,anchor=west,inner sep=0pt] at (0.070,-1.734) {appN-envelope};
  \node[fd2 tag=RoyalBlue,anchor=east] at (2.950,-1.734) {A};
  \node[fd2 asm,minimum width=3.00cm] at (1.500,-2.164) {};
  \node[font=\tiny\ttfamily,anchor=west,inner sep=0pt] at (0.070,-2.164) {appN-multren};
  \node[fd2 tag=RoyalBlue,anchor=east] at (2.950,-2.164) {A};
  \node[fd2 asm,minimum width=3.00cm] at (1.500,-2.594) {};
  \node[font=\tiny\ttfamily,anchor=west,inner sep=0pt] at (0.070,-2.594) {appN-separation};
  \node[fd2 tag=RoyalBlue,anchor=east] at (2.950,-2.594) {A};
  \node[fd2 asm,minimum width=3.00cm] at (1.500,-3.024) {};
  \node[font=\tiny\ttfamily,anchor=west,inner sep=0pt] at (0.070,-3.024) {appN-ward-product};
  \node[fd2 tag=RoyalBlue,anchor=east] at (2.950,-3.024) {A};
  \node[fd2 asm,minimum width=3.00cm] at (1.500,-3.454) {};
  \node[font=\tiny\ttfamily,anchor=west,inner sep=0pt] at (0.070,-3.454) {appN-values};
  \node[fd2 asm,minimum width=3.00cm] at (1.500,-3.884) {};
  \node[font=\tiny\ttfamily,anchor=west,inner sep=0pt] at (0.070,-3.884) {appN-brst-nonlocal};
  \node[fd2 tag=Plum,anchor=east] at (2.950,-3.884) {C};
  \node[fd2 asm,minimum width=3.00cm] at (1.500,-4.314) {};
  \node[font=\tiny\ttfamily,anchor=west,inner sep=0pt] at (0.070,-4.314) {appN-contour};
  \node[fd2 tag=RoyalBlue,anchor=east] at (2.950,-4.314) {A};
  \node[fd2 asm,minimum width=3.00cm] at (1.500,-4.744) {};
  \node[font=\tiny\ttfamily,anchor=west,inner sep=0pt] at (0.070,-4.744) {appN-collins};
  \node[fd2 tag=RoyalBlue,anchor=east] at (2.950,-4.744) {A};
  \node[fd2 asm,minimum width=3.00cm] at (1.500,-5.174) {};
  \node[font=\tiny\ttfamily,anchor=west,inner sep=0pt] at (0.070,-5.174) {appN-ir-regulator};
  \node[fd2 tag=Plum,anchor=east] at (2.950,-5.174) {C};
  \node[fd2 asm,minimum width=3.00cm] at (1.500,-5.604) {};
  \node[font=\tiny\ttfamily,anchor=west,inner sep=0pt] at (0.070,-5.604) {appN-cutsum};
  \node[fd2 tag=RoyalBlue,anchor=east] at (2.950,-5.604) {A};
  \node[fd2 asm,minimum width=3.00cm] at (1.500,-6.034) {};
  \node[font=\tiny\ttfamily,anchor=west,inner sep=0pt] at (0.070,-6.034) {appN-tmd};
  \node[fd2 tag=RoyalBlue,anchor=east] at (2.950,-6.034) {A};
  \node[fd2 asm,minimum width=3.00cm] at (1.500,-10.389) {};
  \node[font=\tiny\ttfamily,anchor=west,inner sep=0pt] at (0.070,-10.389) {intro-bridge};
  \node[fd2 tag=Plum,anchor=east] at (2.950,-10.389) {C};
  \node[fd2 asm,minimum width=3.00cm] at (1.500,-10.819) {};
  \node[font=\tiny\ttfamily,anchor=west,inner sep=0pt] at (0.070,-10.819) {intro-wightman};
  \node[fd2 tag=Plum,anchor=east] at (2.950,-10.819) {C};
  \node[fd2 asm,minimum width=3.00cm] at (1.500,-11.249) {};
  \node[font=\tiny\ttfamily,anchor=west,inner sep=0pt] at (0.070,-11.249) {intro-composite};
  \node[fd2 tag=Plum,anchor=east] at (2.950,-11.249) {C};
  \node[fd2 asm,minimum width=3.00cm] at (1.500,-11.679) {};
  \node[font=\tiny\ttfamily,anchor=west,inner sep=0pt] at (0.070,-11.679) {intro-np-ope};
  \node[fd2 tag=Plum,anchor=east] at (2.950,-11.679) {C};
  \node[fd2 asm,minimum width=3.00cm] at (1.500,-12.109) {};
  \node[font=\tiny\ttfamily,anchor=west,inner sep=0pt] at (0.070,-12.109) {intro-af-coeff};
  \node[fd2 tag=Plum,anchor=east] at (2.950,-12.109) {C};
  \node[fd2 asm,minimum width=3.00cm] at (1.500,-12.539) {};
  \node[font=\tiny\ttfamily,anchor=west,inner sep=0pt] at (0.070,-12.539) {intro-growth};
  \node[fd2 tag=Plum,anchor=east] at (2.950,-12.539) {C};
  \node[fd2 asm,minimum width=3.00cm] at (1.500,-13.589) {};
  \node[font=\tiny\ttfamily,anchor=west,inner sep=0pt] at (0.070,-13.589) {appM-remainder};
  \node[fd2 tag=RoyalBlue,anchor=east] at (2.950,-13.589) {A};
  \node[fd2 asm,minimum width=3.00cm] at (1.500,-14.019) {};
  \node[font=\tiny\ttfamily,anchor=west,inner sep=0pt] at (0.070,-14.019) {appM-brst-ope};
  \node[fd2 tag=RoyalBlue,anchor=east] at (2.950,-14.019) {A};
  \node[fd2 asm,minimum width=3.00cm] at (1.500,-14.449) {};
  \node[font=\tiny\ttfamily,anchor=west,inner sep=0pt] at (0.070,-14.449) {appM-forward};
  \node[fd2 tag=RoyalBlue,anchor=east] at (2.950,-14.449) {A};
  \node[fd2 asm,minimum width=3.00cm] at (14.250,-3.880) {};
  \node[font=\tiny\ttfamily,anchor=west,inner sep=0pt] at (12.820,-3.880) {appO-mixing};
  \node[fd2 tag=RoyalBlue,anchor=east] at (15.700,-3.880) {A};
  \node[fd2 asm,minimum width=3.00cm] at (14.250,-4.310) {};
  \node[font=\tiny\ttfamily,anchor=west,inner sep=0pt] at (12.820,-4.310) {appO-rewritten-cs};
  \node[fd2 tag=RoyalBlue,anchor=east] at (15.700,-4.310) {A};
  \node[fd2 asm,minimum width=3.00cm] at (14.250,-4.740) {};
  \node[font=\tiny\ttfamily,anchor=west,inner sep=0pt] at (12.820,-4.740) {appO-lcopt};
  \node[fd2 tag=RoyalBlue,anchor=east] at (15.700,-4.740) {A};
  \node[fd2 asm,minimum width=3.00cm] at (14.250,-5.170) {};
  \node[font=\tiny\ttfamily,anchor=west,inner sep=0pt] at (12.820,-5.170) {appO-envelope};
  \node[fd2 tag=RoyalBlue,anchor=east] at (15.700,-5.170) {A};
  \node[fd2 asm,minimum width=3.00cm] at (14.250,-5.600) {};
  \node[font=\tiny\ttfamily,anchor=west,inner sep=0pt] at (12.820,-5.600) {appO-deformation};
  \node[fd2 tag=RoyalBlue,anchor=east] at (15.700,-5.600) {A};
  \node[fd2 asm,minimum width=3.00cm] at (14.250,-6.650) {};
  \node[font=\tiny\ttfamily,anchor=west,inner sep=0pt] at (12.820,-6.650) {lamet-aux-renorm};
  \node[fd2 tag=RoyalBlue,anchor=east] at (15.700,-6.650) {A};
  \node[fd2 asm,minimum width=3.00cm] at (14.250,-7.080) {};
  \node[font=\tiny\ttfamily,anchor=west,inner sep=0pt] at (12.820,-7.080) {lamet-ope};
  \node[fd2 tag=RoyalBlue,anchor=east] at (15.700,-7.080) {A};
  \node[fd2 asm,minimum width=3.00cm] at (14.250,-7.510) {};
  \node[font=\tiny\ttfamily,anchor=west,inner sep=0pt] at (12.820,-7.510) {lamet-uniform};
  \node[fd2 tag=Plum,anchor=east] at (15.700,-7.510) {C};
  \node[fd2 asm,minimum width=3.00cm] at (14.250,-7.940) {};
  \node[font=\tiny\ttfamily,anchor=west,inner sep=0pt] at (12.820,-7.940) {lamet-twist4};
  \node[fd2 tag=Plum,anchor=east] at (15.700,-7.940) {C};
  \node[fd2 asm,minimum width=3.00cm] at (14.250,-8.370) {};
  \node[font=\tiny\ttfamily,anchor=west,inner sep=0pt] at (12.820,-8.370) {lamet-large-z};
  \node[fd2 tag=Plum,anchor=east] at (15.700,-8.370) {C};
  \node[fd2 asm,minimum width=3.00cm] at (14.250,-8.800) {};
  \node[font=\tiny\ttfamily,anchor=west,inner sep=0pt] at (12.820,-8.800) {lamet-gap};
  \node[fd2 tag=Plum,anchor=east] at (15.700,-8.800) {C};
  \node[fd2 asm,minimum width=3.00cm] at (14.250,-9.230) {};
  \node[font=\tiny\ttfamily,anchor=west,inner sep=0pt] at (12.820,-9.230) {lamet-staple};
  \node[fd2 tag=RoyalBlue,anchor=east] at (15.700,-9.230) {A};
  \node[fd2 asm,minimum width=3.00cm] at (14.250,-9.660) {};
  \node[font=\tiny\ttfamily,anchor=west,inner sep=0pt] at (12.820,-9.660) {lamet-soft-rapidity};
  \node[fd2 tag=RoyalBlue,anchor=east] at (15.700,-9.660) {A};
  \node[fd2 asm,minimum width=3.00cm] at (14.250,-10.090) {};
  \node[font=\tiny\ttfamily,anchor=west,inner sep=0pt] at (12.820,-10.090) {lamet-order-limits};
  \node[fd2 tag=RoyalBlue,anchor=east] at (15.700,-10.090) {A};
  \node[fd2 asm,minimum width=3.00cm] at (14.250,-11.140) {};
  \node[font=\tiny\ttfamily,anchor=west,inner sep=0pt] at (12.820,-11.140) {yzliu-heavy-pair};
  \node[fd2 tag=Plum,anchor=east] at (15.700,-11.140) {C};
  \node[fd2 asm,minimum width=3.00cm] at (14.250,-11.570) {};
  \node[font=\tiny\ttfamily,anchor=west,inner sep=0pt] at (12.820,-11.570) {yzliu-wightman-pt};
  \node[fd2 tag=RoyalBlue,anchor=east] at (15.700,-11.570) {A};
  \node[fd2 asm,minimum width=3.00cm] at (14.250,-12.000) {};
  \node[font=\tiny\ttfamily,anchor=west,inner sep=0pt] at (12.820,-12.000) {yzliu-lightcone-uniform};
  \node[fd2 tag=BurntOrange,anchor=east] at (15.700,-12.000) {P};
  \node[fd2 asm,minimum width=3.00cm] at (14.250,-13.050) {};
  \node[font=\tiny\ttfamily,anchor=west,inner sep=0pt] at (12.820,-13.050) {front-mor-complete};
  \node[fd2 tag=Plum,anchor=east] at (15.700,-13.050) {C};
  \node[fd2 asm,minimum width=3.00cm] at (14.250,-13.480) {};
  \node[font=\tiny\ttfamily,anchor=west,inner sep=0pt] at (12.820,-13.480) {front-scet-modes};
  \node[fd2 tag=RoyalBlue,anchor=east] at (15.700,-13.480) {A};
  \node[fd2 thm=RoyalBlue,minimum width=3.75cm] at (7.950,-0.500) {};
  \node[font=\tiny\ttfamily,anchor=west,inner sep=0pt] at (6.145,-0.500) {scet-thrust};
  \node[fd2 thm=RoyalBlue,minimum width=3.75cm] at (7.950,-0.930) {};
  \node[font=\tiny\ttfamily,anchor=west,inner sep=0pt] at (6.145,-0.930) {heavy-bdpi};
  \node[fd2 thm=RoyalBlue,minimum width=3.75cm] at (7.950,-1.360) {};
  \node[font=\tiny\ttfamily,anchor=west,inner sep=0pt] at (6.145,-1.360) {excl-dvcs};
  \node[fd2 thm=RoyalBlue,minimum width=3.75cm] at (7.950,-1.790) {};
  \node[font=\tiny\ttfamily,anchor=west,inner sep=0pt] at (6.145,-1.790) {excl-pion-ff};
  \node[fd2 thm=RoyalBlue,minimum width=3.75cm] at (7.950,-2.220) {};
  \node[font=\tiny\ttfamily,anchor=west,inner sep=0pt] at (6.145,-2.220) {disall-factorization};
  \node[fd2 thm=RoyalBlue,minimum width=3.75cm] at (7.950,-2.650) {};
  \node[font=\tiny\ttfamily,anchor=west,inner sep=0pt] at (6.145,-2.650) {coll-sia};
  \node[fd2 thm=RoyalBlue,minimum width=3.75cm] at (7.950,-3.080) {};
  \node[font=\tiny\ttfamily,anchor=west,inner sep=0pt] at (6.145,-3.080) {break-ddis};
  \node[fd2 rcond,minimum width=3.75cm] at (7.950,-3.510) {};
  \node[fd2 tag=Plum,anchor=east] at (9.775,-3.510) {C};
  \node[font=\tiny\ttfamily,anchor=west,inner sep=0pt] at (6.145,-3.510) {appN-structure};
  \node[fd2 rcond,minimum width=3.75cm] at (7.950,-3.940) {};
  \node[fd2 tag=Plum,anchor=east] at (9.775,-3.940) {C};
  \node[font=\tiny\ttfamily,anchor=west,inner sep=0pt] at (6.145,-3.940) {disall-renorm-compat};
  \node[fd2 thm=RoyalBlue,minimum width=3.75cm] at (7.950,-5.170) {};
  \node[font=\tiny\ttfamily,anchor=west,inner sep=0pt] at (6.145,-5.170) {tmd-factorization};
  \node[fd2 rcond,minimum width=3.75cm] at (7.950,-5.600) {};
  \node[fd2 tag=RoyalBlue,anchor=east] at (9.775,-5.600) {A};
  \node[font=\tiny\ttfamily,anchor=west,inner sep=0pt] at (6.145,-5.600) {appO-main};
  \node[fd2 thm=RoyalBlue,minimum width=3.75cm] at (7.950,-6.030) {};
  \node[font=\tiny\ttfamily,anchor=west,inner sep=0pt] at (6.145,-6.030) {dy-factorization};
  \node[fd2 thm=RoyalBlue,minimum width=3.75cm] at (7.950,-6.460) {};
  \node[font=\tiny\ttfamily,anchor=west,inner sep=0pt] at (6.145,-6.460) {break-double-dy};
  \node[fd2 rcond,minimum width=3.75cm] at (7.950,-7.690) {};
  \node[fd2 tag=RoyalBlue,anchor=east] at (9.775,-7.690) {A};
  \node[font=\tiny\ttfamily,anchor=west,inner sep=0pt] at (6.145,-7.690) {lamet-coordinate};
  \node[fd2 thm=RoyalBlue,minimum width=3.75cm] at (7.950,-8.120) {};
  \node[font=\tiny\ttfamily,anchor=west,inner sep=0pt] at (6.145,-8.120) {lattice-lamet};
  \node[fd2 rcond,minimum width=3.75cm] at (7.950,-8.550) {};
  \node[fd2 tag=Plum,anchor=east] at (9.775,-8.550) {C};
  \node[font=\tiny\ttfamily,anchor=west,inner sep=0pt] at (6.145,-8.550) {lamet-hadronic};
  \node[fd2 rcond,minimum width=3.75cm] at (7.950,-8.980) {};
  \node[fd2 tag=Plum,anchor=east] at (9.775,-8.980) {C};
  \node[font=\tiny\ttfamily,anchor=west,inner sep=0pt] at (6.145,-8.980) {lamet-exp-decay};
  \node[fd2 thm=RoyalBlue,minimum width=3.75cm] at (7.950,-9.410) {};
  \node[font=\tiny\ttfamily,anchor=west,inner sep=0pt] at (6.145,-9.410) {lamet-quasi-tmd};
  \node[fd2 rcond,minimum width=3.75cm] at (7.950,-9.840) {};
  \node[fd2 tag=Plum,anchor=east] at (9.775,-9.840) {C};
  \node[font=\tiny\ttfamily,anchor=west,inner sep=0pt] at (6.145,-9.840) {{\color{gray!50!black}cor:}yzliu-soft-form-factor};
  \node[fd2 thm=RoyalBlue,minimum width=3.75cm] at (7.950,-10.270) {};
  \node[font=\tiny\ttfamily,anchor=west,inner sep=0pt] at (6.145,-10.270) {yzliu-analyticity};
  \node[fd2 thm=RoyalBlue,minimum width=3.75cm] at (7.950,-10.700) {};
  \node[font=\tiny\ttfamily,anchor=west,inner sep=0pt] at (6.145,-10.700) {yzliu-time-ordering};
  \node[fd2 thm=BurntOrange,minimum width=3.75cm] at (7.950,-11.130) {};
  \node[font=\tiny\ttfamily,anchor=west,inner sep=0pt] at (6.145,-11.130) {{\color{gray!50!black}prop:}yzliu-universality-classes};
  \node[fd2 rcond,minimum width=3.75cm] at (7.950,-11.560) {};
  \node[fd2 tag=Plum,anchor=east] at (9.775,-11.560) {C};
  \node[font=\tiny\ttfamily,anchor=west,inner sep=0pt] at (6.145,-11.560) {{\color{gray!50!black}prop:}front-commuting};
  \node[fd2 thm=RoyalBlue,minimum width=3.75cm] at (7.950,-11.990) {};
  \node[font=\tiny\ttfamily,anchor=west,inner sep=0pt] at (6.145,-11.990) {front-hgg};
  \node[fd2 thm=RoyalBlue,minimum width=3.75cm] at (7.950,-13.220) {};
  \node[font=\tiny\ttfamily,anchor=west,inner sep=0pt] at (6.145,-13.220) {dis-ope};
  \node[fd2 rcond,minimum width=3.75cm] at (7.950,-13.650) {};
  \node[fd2 tag=OliveGreen,anchor=east] at (9.775,-13.650) {R};
  \node[font=\tiny\ttfamily,anchor=west,inner sep=0pt] at (6.145,-13.650) {dis-moments-np};
  \node[fd2 asm,minimum width=0.9cm,font=\tiny] at (0.5,-15.199) {asm};
  \node[font=\scriptsize,anchor=west] at (1.0,-15.199) {编号假设（右端字母为其标签）};
  \node[fd2 thm=RoyalBlue,minimum width=0.75cm,font=\tiny] at (5.45,-15.199) {[A]};
  \node[fd2 thm=BurntOrange,minimum width=0.75cm,font=\tiny] at (6.3,-15.199) {[P]};
  \node[font=\scriptsize,anchor=west] at (6.75,-15.199) {定理（填色为其标签）};
  \node[fd2 rcond,minimum width=1.1cm,font=\tiny] at (10.25,-15.199) {$R_{\mathrm{cond}}$};
  \node[font=\scriptsize,anchor=west] at (10.85,-15.199) {条件定理（右端字母为内容等级）};
  \draw[fd2 e] (0.05,-15.649) -- (0.95,-15.649);
  \node[font=\scriptsize,anchor=west] at (1.0,-15.649) {直接依赖（\texttt{\string\uses} 或前提句）};
  \draw[fd2 b] (5.1,-15.649) -- (6.0,-15.649);
  \node[font=\scriptsize,anchor=west] at (6.05,-15.649) {瓶颈假设（可经由所用陈述进入）};
\end{tikzpicture}
\caption{主要全阶论证与条件定理所依赖的编号假设（节点为去掉类型前缀的规范标签；箭头由假设指向使用它的定理）。细灰线为直接依赖（假设出现在该定理的 \texttt{\string\uses} 列表或前提句中，即表~\ref{tab:summary-assumptions} 中的 $\checkmark$）；粗红线为该定理的瓶颈假设（总表~\ref{tab:summary-master} 的瓶颈栏；条件定理取前提中等级最低者），它们也可能经由所用的陈述进入；不是瓶颈的间接依赖从略，完整的依赖闭包由 \texttt{\string\uses} 数据生成，见附录~P（\cref{fig:appP-dependency}）。全局假设 \texttt{asm:intro-*} 不降低普通陈述的等级；\texttt{appN-brst-nonlocal} 是具名缺口，只作为 \texttt{thm:appN-structure} 的前提。}
\label{fig:summary-hypotheses}
\end{figure}
\clearpage
\chapter{外部引用的占位目标}
\begin{assumption}[占位]\label{asm:appN-ir-regulator}占位。\end{assumption}
\begin{assumption}[占位]\label{asm:intro-growth}占位。\end{assumption}
\begin{assumption}[占位]\label{asm:intro-np-ope}占位。\end{assumption}
\begin{assumption}[占位]\label{asm:lamet-gap}占位。\end{assumption}
\begin{assumption}[占位]\label{asm:lamet-large-z}占位。\end{assumption}
\begin{assumption}[占位]\label{asm:lamet-ope}占位。\end{assumption}
\begin{assumption}[占位]\label{asm:lamet-twist4}占位。\end{assumption}
\begin{assumption}[占位]\label{asm:lamet-uniform}占位。\end{assumption}
\begin{assumption}[占位]\label{asm:yzliu-heavy-pair}占位。\end{assumption}
\begin{assumption}[占位]\label{asm:yzliu-lightcone-uniform}占位。\end{assumption}
\begin{assumption}[占位]\label{asm:yzliu-wightman-pt}占位。\end{assumption}
\begin{conjecture}[占位]\label{conj:lamet-decay-qcd}占位。\end{conjecture}
\begin{conjecture}[占位]\label{conj:power-renormalon-ht}占位。\end{conjecture}
\begin{conjecture}[占位]\label{conj:yzliu-renormalon-norm}占位。\end{conjecture}
\begin{corollary}[占位]\label{cor:lamet-endpoint}占位。\end{corollary}
\begin{definition}[占位]\label{def:power-borel}占位。\end{definition}
\begin{figure}[h]\centering\caption{占位}\label{fig:appP-dependency}\end{figure}
\begin{figure}[h]\centering\caption{占位}\label{fig:tmd-staple}\end{figure}
\begin{lemma}[占位]\label{lem:lamet-strip}占位。\end{lemma}
\begin{lemma}[占位]\label{lem:yzliu-partial-exp}占位。\end{lemma}
\begin{proposition}[占位]\label{prop:heavy-qcdf-endpoint}占位。\end{proposition}
\begin{proposition}[占位]\label{prop:heavy-qcdf-lp}占位。\end{proposition}
\begin{proposition}[占位]\label{prop:lamet-hybrid-kernel}占位。\end{proposition}
\begin{proposition}[占位]\label{prop:lamet-truncation}占位。\end{proposition}
\begin{proposition}[占位]\label{prop:lattice-renormalizability}占位。\end{proposition}
\begin{proposition}[占位]\label{prop:power-adler-large-beta0}占位。\end{proposition}
\begin{proposition}[占位]\label{prop:yzliu-linear-renormalon}占位。\end{proposition}
\begin{proposition}[占位]\label{prop:yzliu-pinch-ultrasoft}占位。\end{proposition}
\begin{proposition}[占位]\label{prop:yzliu-qtmd-powercount}占位。\end{proposition}
\begin{remark}[占位]\label{rem:lamet-domain}占位。\end{remark}
\begin{remark}[占位]\label{rem:lamet-ratio-large-z}占位。\end{remark}
\begin{table}[h]\centering\caption{占位}\label{tab:lamet-dependency}\end{table}
\begin{table}[h]\centering\caption{占位}\label{tab:summary-assumptions}\end{table}
\begin{table}[h]\centering\caption{占位}\label{tab:summary-master}\end{table}
\begin{theorem}[占位]\label{thm:dis-ope}占位。\end{theorem}
\begin{theorem}[占位]\label{thm:heavy-bdpi}占位。\end{theorem}
\begin{theorem}[占位]\label{thm:lamet-coordinate}占位。\end{theorem}
\begin{theorem}[占位]\label{thm:lamet-hadronic}占位。\end{theorem}
\begin{theorem}[占位]\label{thm:lamet-parton-momentum}占位。\end{theorem}
\begin{theorem}[占位]\label{thm:lattice-lamet}占位。\end{theorem}
\begin{theorem}[占位]\label{thm:pdf-dglap}占位。\end{theorem}
\begin{theorem}[占位]\label{thm:pdf-moments-local}占位。\end{theorem}
\begin{theorem}[占位]\label{thm:yzliu-analyticity}占位。\end{theorem}
\begin{theorem}[占位]\label{thm:yzliu-soft-hqet}占位。\end{theorem}
\begin{theorem}[占位]\label{thm:yzliu-threshold}占位。\end{theorem}
\begin{theorem}[占位]\label{thm:yzliu-time-ordering}占位。\end{theorem}
\begin{thebibliography}{99}
\bibitem{IzubuchiEtAl2018} 占位文献 \texttt{IzubuchiEtAl2018}。
\end{thebibliography}
\end{document}
